\section{Properties of theta-lift rows}\label{sec:pf:rowprops}

\providecommand{\rowpropsBig}{\mathsf{Big}}
\providecommand{\rowpropsB}{\mathcal{B}_+}
\providecommand{\rowpropsL}{\mathcal{L}}
\providecommand{\rowpropsth}{\vartheta}
\providecommand{\rowpropsfd}{\mathcal{F}}
\providecommand{\rowpropsS}{\mathcal{S}}
\providecommand{\rowpropsSupp}{\mathrm{supp}}
\providecommand{\rowpropsctot}{c_{\mathrm{Tot}}}
\providecommand{\rowpropsXi}{\Xi}
\providecommand{\rowpropsemb}{\iota}
\providecommand{\rowpropsLL}{\mathit{LL}}
\providecommand{\rowpropsG}{\mathbf{G}}
\providecommand{\rowpropsgiven}{\gamma}
\providecommand{\rowpropsopp}{\gamma'}
\providecommand{\rowpropsdef}{\mathrm{def}}
\providecommand{\rowpropstot}{\mathrm{tot}}
\providecommand{\rowpropsper}{\mathrm{per}}
\providecommand{\rowpropsNIter}{\mathrm{It}}
\providecommand{\rowpropscw}{\mathrm{cw}}
\providecommand{\rowpropsclass}{\mathrm{cs}}
\providecommand{\rowpropsKbig}{K^{\mathrm{big}}}
\providecommand{\rowpropsKsmall}{K^{\mathrm{sm}}}
\providecommand{\rowpropsdeg}{\deg}
\providecommand{\rowpropsAf}{\A_{\mathrm f}}

This section proves the theorem on the properties of rows of twisted theta lifts on which the rationality of rows
(\cref{prop:families:rowrat}) rests, and deduces \cref{prop:families:rowrat} from it.

\paragraph{What is proved.}
The row test $\sfname{RowTest}$ of \cref{sec:families:theta} asks for linear forms on a space of automorphic forms
which separate points and take values in the coefficient field on the theta vectors of rational points. The forms
used are the \emph{theta pairings} $\psi_K(a,\cdot)$: an automorphic form $u$ on the unitary group of the first
column space $Y_0$ of a datum is lifted along a row of $m$ theta correspondences to a form on a unitary group
$\Uni(V_{m,0})$, and paired, by a Bessel period, with one of finitely many explicit vector-valued forms on a
definite subgroup $\Uni(V_{\mathrm t_0})\subset\Uni(V_{m,0})$. The main result of the section
(\cref{pf:rowprops:rowthm}) establishes, uniformly for all data at large scale, five properties of these pairings:
\begin{enumerate}[label=(\roman*)]
\item the class sums $v_K$ which will turn out to compute the pairings are given by meaningful finite sums
($\sfname{SumOKK}$; this is proved elsewhere and only quoted);
\item the constants $c_{ijt}$ attached to the $m^2$ squares of a grid of seesaw pairs are well defined and non-zero
($\sfname{ConstK}$);
\item all the integrals defining the pairings converge ($\sfname{LinK}$);
\item the \emph{lift formula}: if the field condition holds, then there are a constant $\kappa$ and roots of unity
$\zeta_t$ in the coefficient field such that for all good rational points $x$
\begin{equation}\label{eq:pf:rowprops:formula}
\zeta_t\, \rowpropsctot(t)\, \psi_K\bigl(a,\Theta_t(j_X x)\bigr)=\kappa\, v_K(t,x,a),\qquad
\rowpropsctot(t)=\prod_{i,j}c_{ijt}
\end{equation}
($\sfname{Formula2K}$);
\item if the field condition holds, all class sums $v_K(t,x,a)$ lie in the coefficient field ($\sfname{SumRatK}$).
\end{enumerate}
Together with a non-vanishing statement for Hecke $L$-values, a Witt-tower norm formula, the injectivity coming from
$\sfname{FunctionalsSeparate}$ and the existence of an extension field (all proved in other sections of Part~II), these
properties give the row test at all pinned data (\cref{pf:rowprops:main}).

\paragraph{The idea.}
The lift formula is obtained in four steps; the method, unfolding theta integrals into orbit sums, goes back to Rallis
\cite{Rallis1984}, and its ingredients are seesaw identities \cite{Kudla1984,HarrisKudlaSweet1996} and the fact that
automorphic forms on totally definite unitary groups are functions on finite class sets
\cite{Gross1999}. The theta lifts are organised in a grid of hermitian spaces $V_{i,j}$, $0\le i,j\le m$, with seesaw
squares
\[
\begin{array}{ccc}
V_{i,j+1}&\longrightarrow&V_{i+1,j+1}\\
\big\downarrow&&\big\downarrow\\
V_{i,j}&\longrightarrow&V_{i+1,j}
\end{array}
\]
in which horizontal arrows are theta lifts twisted by characters $\chi_{m-i}$ and vertical arrows theta lifts twisted
by $\mu_t(j+1)$. The function lifted to $\Uni(V_{m,0})$ in $\psi_K$ is the iterated lift along the \emph{given path}
(down the first column, then along the first row).
\begin{enumerate}[label=\arabic*.]
\item \emph{Projection.} The theta vector of a good rational point is fixed by the Hecke projector, so the iterated lift
is the function paired (\cref{sec:pf:rowprops:proj}).
\item \emph{Flip.} For each square, the theta integrals along the two routes around the square agree up to the square
constant $c_{ijt}$ of the choice (by Fubini's theorem, the defining property of the square constants, and the
convergence and growth statements of Theorem~K and of the square statements). Flipping the $m^2$ squares one at a
time moves the given path to the \emph{opposite path} (along the last row, then down the last column), at the cost of
the factor $(z_t\rowpropsctot(t))^{-1}$, with $z_t$ a root of unity (\cref{sec:pf:rowprops:flip}).
\item \emph{Unfolding.} The Bessel pairing of a column form with a lift along the opposite path is an iterated integral
over the groups of the last column. Its integrand is bi-invariant and supported on finitely many cells, indexed by the
support of a local unfolding of the flipped tensor; on each cell the integrand is computed, and the pairing becomes a
finite sum over the cells of rational local coefficients times definite integrals over $\Uni(V_{m,m})$
(\cref{sec:pf:rowprops:unfold}).
\item \emph{Jets.} Each definite integral pairs a lift up the column of definite spaces $V_{\mathrm t_j}$ with a lift along
the last row of definite spaces $V_{i,m}$. On definite groups the lifts act on finite-dimensional algebraic data
(``jets'') by explicit maps; the definite integral equals a universal constant $\kappa$ times a class sum
(\cref{sec:pf:rowprops:jets}).
\end{enumerate}
The convergence of all these integrals and the passage of cuspidality and smoothness along the arrows of the grid
rest on a theorem on smooth cuspidal lifts (\cref{sec:pf:rowprops:cusp}). The rationality of the class sums is a
bookkeeping of fields (\cref{sec:pf:rowprops:sumrat}), and the non-vanishing of the square constants follows from the
non-vanishing of one theta integral attached to each square (\cref{sec:pf:rowprops:const}).

\paragraph{Organisation.}
\Cref{sec:pf:rowprops:setting} describes the objects: grids, Weil data, lifts along paths, the Bessel pairing, choices
and class sums. \Cref{sec:pf:rowprops:results} states the main theorem and proves \cref{prop:families:rowrat} from it.
\Cref{sec:pf:rowprops:assembly} reduces the main theorem to the statements proved in
\cref{sec:pf:rowprops:const,sec:pf:rowprops:sumrat,sec:pf:rowprops:formula,sec:pf:rowprops:proj,sec:pf:rowprops:flip,sec:pf:rowprops:unfold,sec:pf:rowprops:jets,sec:pf:rowprops:cusp}.
\Cref{sec:pf:rowprops:remarks} collects remarks on the formal statements.

\paragraph{Conventions.}
All statements are uniform in the data and are proved \emph{at large scale} in the sense of \cref{def:pf:rowprops:big}
below. The formal objects are large records; we name the components that are used and describe them as precisely as
needed, with the formal names of the definitions at the end of each definition. A ``finite combination'' of elements of
a set $X$ is a finitely supported function $X\to\C$, and a linear map defined on $X$ is extended to finite combinations
by linearity. For a hermitian space $V$ over the CM field $E$ of a datum we write $G_V=\Uni(V)$, $G_V(\A)$ for its adelic
points, $G_V(F)$ for its rational points and $G_V(\A)=G_{V,\infty}\times G_{V,\mathrm f}$ for the decomposition into
archimedean and finite parts. Integrals over quotients $G_V(F)\backslash G_V(\A)$ are integrals over a measurable
fundamental domain $\rowpropsfd_V$ for the left action of $G_V(F)$, with respect to a Haar measure $\nu_V$.

% ---------------------------------------------------------------------------
\subsection{The setting}\label{sec:pf:rowprops:setting}

\subsubsection{Data, realisations and the column of a datum}
Let $D$ be a global datum of the development (\cref{sec:families:rows}). We use the following components: a CM field
$E$ with totally real subfield $F$, a coefficient field $\rowpropsLL$ (a number field) with a fixed embedding
$\rowpropsemb:\rowpropsLL\to\C$, ranks $r,t,m$ with $m\ge2$ even and $r,t$ odd, hermitian spaces $Y_0,\dots,Y_{m-1}$ and
$X'$ over $E$ of ranks $r+jt$ ($0\le j<m$) and $r+mt$, and integers attached to the host (a splitting bound, level
primes, an auxiliary prime).

\begin{definition}[Realisations; good points]\label{def:pf:rowprops:real}
A \emph{realisation} $R$ of $D$ (a \emph{real datum} of the development) is given over a family input $A$ consisting of
a finite group $\Gamma$, an index set $T$ of \emph{twists} with a base twist $\mu_0\in T$, finite-dimensional spaces
$\mathrm{FR}(t)$ over $\rowpropsLL$ of \emph{rational points} ($t\in T$) with an action of $\Gamma$, and a finite set $X$
of characters of $\Gamma$. Its components used here are:
\begin{itemize}
\item a standard additive character $\psi$ of $\A_F/F$, archimedean data $\theta$ (Cartan involutions at the real
places), Hecke characters $\mu_t(j)$ ($t\in T$, $1\le j\le m$) and $\eta$, and for every $t$ a twisting Hecke
character $\xi_t$;
\item Weil representations $\omega_Y(t,k)$ of $\Uni(Y_k)\times\Uni(Y_{k+1})$ ($k+1<m$) and $\omega_X(t)$ of
$\Uni(Y_{m-1})\times\Uni(X')$, with splitting characters built from $\mu_t(k+1)$ and $\eta$;
\item discrete series models $M_Y(j)$, $M_X$, multiplicity spaces $H_X(t)$ of automorphic forms on $\Uni(X')$ and $H_E$
of automorphic forms on $\Uni(Y_0)$, theta maps $\Theta_t:H_X(t)\to H_E$, and $\rowpropsemb$-semilinear maps
$j_X(t):\mathrm{FR}(t)\to H_X(t)$;
\item the Hecke projector $e_0$, a linear endomorphism of the space of functions on $G_{Y_0}(\A)$.
\end{itemize}
A rational point $x\in\mathrm{FR}(t)$ is \emph{good} if it lies in the image of the idempotent
$e_X=\sum_{\chi\in X}e_\chi$ of the group algebra of $\Gamma$ acting on $\mathrm{FR}(t)$.
\formal{Formalism.L2.DIV.BlockDat.goodF}
\end{definition}

\begin{definition}[Chain vectors; row forms]\label{def:pf:rowprops:chain}
A \emph{chain vector} $d$ consists of polynomial data $P_Y(k)$, $P_X$ and a smooth vector $v$ of $M_X$. For
$y\in H_X(t)$ let $F_{t,y}=\mathrm{form}_X(t,y)$ be the automorphic form on $\Uni(X')$ attached to $y$ and $v$, and let
$\mathrm{lift}_t(y)$ be the function on $G_{Y_0}(\A)$ obtained from $F_{t,y}$ by the theta integrals down the column
$X'\to Y_{m-1}\to\dots\to Y_0$ with the Weil representations above and the Fock vectors of the polynomial data. The
\emph{row form} of $d$ is the linear map $\mathrm{form}:H_E\to\{\text{functions on }G_{Y_0}(\A)\}$ obtained by evaluating
the forms of the multiplicity space on the vector $w=w(d)$ obtained from $v$ by the iterated intertwiners of the
column.
\formal{GlobalTheory.RowTest.ChainVec}
\end{definition}

\begin{lemma}[Row forms of theta vectors]\label{pf:rowprops:formTh}
If $R$ is a chain (a condition of the standing hypotheses below), then for all $t$ and $y\in H_X(t)$
\begin{equation}\label{eq:pf:rowprops:formTh}
\mathrm{form}\bigl(\Theta_t(y)\bigr)=e_0\bigl(\mathrm{lift}_t(y)\bigr).
\end{equation}
\formallabel{GlobalTheory.RowTest.form_Th}
\end{lemma}

\begin{proof}
The chain condition identifies $\Theta_t$, evaluated on smooth vectors, with $e_0$ composed with the iterated column
lift, and identifies the iterated lift of the vector $w(d)$ with $\mathrm{lift}_t$ (each step of the column is the theta
integral with the Fock vector of the corresponding polynomial datum). Evaluating at $w(d)$ gives the claim.
\end{proof}

\subsubsection{Grids of hermitian spaces}
\begin{definition}[Grids]\label{def:pf:rowprops:grid}
The \emph{nodes} are the symbols $(i,j)$, $\mathrm c_i$ and $\mathrm t_j$ ($i,j\in\N$); a node is \emph{in range} if its
indices are at most $m$. A \emph{grid datum} $\rowpropsG$ over the level characters of $R$ and constants
$\mathrm{sc}=(c,C_F,b)$ consists of a finite set of places $S_0$ of $F$ containing a base place $v_0$ (with chosen
places $w(v)$ of $E$ above them and \emph{depths}, a large depth at $v_0$ and depth $2$ elsewhere, defining a modulus
$\mathfrak m$), characters $\chi_1,\dots,\chi_m$ of the ideles of $E$, and for each node $X$ a hermitian space $V_X$
with a lattice $\Lambda_X$, such that:
\begin{enumerate}[label=(\alph*)]
\item the first column is the column of the datum: $V_{0,j}=Y_j$ for $j<m$ and $V_{0,m}=X'$, with the lattices of $R$;
\item for $1\le i\le m$ the space $V_{i,j}$ has rank $\mathrm{rk}(i,j)$ and $\lfloor (m-j)/2\rfloor$ negative eigenvalues
at every real place; $V_{\mathrm c_i}$ has rank $\mathrm{rk}(\mathrm c_i)$ and $\lfloor(m-i+1)/2\rfloor$ negative
eigenvalues; $V_{\mathrm t_j}$ is totally definite; the lattices are self-dual in the appropriate sense;
\item $V_{\mathrm t_0}=V_{\mathrm c_m}$; $V_{\mathrm c_i}$ embeds in $V_{i,0}$ with orthogonal complement of rank $i-1$;
$V_{\mathrm t_j}$ embeds in $V_{m,j}$ with orthogonal complement of rank $m-1-j$ ($j<m$); and $V_{m,m}$ embeds in
$V_{\mathrm t_m}$ with complement of rank one; the lattices are compatible with these decompositions.
\end{enumerate}
The ranks increase by $t$ along columns and by an odd integer $s$ along rows: $\mathrm{rk}(i,j+1)=\mathrm{rk}(i,j)+t$ and
$\mathrm{rk}(i+1,j)=\mathrm{rk}(i,j)+s$. The signature condition of (b) also holds in the first column (for the column
of the datum, under the level conditions), so the spaces $V_{i,j}$ with $j\ge m-1$ and all $V_{\mathrm t_j}$ are totally
definite, and $G_{V_{\mathrm t_0}}$ is the stabiliser in $G_{V_{m,0}}$ of a subspace of rank $m-1$.
\formal{GlobalTheory.ChiGrid.GridData}
\end{definition}

\begin{definition}[Pairs, arrows and paths]\label{def:pf:rowprops:arrows}
The \emph{pairs} of a grid are $\chi(i,j)=\{(i,j),(i+1,j)\}$, $\mu(t,i,j)=\{(i,j),(i,j+1)\}$ ($t\in T$),
$\mathrm c(i)=\{\mathrm c_i,\mathrm c_{i+1}\}$ and $\mathrm t(t,j)=\{\mathrm t_j,\mathrm t_{j+1}\}$; the pair $\chi(i,j)$
carries the character $\chi_{m-i}$, the pairs $\mu(t,i,j)$ and $\mathrm t(t,j)$ the character $\mu_t(j+1)$. The
\emph{arrows} are $\rightarrow_{i,j}:(i,j)\to(i+1,j)$, $\downarrow^t_{i,j}:(i,j+1)\to(i,j)$,
$\uparrow^t_j:\mathrm t_j\to\mathrm t_{j+1}$ and two arrows between $\mathrm c_i$ and $\mathrm c_{i+1}$; each arrow has a
pair, a source and a target. A \emph{path} is a composable sequence of arrows. We use:
\begin{itemize}
\item the row path $\rho_j:(0,j)\to(1,j)\to\dots\to(m,j)$ and the down path $\delta_i^t:(i,m)\to(i,m-1)\to\dots\to(i,0)$;
\item the \emph{given path} $\rowpropsgiven_t=\delta^t_0$ followed by $\rho_0$, and the \emph{opposite path}
$\rowpropsopp_t=\rho_m$ followed by $\delta^t_m$, both from $(0,m)$ to $(m,0)$;
\item the column path $\kappa^t:\mathrm t_0\to\mathrm t_1\to\dots\to\mathrm t_m$;
\item for a square $(i,j)$ the \emph{small route} $(i,j+1)\to(i,j)\to(i+1,j)$ and the \emph{big route}
$(i,j+1)\to(i+1,j+1)\to(i+1,j)$.
\end{itemize}
\formal{GlobalTheory.ChiGrid.Pair}\formal{GlobalTheory.ChiGrid.Arrow}\formal{GlobalTheory.FunObj.oppPath}
\end{definition}

\begin{definition}[Weil data; lifts along paths]\label{def:pf:rowprops:weil}
A \emph{Weil datum} $W$ for $\rowpropsG$ consists of a standard additive character $\psi$, archimedean data $\theta$,
Cartan involutions and orthonormal frames at the real places for every node, and for every pair $p$ a representation
$\omega_p$ of $G_{V_{S(p)}}(\A)\times G_{V_{T(p)}}(\A)$ on a space $\rowpropsS_p$ of functions, which is a Weil
representation with the splitting characters of $p$ whenever $p$ is in range; here
$(V_{S(p)},V_{T(p)})$ is $(-V_{i+1,j},V_{i,j})$ for $p=\chi(i,j)$, $(V_{i,j},V_{i,j+1})$ for $p=\mu(t,i,j)$,
$(-V_{\mathrm c_{i+1}},V_{\mathrm c_i})$ for $p=\mathrm c(i)$ and $(V_{\mathrm t_j},V_{\mathrm t_{j+1}})$ for
$p=\mathrm t(t,j)$, where $-V$ denotes $V$ with the negated form. It also contains vacuum vectors and, for every coset
$\kappa$ of $\Lambda_{(i+1,0)}$ modulo $\mathfrak m$ (relative to $\Lambda_{(i,0)}$), a \emph{coset vector} in
$\rowpropsS_{\chi(i,0)}$, obtained from the vacuum vector and the coset through a fixed idempotent of the finite adeles
supported at the places of $S_0$. For $f\in\rowpropsS_p$ the
\emph{theta kernel} is $K_p(f;g,h)=\theta(\omega_p(g,h)f)$, the theta distribution applied to the translate; for an
arrow $a$ we write $K_a(f;g,h)$ with $g$ on the source and $h$ on the target. For a polynomial $Q$ the \emph{Fock vector}
$\mathrm{fock}_p(f,Q)\in\rowpropsS_p$ is the pure tensor obtained from $f$ and $Q$ with the frames at the real places.

A \emph{volume datum} $M$ consists of Haar measures $\nu_X$ on $G_{V_X}(\A)$ and fundamental domains $\rowpropsfd_X$ for
$G_{V_X}(F)$, and a measure and fundamental domain for the unipotent group $N$ of the Bessel datum below. For
$t\in T$, $M$ is \emph{of level $t$} if at every node in range the measure is normalised at the level subgroup of the
node attached to $t$ (principal congruence subgroups of $\Lambda_X$ modulo $\mathfrak m$ or maximal compact
subgroups, and in the first column the levels of $R$); $M$ is \emph{of row type} if this holds at the nodes
$(i,0)$ and $\mathrm t_0$.

For an arrow $a$, $f\in\rowpropsS_{p(a)}$ and a function $F$ on the source group, the \emph{lift}
\[
\rowpropsL_a(f)F(h)=\int_{\rowpropsfd_{\mathrm{src}(a)}}K_a(f;g,h)\,F(g)\,d\nu_{\mathrm{src}(a)}(g)
\]
is a function on the target group, and $F$ is \emph{convergent} for $(a,f)$ if the integrand is integrable for every
$h$. For a path $\gamma$ and a family $d=(d_a)$ of elements $d_a\in\rowpropsS_{p(a)}$ indexed by the arrows of $\gamma$
(a \emph{pure tensor along $\gamma$}), $\rowpropsL_\gamma(d)$ is the composite of the lifts, and for a finite
combination $\tau$ of pure tensors, $\rowpropsL_\gamma(\tau)=\sum_d\tau(d)\rowpropsL_\gamma(d)$; convergence along
$\gamma$ means convergence at each step. Composition of paths corresponds to composition of lifts.
\formal{GlobalTheory.FunObj.liftT}
\end{definition}

Volume data of every level exist: by \lref{GlobalTheory.FunObj.exists_vol}, for every grid datum satisfying the
level conditions and every $t$ there is a volume datum of level $t$ (Haar measures and fundamental domains exist at the
open compact levels of all nodes in range).

\begin{definition}[Row lifts and the Bessel pairing]\label{def:pf:rowprops:bessel}
A \emph{coset datum} $\Phi$ is a family $(\Phi_i)_{i<m}$ of cosets $\Phi_i$ as above, and a \emph{real vector} $\varrho$
is a family of polynomials $\varrho_i$. For a function $u$ on $G_{Y_0}(\A)$ the \emph{row lift} is
\[
\rowpropsth(\varrho,\Phi;u)=\rowpropsL_{\rho_0}\bigl(\mathrm{fock}_{\chi(i,0)}(\mathrm{cv}_i(\Phi_i),\varrho_i)\bigr)_{i<m}\,u,
\]
a function on $G_{V_{m,0}}(\A)$, where $\mathrm{cv}_i(\Phi_i)$ is the coset vector. Let $\iota_B:G_{V_{\mathrm t_0}}\to
G_{V_{m,0}}$ be the stabiliser embedding of \cref{def:pf:rowprops:grid}(c), $N$ the unipotent radical of the Bessel
datum attached to the decomposition $V_{m,0}=V_{\mathrm t_0}\oplus V_{\mathrm t_0}^\perp$ and $\chi_B$ its Bessel
character. For functions $\phi$ on $G_{V_{\mathrm t_0}}(\A)$ and $F$ on $G_{V_{m,0}}(\A)$ the \emph{Bessel pairing} is
\begin{equation}\label{eq:pf:rowprops:bessel}
\rowpropsB(\phi,F)=\int_{\rowpropsfd_{\mathrm t_0}}\int_{\rowpropsfd_N}\phi(h)\,F\bigl(n\,\iota_B(h)\bigr)\,
\overline{\chi_B(n)}\,dn\,dh .
\end{equation}
\formal{GlobalTheory.FunObj.thetaChi}\formal{GlobalTheory.FunObj.Bplus}
\end{definition}

\begin{definition}[Columns]\label{def:pf:rowprops:columns}
The \emph{columns} are vector-valued functions $\Psi$ on $G_{V_{\mathrm t_0},\mathrm f}$ with values in an
algebraic representation $W_{\mathrm{coef}}$ (built over $\rowpropsLL$ from the archimedean weights of the node
$\mathrm t_0$ and made unitary by a positive definite form), which are left equivariant under $G_{V_{\mathrm
t_0}}(F)$ and right invariant under the principal congruence subgroup $K_+$. From $\Psi$ and a linear form $\xi$ on
$W_{\mathrm{coef}}$ one obtains the scalar function $\phi(h)=\xi\bigl(\varrho_{\mathrm{coef}}(h)^{-1}\Psi(h_{\mathrm f})\bigr)$ on
$G_{V_{\mathrm t_0}}(\A)$, where $\varrho_{\mathrm{coef}}$ is the action of $G_{V_{\mathrm t_0}}(\A)$ on $W_{\mathrm{coef}}$.
\formal{GlobalTheory.IndefiniteRowSideFullHecke.LevelPlusSubgroup.Cols}\formal{GlobalTheory.FunObj.pairVec}
\end{definition}

\subsubsection{Square constants and real matches}
For a square $(i,j)$ ($i,j<m$) and finite combinations $\Theta$, $\Theta'$ of pure tensors along the big and the small
route, the \emph{route kernels} are the functions
\begin{align*}
\rowpropsKbig_\Theta(g,g_1)(u)&=\sum_d\Theta(d)\,K_{\rightarrow_{i,j+1}}(d_1;g,u)\,K_{\downarrow^t_{i+1,j}}(d_2;u,g_1),\\
\rowpropsKsmall_{\Theta'}(g,g_1)(u)&=\sum_d\Theta'(d)\,K_{\downarrow^t_{i,j}}(d_1;g,u)\,K_{\rightarrow_{i,j}}(d_2;u,g_1)
\end{align*}
of $u\in G_{V_{i+1,j+1}}(\A)$, resp.\ $u\in G_{V_{i,j}}(\A)$, for $g\in G_{V_{i,j+1}}(\A)$, $g_1\in G_{V_{i+1,j}}(\A)$.
For a finite combination $z$ of elements of $G_{V_{i,j}}(\A)$ put $\rowpropsdeg z=\sum_k z(k)$ and
$(z\cdot\varphi)(u)=\sum_kz(k)\varphi(uk)$.

\begin{definition}[Real matches; square constants]\label{def:pf:rowprops:square}
A \emph{real match} $\mathrm{rm}$ for $W$ consists of vacuum vectors on all pairs $\chi(i,j)$ and $\mu(t,i,j)$
(those on $\chi(i,0)$ being the vacuum vectors of $W$) and, for every square, a subspace $\mathrm{rel}_{ij}$ of pairs
(polynomial tensor along the big route, polynomial tensor along the small route) which projects onto both factors and
is a proper subspace. A pair $(\Theta,\Theta')$ is \emph{matched} at $(i,j,t)$ if $\Theta$, $\Theta'$ are obtained from
polynomial tensors related by $\mathrm{rel}_{ij}$ and from finite local data whose local components at the places
$\mathrm{loc}(t)$ (the places of $S_0$ and the places of the twist $t$) are given by a natural datum of depth at least the
local depth and its \emph{flip} (a local integral transform). A number $c$ is a
\emph{square constant} at $(i,j,t)$ if for every volume datum $M$ of level $t$, every matched pair, every $z$, $g$, $g_1$
such that $\rowpropsKbig_\Theta(g,g_1)$ and $z\cdot\rowpropsKsmall_{\Theta'}(g,g_1)$ are integrable on the fundamental
domains,
\begin{equation}\label{eq:pf:rowprops:sqconst}
\rowpropsdeg z\int_{\rowpropsfd_{i+1,j+1}}\rowpropsKbig_\Theta(g,g_1)\,d\nu
=c\int_{\rowpropsfd_{i,j}}z\cdot\rowpropsKsmall_{\Theta'}(g,g_1)\,d\nu .
\end{equation}
The square is \emph{non-zero} at $(i,j,t)$ if for some $M$ of level $t$, some matched pair and some $z,g,g_1$ with these
integrability properties, the right-hand integral is non-zero. If the square is non-zero, a square constant is unique
(cancel the non-zero integral in \eqref{eq:pf:rowprops:sqconst}).
\formal{GlobalTheory.FunObj.RealMatch}\formal{GlobalTheory.FunObj.flipAt}\formal{GlobalTheory.FunObj.IsSquareConst}\formal{GlobalTheory.FunObj.SquareNonzero}
\end{definition}

\subsubsection{Choices and class sums}
\begin{definition}[Choices]\label{def:pf:rowprops:choice}
A \emph{choice} $\mathrm{ch}$ at $R$ consists of constants $\mathrm{sc}$, a grid datum $\rowpropsG$ over the level
characters of $R$, a Weil datum $W$, a volume datum $M$, a real vector $\varrho$, a linear form $\xi$ on $W_{\mathrm
coef}$, a chain vector $d$, a real match $\mathrm{rm}$, a sum choice $S$ (\cref{def:pf:rowprops:classsum}), class
representatives, a finite enumeration $a\mapsto(\Phi_a,\Psi_a)$, $a\in\{1,\dots,n\}$, of pairs (coset datum, column
function) built from test lattices and a basis of $W_{\mathrm{coef}}$, and numbers $c_{ijt}\in\C$ ($i,j<m$, $t\in T$).
We write $\phi_a$ for the scalar column function of $(\Psi_a,\xi)$, $\mathrm{form}$ for the row form of the chain
vector, and
\begin{align}
\psi_K(a,u)&=\rowpropsB\bigl(\phi_a,\rowpropsth(\varrho,\Phi_a;\mathrm{form}(u))\bigr)\qquad(u\in H_E),\label{eq:pf:rowprops:psiK}\\
\rowpropsctot(t)&=\textstyle\prod_{i<m}\prod_{j<m}c_{ijt}.\notag
\end{align}
The pairing \emph{converges} at $(a,u)$ if $\mathrm{form}(u)$ is convergent along the row path and the
double integral \eqref{eq:pf:rowprops:bessel} converges absolutely in the iterated sense. The choice is
\emph{admissible} ($\mathrm{ch}.\mathsf{Is}$) if: $R$ satisfies a guard condition (no prime of a twist divides a fixed
ideal of the host); $W$ lies over $R$ (same $\psi$ and $\theta$, and the Weil representations, Cartan involutions and
frames of the first column are those of $R$); $M$ is of row type; $\mathrm{rm}$ lies over $R$ (its vacuum vectors in
the first column are those of $R$); $w(d)\ne0$; each $c_{ijt}$ is a square constant at $(i,j,t)$; and the real
pairing condition holds (if all $\psi_K(a,u)$ vanish, so do the pairings with every other real vector and linear form).
It is \emph{at scale} $(C_F,b)$ if $\mathrm{sc}=(1,C_F,b)$.
\formal{GlobalTheory.RowTest.Choice}\formal{GlobalTheory.RowTest.Choice.Is}\formal{GlobalTheory.RowTest.psiK}
\end{definition}

The column function $\Psi_a$ of a choice is a column in the sense of \cref{def:pf:rowprops:columns}: by
\lref{GlobalTheory.RowTest.col_mem}, it is left equivariant and right $K_+$-invariant. The pairing $\psi_K(a,\cdot)$
is homogeneous, and additive on convergent arguments.

\begin{definition}[Local unfolding terms]\label{def:pf:rowprops:local}
At each place $w\in\mathrm{loc}(t)$ the grid has a local model $\mathrm{Grid}_w$ (the local completions of the spaces,
with the local depth $e_w$ and a uniformiser), and the coset datum $\Phi$ determines a local tensor
$T^{\mathrm{flip}}_w(\Phi)$ on the arguments of $\mathrm{Grid}_w$: the image of the local tensor of the given path under
the $m^2$ local square flips. The local unfolding of $T^{\mathrm{flip}}_w(\Phi)$ (an explicit expansion over double cosets
of the local group of the last row by local level subgroups, computed from staircase representatives) has a
\emph{support}, a set of cosets, and for each coset $\beta_w$ a \emph{term} with a rational coefficient
$\mathrm{coef}_w(\beta_w)\in\Q$, a local \emph{datum} (a function of the local arguments of the column and the last row)
and local levels.
\formal{GlobalTheory.FunObj.Tflip}\formal{GlobalTheory.FunObj.GridChoice.term}
\end{definition}

\begin{definition}[Sum choices and class sums]\label{def:pf:rowprops:classsum}
A \emph{grid choice} consists of frames at the places of $\mathrm{loc}(t)$, lattice frames $\beta'_j$ of the top
spaces, standard and normal data (the maps $k_{n}$ described below), and a finite set $P_{AA}$ of places dividing the
ideals $\mathfrak a,\bar{\mathfrak a}$ of the host. A \emph{sum choice} $S$ adds a linear map from the forms of $R$ on
$\Uni(X')$ to the algebraic representation attached to the node $(0,m)$, compatible with the rational structure of
$R$; we write $\mathbf{f}(t,y)$ for the resulting vector-valued function on the finite points of $G_{X'}$. For
$t\in T$ and a coset datum $\Phi$:
\begin{itemize}
\item $\rowpropsSupp(t,\Phi)$ is the set of families $\beta=(\beta_w)_{w\in\mathrm{loc}(t)}$ with $\beta_w$ in the
support of the local unfolding of $T^{\mathrm{flip}}_w(\Phi)$;
\item for $\beta\in\rowpropsSupp(t,\Phi)$, an \emph{expansion} of $\beta$ is a finite combination $\tau$ of pure finite
data $d=(d_1,d_2)$ along the column path $\kappa^t$ and the last row $\rho_m$ whose product function equals
$\prod_w\mathrm{datum}_w(\beta_w)$ on the half-integral arguments;
\item $\mathcal C(\beta)$ is the set of families $c$ of double cosets $G_V(F)\backslash G_{V,\mathrm f}/K$, one for
each top node $\mathrm t_j$ ($j<m$) and each node $(i,m)$ of the last row, with the levels $K$ determined by the terms of
$\beta$; $\mathrm{pt}(c)$ denotes chosen representatives, and $\gamma(c)=\prod|\Gamma_c|^{-1}$ the product of the
inverses of the orders of the finite groups $G_V(F)\cap xKx^{-1}$ at the representatives;
\item for an expansion $\tau$ and $c\in\mathcal C(\beta)$, the \emph{pairing at $c$} is the bilinear form
\[
\langle v,v'\rangle_{\tau,c}=\sum_d\tau(d)\,\mathrm{br}\bigl(\mathrm{J}^{\mathrm{col}}_d(\mathrm{pt}(c))v,\ \mathrm{J}^{\mathrm{row}}_d(\mathrm{pt}(c))v'\bigr),
\]
where $\mathrm{J}^{\mathrm{col}}_d$ and $\mathrm{J}^{\mathrm{row}}_d$ are the composites of the jet maps $k_n$ of the normal datum along
the column and along the last row (linear maps between algebraic representations of the definite groups, depending on
the finite data $d$ and on the representatives), and $\mathrm{br}$ is the branching pairing between the representations
of $G_{V_{\mathrm t_m}}$ and $G_{V_{m,m}}$;
\item the \emph{class sum} is
\begin{equation}\label{eq:pf:rowprops:classsum}
\rowpropsclass_\Phi(\beta,\tau;\Psi',F')=\sum_{c\in\mathcal C(\beta)}\gamma(c)\,
\bigl\langle\Psi'(\mathrm{pt}(c)_0),F'(\mathrm{pt}(c)_0)\bigr\rangle_{\tau,c}
\end{equation}
for vector-valued functions $\Psi'$ on the finite points of $G_{V_{\mathrm t_0}}$ and $F'$ on those of $G_{X'}$;
\item $s_{AA}=\prod_{w\in P_{AA}}s_w$ is a product of local scalars, and
\begin{equation}\label{eq:pf:rowprops:vcore}
v_{\mathrm{core}}(t,\Phi;\Psi',F')=\sum_{\beta\in\rowpropsSupp(t,\Phi)}\Bigl(\prod_w\mathrm{coef}_w(\beta_w)\Bigr)
\rowpropsclass_\Phi(\beta,\tau_\beta;\Psi',F'),
\end{equation}
with $\tau_\beta$ some expansion of $\beta$ (fixed by a choice function).
\end{itemize}
For a choice $\mathrm{ch}$ the \emph{class sums} are
\[
v_K(t,x,a)=s_{AA}\cdot v_{\mathrm{core}}\bigl(t,\Phi_a;\psi_C(\Psi_a),\mathbf{f}(t,j_X(t)x)\bigr)
\qquad(t\in T,\ x\in\mathrm{FR}(t),\ 1\le a\le n),
\]
where $\psi_C(\Psi_a)$ is $\Psi_a$ with coefficients mapped to $\C$ by $\rowpropsemb$. The sums are \emph{meaningful} at
$(t,\Phi)$ ($\mathsf{SumOK}$) if the local tensors are flat-integrable, $\rowpropsSupp(t,\Phi)$ is finite, every
$\beta$ in the support has an expansion and expansions consist of data with vectors, the pairing at $c$ does not depend
on the expansion, all $\mathcal C(\beta)$ are finite and all $\gamma(c)\ne0$, and weight, interlacing and readability
conditions for the representations hold. Under $\mathsf{SumOK}$, the sum \eqref{eq:pf:rowprops:vcore} does not depend
on the choice of the expansions $\tau_\beta$, since the pairings at $c$ do not.
\formal{GlobalTheory.FunObj.brPair}\formal{GlobalTheory.FunObj.GridChoice.classSum}\formal{GlobalTheory.FunObj.GridChoice.SumOK}\formal{GlobalTheory.FunObj.vSum}
\end{definition}

\begin{definition}[Field condition]\label{def:pf:rowprops:field}
A choice satisfies the \emph{field condition} ($\mathrm{ch}.\mathsf{FieldOK}$) if the following numbers lie in
$\rowpropsemb(\rowpropsLL)$: the values of the characters of the grid and of $R$ (in particular all local values of
the twisting characters $\xi_t$), the eighth roots of unity, the roots of unity of order $\Nm(w)^{e_w}$ for
$w\in\mathrm{loc}(t)$, the phases and jets of the local unfolding terms, the square roots $\Nm(w)^{1/2}$ at the places
used, certain ratios of Hecke $L$-values and the bottom constants of the squares. For $t\in T$ let
$\rowpropsXi_t\subset\C^\times$ be the subgroup generated by the local values of the twisting character $\xi_t$.
\formal{GlobalTheory.RowTest.NumbersIn}\formal{GlobalTheory.RowTest.Choice.FieldOK}\formal{GlobalTheory.RowClaim.xiVals}
\end{definition}

\subsubsection{Statements at large scale}
\begin{definition}[Hida setting]\label{def:pf:rowprops:hida}
The standing hypotheses on $(D,R,F)$, for size data of $D$ and a face $F$, are collected in the \emph{Hida setting} of
the development: the size conditions of the size data, a tie of $E$-series, hypotheses on the host, $F$ is the maximal
totally real subfield of $E$, a core condition on $(R,F)$ (in particular the level characters of $R$ satisfy the
conditions of a level character at $\mu_0$, among them that each $\xi_t$ has finite order; the Weil representations of
$R$ form a chain; and the discrete series models are as required), conditions on the characters $\mu_t$, a group
record, a guard condition, finiteness of $T$, fullness of the Weil representations of the column, standardness of $\psi$
and of the frame of $X'$, and the triviality of one splitting character.
\formal{GlobalTheory.RowClaim.HidaSetting}
\end{definition}

\begin{definition}[Large scale]\label{def:pf:rowprops:big}
For a property $Q$ of choices, $\rowpropsBig(Q)$ means: there is $C_{F0}$ such that for every $\eps$ and every
$C_F\ge C_{F0}$ there is $x_0$ such that for every global datum $D$, all size data with exponent $\eps$ and scale
$x\ge x_0$, every family input $A$ (with finite group $\Gamma$, twists $T$ and rational points), every number field $K_s$
over which $\rowpropsLL$ is an algebra, every realisation $R$ and face $F$ in Hida setting, and every admissible choice
$\mathrm{ch}$ at $R$ at scale $(C_F,b)$, where $b$ is the scale quantity of the size data, $Q(\mathrm{ch})$ holds.
\formal{GlobalTheory.RowClaim.Big}
\end{definition}

If $\rowpropsBig(Q)$ and $\rowpropsBig(Q')$ then $\rowpropsBig(Q\wedge Q')$: this is \lref{GlobalTheory.RowClaim.Big.and}
(take the maxima of the constants $C_{F0}$ and of the thresholds $x_0$). We use it without comment to combine finitely
many statements at large scale.

\begin{definition}[The five row properties]\label{def:pf:rowprops:five}
\begin{enumerate}[label=(\alph*)]
\item $\sfname{SumOKK}=\rowpropsBig\bigl(\forall t,a:\ \mathsf{SumOK}(t,\Phi_a)\bigr)$;
\item $\sfname{ConstK}=\rowpropsBig\bigl(\forall i,j,t:$ the square $(i,j)$ is non-zero at $(i,j,t)$ and
$c_{ijt}\ne0\bigr)$;
\item $\sfname{LinK}=\rowpropsBig\bigl(\forall a,u:\ \psi_K$ converges at $(a,u)\bigr)$;
\item $\sfname{Formula2K}=\rowpropsBig\bigl(\mathrm{ch}.\mathsf{FieldOK}\Rightarrow\exists\kappa\in\C,\
\zeta:T\to\rowpropsLL$ with $\mathsf{FormulaAt}(\mathrm{ch},\kappa,\zeta)\bigr)$, where $\mathsf{FormulaAt}(\mathrm{ch},\kappa,\zeta)$
means: every $\zeta_t$ is a root of unity, and for all $t$, all good $x\in\mathrm{FR}(t)$ and all $a$,
\[
\rowpropsemb(\zeta_t)\,\rowpropsctot(t)\,\psi_K\bigl(a,\Theta_t(j_X(t)x)\bigr)=\kappa\,v_K(t,x,a);
\]
\item $\sfname{SumRatK}=\rowpropsBig\bigl(\mathrm{ch}.\mathsf{FieldOK}\Rightarrow\forall t,x,a:\
v_K(t,x,a)\in\rowpropsemb(\rowpropsLL)\bigr)$.
\end{enumerate}
\formal{GlobalTheory.RowClaim.SumOKK}\formal{GlobalTheory.RowClaim.ConstK}\formal{GlobalTheory.RowClaim.LinK}\formal{GlobalTheory.RowClaim.FormulaAt}\formal{GlobalTheory.RowClaim.Formula2K}\formal{GlobalTheory.RowClaim.SumRatK}
\end{definition}

% ---------------------------------------------------------------------------
\subsection{Main results}\label{sec:pf:rowprops:results}

The hypotheses are the bundle $F=\sfname{ThetaLocalInputs}$ of classical results of \cref{sec:families:theta}, the two
statements of Theorem~K (\cref{def:families:K,thm:theorem-K}), and the three \emph{square statements}: the
convergence and moderate growth of theta integrals for the Weil representations of the grid squares
($\sfname{SquareSt.WeilConvergence}$), the local matching for the two Fock spaces of a square
($\sfname{SquareSt.LocalMatching}$), and the agreement, up to a non-zero constant, of matched theta integrals of a
square ($\sfname{SquareSt.ThetaIntegralsAgree}$).

\begin{theorem}[Row properties of the twisted theta lift]\label{pf:rowprops:rowthm}
Assume $\sfname{ThetaLocalInputs}$, $\sfname{ThetaIntegralsAgree}$, $\sfname{LocalMatching}$ and the three square
statements. Then $\sfname{SumOKK}$, $\sfname{ConstK}$, $\sfname{LinK}$, $\sfname{Formula2K}$ and $\sfname{SumRatK}$ hold
(\cref{def:pf:rowprops:five}).
\formallabel{GlobalTheory.RowClaim.Join.thetaLift_row_properties_conj}
\end{theorem}

The proof occupies \cref{sec:pf:rowprops:assembly,sec:pf:rowprops:const,sec:pf:rowprops:sumrat,sec:pf:rowprops:formula,sec:pf:rowprops:proj,sec:pf:rowprops:flip,sec:pf:rowprops:unfold,sec:pf:rowprops:jets,sec:pf:rowprops:cusp}.
Two of the hypotheses play no role (\cref{pf:rowprops:rem:unused}).

\begin{definition}[Row test at pinned data]\label{def:pf:rowprops:atdat}
The row test of \cref{sec:families:theta} is the statement $\sfname{RowTest}(X)$ for the complex realisation data $X$
of a real-place datum: there are finitely many linear forms on the space of $X$ which jointly separate points and take
values in $\rowpropsemb(\rowpropsLL)$ on all theta vectors of rational points. By $\sfname{AtDat}(\sfname{RowTest})$ we mean
that for every pinned starting configuration of the construction of \cref{sec:families:nonempty} and every real-place
datum $z$ at it satisfying the group condition, the existence condition and the geometric condition, $\sfname{RowTest}$
holds for the complex realisation data of $(z.R,z.F)$.
\formal{Formalism.IndefR.Rat.RowTest}\formal{GlobalTheory.Mem.Top.Bl.AtDat}
\end{definition}

\begin{theorem}[\cref{prop:families:rowrat}]\label{pf:rowprops:main}
Assume the existence of Siegel sets for unitary groups ($\sfname{SiegelSetUnitary}$), Howe duality at the real place
($\sfname{HoweDualityReal}$), Green's theorem on cuspidal characters ($\sfname{GreenCuspidal}$), the bundle
$F=\sfname{ThetaLocalInputs}$, the statement $\sfname{JoinOut}$ (toric denominators, lattice-norm transport, twisted
toric transport, zeta denominators and a reflection statement), the functional equation of Hecke $L$-functions of CM
fields ($\sfname{HeckeTateFunctionalEquation}$), $\sfname{ThetaIntegralsAgree}$, $\sfname{LocalMatching}$ and
$\sfname{FunctionalsSeparate}$. Then $\sfname{AtDat}(\sfname{RowTest})$ holds. Consequently $\sfname{RowRat}$ holds.
\formallabel{GlobalTheory.T.test_vector_assembly_holds}
\formallabel{GlobalTheory.RowTestAssembly.test_vector_at_datum_of_join}
\end{theorem}

\begin{proof}
We collect the inputs.
\begin{enumerate}[label=(\arabic*)]
\item By \lref{GlobalTheory.WeilConv.twisted_weil_representation_convergence}, $\sfname{SiegelSetUnitary}$ and $F$
imply $\sfname{SquareSt.WeilConvergence}$.
\item By \lref{GlobalTheory.SquareLM.twisted_square_local_matching}, $\sfname{HoweDualityReal}$, $F$ and Theorem~K
imply $\sfname{SquareSt.LocalMatching}$.
\item By \lref{GlobalTheory.SquareAgree.twisted_theta_integrals_agree}, $F$ and Theorem~K imply
$\sfname{SquareSt.ThetaIntegralsAgree}$.
\item By \cref{pf:rowprops:rowthm}, (1)--(3), $F$ and Theorem~K give $\sfname{SumOKK}$, $\sfname{ConstK}$, $\sfname{LinK}$,
$\sfname{Formula2K}$ and $\sfname{SumRatK}$.
\item The functional equation implies its form for absolute values
(\lref{GlobalTheory.Published.HeckeTateFunctionalEquation.abs}); completed Hecke $L$-functions are bounded in
vertical strips (\lref{GlobalTheory.HeckeStrip.completedLBoundedInStrips}); and a field datum exists
(\lref{GlobalTheory.SqTop.fieldDat}). By \lref{GlobalTheory.SqTop.lNe_of_functionalEquation} these three
statements give $\sfname{LNe}$: at large scale, for all squares $(i,j)$ the Hecke $L$-functions of the characters
$\nu_H(t,i,j)$ do not vanish at the integers $s$ with $1-e_0(i,j)\le s\le h(i,j)$ (a half rank), and those of the
family at $\mu_0$ do not vanish at the integers $s\ge1$.
\item By \lref{GlobalTheory.SqTop.triEx_holds_alt}, $\sfname{HoweDualityReal}$, $F$, Theorem~K and $\sfname{LNe}$
give $\sfname{TriEx}$ (every square has a bottom triple with constants), and by \lref{GlobalTheory.SqTower.tower_alt_bis}
$F$, Theorem~K and $\sfname{LNe}$ give the Witt-tower norm formula $\sfname{Tower}$, which expresses the ratio of the square
constants at $t$ and at $\mu_0$, multiplied by a ratio of Hecke $L$-values, as an away-unit of $\rowpropsLL$ times step
ratios times the ratio of the bottom constants times another ratio of $L$-values. By
\lref{GlobalTheory.RowClaim.twisted_local_factor_ratio_holds}, $F$, Theorem~K, $\sfname{ConstK}$, $\sfname{LNe}$,
$\sfname{TriEx}$ and $\sfname{Tower}$ give $\sfname{RatioK}$: at large scale, if the field condition holds then
$c_{ijt}/c_{ij\mu_0}\in\rowpropsemb(\rowpropsLL)$ for all $i,j,t$.
\item By \lref{GlobalTheory.PairLift.thetaPairLift_cuspidal_of_hypotheses}, $F$ implies that the theta lifts along
the arrows of the grids carry vector-valued cusp forms of the multiplicity spaces to such cusp forms
($\sfname{LiftsCuspidal}$); by \lref{GlobalTheory.RowEx.twisted_linear_bound}, $F$ and $\sfname{LiftsCuspidal}$ give
$\sfname{LinB}$, the convergence of all pairings $\psi_K$ for choices satisfying the admissibility conditions other than the
real pairing condition; and by \lref{GlobalTheory.RowEx.injK_of_separates_linB}, $\sfname{FunctionalsSeparate}$ and
$\sfname{LinB}$ give $\sfname{InjK}$: at large scale, for admissible choices, $\psi_K(a,u)=0$ for all $a$ implies $u=0$.
\item By \lref{GlobalTheory.T.theta_lift_extension_field_exists}, $\sfname{HoweDualityReal}$,
$\sfname{LocalMatching}$, $F$, $\sfname{JoinOut}$, $\sfname{HeckeTateFunctionalEquation}$, $\sfname{ThetaIntegralsAgree}$,
(1)--(3) and $\sfname{FunctionalsSeparate}$ give $\sfname{ExField}$: at large scale, for every pre-datum, character
frame and finite set $\Xi$ of Hecke characters there is a finite set $S$ of algebraic numbers such that every realisation
in Hida setting with this frame, with twist places off the auxiliary prime and with twisting characters in $\Xi$, whose
coefficient field contains $S$ (via $\rowpropsemb$), admits an admissible choice at the given scale satisfying the field
condition.
\end{enumerate}
Finally, \lref{GlobalTheory.Twist.atDat_rowTest_of_exField} states that $\sfname{GreenCuspidal}$, the rapid decay of
cusp forms (a field of $F$), $\sfname{LinK}$, $\sfname{InjK}$, $\sfname{ConstK}$, $\sfname{SumOKK}$, $\sfname{Formula2K}$,
$\sfname{SumRatK}$, $\sfname{RatioK}$ and $\sfname{ExField}$ imply $\sfname{AtDat}(\sfname{RowTest})$. This proves the
first assertion. The second follows from \lref{GlobalTheory.Mem.LiteralAutomorphy.rowRatW_of_rowTest}, which states
that $\sfname{AtDat}(\sfname{RowTest})$ implies $\sfname{RowRatW}$; the definitions of $\sfname{RowRatW}$ and of
$\sfname{RowRat}$ have the same body.
\end{proof}

\begin{remark}\label{pf:rowprops:rem:partI}
The hypotheses $\sfname{SiegelSetUnitary}$ and $\sfname{HoweDualityReal}$ of \cref{pf:rowprops:main} are also fields
of the bundle $F$, and $\sfname{GreenCuspidal}$ is a field of the bundle $P$; they are listed separately in the formal
statement, as are $\sfname{JoinOut}$ and $\sfname{HeckeTateFunctionalEquation}$. \Cref{prop:families:rowrat} summarises
the hypotheses as ``$P$, $F$, Theorem~K and $\sfname{FunctionalsSeparate}$''; the precise list is the one above.
\end{remark}

% ---------------------------------------------------------------------------
\subsection{Proof of the main theorem}\label{sec:pf:rowprops:assembly}

\begin{proof}[Proof of \cref{pf:rowprops:rowthm}]
Assume the six hypotheses.
\begin{enumerate}[label=(\arabic*)]
\item $\sfname{ConstK}$ holds by \cref{pf:rowprops:constK}.
\item By \lref{GlobalTheory.T.twisted_thetaLift_cuspidal}, $F$ implies $\sfname{LiftsCuspidal}$ (the cuspidality
statement recalled in item (7) of the proof of \cref{pf:rowprops:main}).
\item By \lref{GlobalTheory.T.sumOKK_of_theta_local_inputs}, $F$ implies $\sfname{SumOKK}$: at large scale the
class sums of all admissible choices are meaningful.
\item $\sfname{LinK}$ follows from $F$ and $\sfname{LiftsCuspidal}$ by \cref{pf:rowprops:linK}.
\item By \cref{pf:rowprops:proj}, \cref{pf:rowprops:flipmain}, \cref{pf:rowprops:unfoldmain} and \cref{pf:rowprops:jetsmain},
the four steps $\sfname{StepProj}$, $\sfname{StepFlip}$, $\sfname{StepUnfold}$ and $\sfname{StepJets}$ of
\cref{def:pf:rowprops:steps} hold; here $\sfname{StepProj}$ uses $F$; $\sfname{StepFlip}$ uses $F$, $\sfname{LiftsCuspidal}$,
$\sfname{SquareSt.WeilConvergence}$, $\sfname{ThetaIntegralsAgree}$ and $\sfname{ConstK}$; $\sfname{StepUnfold}$ uses $F$,
$\sfname{LiftsCuspidal}$ and $\sfname{SumOKK}$; and $\sfname{StepJets}$ uses $F$ and $\sfname{SumOKK}$. By
\cref{pf:rowprops:formula2K} the four steps, $\sfname{ConstK}$ and $\sfname{SumOKK}$ give $\sfname{Formula2K}$.
\item $\sfname{SumRatK}$ follows from $\sfname{SumOKK}$ by \cref{pf:rowprops:sumrat}.\qedhere
\end{enumerate}
\end{proof}

% ---------------------------------------------------------------------------
\subsection{The square constants}\label{sec:pf:rowprops:const}

\begin{definition}[Square standing]\label{def:pf:rowprops:standing}
For a grid datum and $t\in T$, the \emph{square standing condition} says that the splitting bound of the host is at
least $11$ and that no place $w\in\mathrm{loc}(t)$ divides $Q_S\,q\,d\,\Nm(\mathfrak c)$, where $Q_S$, $q$, $d$ and the
ideal $\mathfrak c$ are invariants of the host.
\formal{GlobalTheory.FunObj.SquareStanding}
\end{definition}

When the square standing condition holds, every square $(i,j)$ is an instance of the configuration of complementary
spaces of Theorem~K: by \lref{GlobalTheory.T.squareThetaTripleHypotheses_of_standing}, for every volume datum $M$ of
level $t$ there are a base, geometric data, characters and Weil representations forming an admissible triple realising
the square, with the auxiliary split place of the triple outside $\mathrm{loc}(t)$.

\begin{lemma}[Standing]\label{pf:rowprops:standing}
\begin{enumerate}[label=(\alph*)]
\item At large scale, for every admissible choice and every place $w$ of the base set $W_0=\{w(v):v\in S_0\}$, the integer
$Q_S\,q\,d\,\Nm(\mathfrak c)$ does not lie in $w$.
\item At large scale, for every admissible choice and every $t\in T$, the square standing condition holds.
\end{enumerate}
\formallabel{GlobalTheory.PlacesPrime.places_prime_to_level}
\formallabel{GlobalTheory.T.test_places_prime_to_level}
\formallabel{GlobalTheory.SmallConv.standing}
\formallabel{GlobalTheory.BigNZ.standing}
\formallabel{GlobalTheory.SmallConv.splitBound_ge}
\formallabel{GlobalTheory.BigNZ.splitBound_ge}
\end{lemma}

\begin{proof}
(a) Let $w=w(v)$ with $v\in S_0$. There is a rational prime $p$ lying in $v$ and coprime to $N=Q_S\,q\,d\,\Nm(\mathfrak c)$:
for $v=v_0$ take the auxiliary prime $q_0$ over which $v_0$ lies, which is coprime to the invariants of the host by its
construction; for the other places of $S_0$, which lie over the primes $g_0,g_1,p_{de},p'_{de}$ of the level, the size
conditions of the Hida setting make these primes coprime to $N$. Since $w$ lies over $v$, $p\in w$; as $w$ is prime and
$\gcd(p,N)=1$, $N\notin w$.

(b) The bound $11\le$ (splitting bound) is one of the size conditions. Let $w\in\mathrm{loc}(t)$. If $w\in W_0$ use (a).
Otherwise $w$ is a place of the twist $t$; the level character conditions give $Q_S\,q\,d\notin w$ and the guard condition
of the Hida setting gives $\Nm(\mathfrak c)\notin w$, so the product is not in the prime ideal $w$.
\end{proof}

\begin{lemma}[Integrability of the big route]\label{pf:rowprops:bigint}
Assume $\sfname{SquareSt.WeilConvergence}$. At large scale (with $x_0=12^{5/\eps}$), for every admissible choice, all
$i,j<m$, $t$, volume data $M$ of level $t$, matched pairs $(\Theta,\Theta')$ and all $g,g_1$, the big route kernel
$\rowpropsKbig_\Theta(g,g_1)$ is integrable on $\rowpropsfd_{i+1,j+1}$.
\formallabel{GlobalTheory.SquareConv.square_convolution_large_bound}
\formallabel{GlobalTheory.T.square_convolution_estimate}
\end{lemma}

\begin{proof}
By the size conditions the square standing condition holds (\lref{GlobalTheory.SquareConv.standing}, which is the
argument of \cref{pf:rowprops:standing}(b)), so the square is realised by an admissible triple
(\lref{GlobalTheory.T.squareThetaTripleHypotheses_of_standing}). For such squares,
\lref{GlobalTheory.T.bigRoute_kernel_integrableOn_unitary_quotient} states that, under
$\sfname{SquareSt.WeilConvergence}$, the big route kernel of a matched pair is integrable on the fundamental domain of
$G_{V_{i+1,j+1}}$, for all $g,g_1$.
\end{proof}

\begin{lemma}[Integrable translates of the small route]\label{pf:rowprops:smallint}
Assume $\sfname{ThetaIntegralsAgree}$. At large scale, for every admissible choice, all $i,j,t$, volume data $M$ of
level $t$, matched pairs $(\Theta,\Theta')$ and all $g,g_1$, there is a finite combination $z$ with $\rowpropsdeg z\neq0$ such
that $z\cdot\rowpropsKsmall_{\Theta'}(g,g_1)$ is integrable on $\rowpropsfd_{i,j}$.
\formallabel{GlobalTheory.SmallConv.small_convolution_exists_integrable_translates}
\formallabel{GlobalTheory.T.exists_integrable_translates_small_convolution}
\end{lemma}

\begin{proof}
By \cref{pf:rowprops:standing} and \lref{GlobalTheory.T.squareThetaTripleHypotheses_of_standing}, the square is
realised by an admissible triple $(\mathrm{Bs},\Gamma',X,\omega,\omega')$ with a block embedding
$\mathrm{blk}:G_{V_{i,j+1}}\times G_{V_{i+1,j}}\to\Uni(W)$ into the unitary group of the skew-hermitian space of the triple.
By \lref{GlobalTheory.T.translatedKernel_le_thetaFunction_routeBigVertex} there are $z$ with $\rowpropsdeg z\ne0$, a
Schwartz function $\Phi'$ in the class for which the theta integrals of the triple converge, a map $u\mapsto a(u)$ and a
constant $C_0$ with
\[
\bigl|z\cdot\rowpropsKsmall_{\Theta'}(g,g_1)(u)\bigr|\le C_0\,\bigl|\theta_{\omega'}(\Phi')\bigl(a(u),\mathrm{blk}(g,g_1)\bigr)\bigr|
\qquad\text{for all }g,u .
\]
The first component of $\sfname{ThetaIntegralsAgree}$ gives the moderate growth of these theta functions: there are $C,k$
with $|\theta_{\omega'}(\Phi')(x,b)|\le C\,\mathrm{ht}(b)^k$ for all $x,b$, $\mathrm{ht}$ being a height on $\Uni(W)(\A)$.
Choose $\gamma\in G_{V_{i,j+1}}(F)$ with $\mathrm{ht}(\gamma g)\le 2\,\mathrm{ht}_{\mathrm{quot}}(g)$
(\lref{GlobalTheory.Grow.exists_height_le_two_mul_quotHeight}). The small route kernel is invariant under rational
left translation of $g$ (\lref{GlobalTheory.ZConvLeft.kerSmall_rat_left}), and the height of a block element is at
most the product of the heights (\lref{GlobalTheory.BlockHeight.ht_blkU_le}). Hence
\[
\bigl|z\cdot\rowpropsKsmall_{\Theta'}(g,g_1)(u)\bigr|\le |C_0|\,|C|\,\bigl(2\,\mathrm{ht}_{\mathrm{quot}}(g)\,\mathrm{ht}(g_1)\bigr)^k
\]
uniformly in $u$. The function $u\mapsto z\cdot\rowpropsKsmall_{\Theta'}(g,g_1)(u)$ is continuous
(\lref{GlobalTheory.SmallRoute.continuous_translates_kerSmall}), and the fundamental domain $\rowpropsfd_{i,j}$ has
finite measure (finite volume of $G_{V_{i,j}}(F)\backslash G_{V_{i,j}}(\A)$, \lref{GlobalTheory.CuspDecay.measure_lt_top},
\cite{Borel1963}). A bounded continuous function on a set of finite measure is integrable.
\end{proof}

\begin{lemma}[A non-vanishing big route integral]\label{pf:rowprops:bignz}
Assume $F$, $\sfname{ThetaIntegralsAgree}$ and $\sfname{LocalMatching}$. At large scale, for every admissible choice and
all $i,j,t$ there are a volume datum $M$ of level $t$, polynomial tensors $a$ along the big route, local data $(n,\Phi)$
at the places of $\mathrm{loc}(t)$ which are natural data of depth at least the local depths, finite tensors
$\tau,\tau'$ along the two routes having the local flip property for $(n,\Phi)$, and $g,g_1$, such that
\[
\int_{\rowpropsfd_{i+1,j+1}}\rowpropsKbig_{\Theta}(g,g_1)\,d\nu\ne0,\qquad \Theta=\mathrm{tens}(a,\tau).
\]
\formallabel{GlobalTheory.BigNZ.large_term_nonzero}
\formallabel{GlobalTheory.T.dominant_term_ne_zero}
\end{lemma}

\begin{proof}
By \cref{pf:rowprops:standing} the square standing condition holds. \Lref{GlobalTheory.BigNZ.exists_integral_ne_zero}
states exactly the conclusion for every grid datum, Weil datum, frames, real match and square satisfying the square
standing condition, given strong approximation for special unitary groups (a field of $F$),
$\sfname{ThetaIntegralsAgree}$ and $\sfname{LocalMatching}$. (Its proof realises the square by an admissible triple, takes
the matched data and the non-zero constant of Theorem~K, and produces a non-zero theta integral with local components of
the required kind at the places of $\mathrm{loc}(t)$, using strong approximation at a split place.)
\end{proof}

\begin{proposition}[Non-vanishing of the square constants]\label{pf:rowprops:constK}
Assume $F$, $\sfname{ThetaIntegralsAgree}$, $\sfname{LocalMatching}$ and the three square statements. Then
$\sfname{ConstK}$ holds: at large scale, for every admissible choice and all $i,j<m$ and $t\in T$, the square $(i,j)$
is non-zero at $(i,j,t)$ and $c_{ijt}\ne0$.
\formallabel{GlobalTheory.ConstNZ.constK_alt}
\end{proposition}

\begin{proof}
Combine \cref{pf:rowprops:bigint,pf:rowprops:smallint,pf:rowprops:bignz} (\cref{def:pf:rowprops:big}). Fix an admissible
choice at large scale and $(i,j,t)$. \Cref{pf:rowprops:bignz} gives $M$, $a$, $\tau$, $\tau'$, $(n,\Phi)$, $g$, $g_1$ with
$\int\rowpropsKbig_{\mathrm{tens}(a,\tau)}(g,g_1)\ne0$. Since $\mathrm{rel}_{ij}$ projects onto its first factor, there is a
polynomial tensor $b'$ along the small route with $(a,b')\in\mathrm{rel}_{ij}$ (after identifying finite combinations
of pure polynomial tensors with elements of the tensor product), and $b'$ comes from a finite combination $b$ of pure
tensors. Then $\Theta=\mathrm{tens}(a,\tau)$ and $\Theta'=\mathrm{tens}(b,\tau')$ form a matched pair at $(i,j,t)$.
\Cref{pf:rowprops:smallint} gives $z$ with $\rowpropsdeg z\ne0$ and $z\cdot\rowpropsKsmall_{\Theta'}(g,g_1)$ integrable,
and \cref{pf:rowprops:bigint} gives the integrability of $\rowpropsKbig_\Theta(g,g_1)$. Since the choice is admissible,
$c_{ijt}$ is a square constant at $(i,j,t)$, and \eqref{eq:pf:rowprops:sqconst} gives
\[
\rowpropsdeg z\int\rowpropsKbig_\Theta(g,g_1)\,d\nu=c_{ijt}\int z\cdot\rowpropsKsmall_{\Theta'}(g,g_1)\,d\nu .
\]
The left side is non-zero. Hence $c_{ijt}\ne0$, and the integral on the right is non-zero, which shows that the square is
non-zero at $(i,j,t)$ (with the witnesses $M$, $(\Theta,\Theta')$, $z$, $g$, $g_1$).
\end{proof}

% ---------------------------------------------------------------------------
\subsection{Convergence and rationality of the class sums}\label{sec:pf:rowprops:sumrat}

\begin{proposition}[Convergence]\label{pf:rowprops:linK}
Assume $F$ and $\sfname{LiftsCuspidal}$. Then $\sfname{LinK}$ holds (with $C_{F0}=0$ and $x_0=0$).
\formallabel{GlobalTheory.RowTest.twist_linear_independence_test}
\formallabel{GlobalTheory.T.twisted_characters_linearly_independent}
\formallabel{GlobalTheory.T.twisted_theta_lift_automorphic}
\formallabel{GlobalTheory.twisted_theta_lift_automorphic_of_convergent}
\end{proposition}

\begin{proof}
Let $\mathrm{ch}$ be an admissible choice at a realisation in Hida setting, and let $a$, $u$ be given.
By \lref{GlobalTheory.T.twisted_base_change_lift_exists} (which uses $F$ and $\sfname{LiftsCuspidal}$), the row form
$\mathrm{form}(u)$ is convergent along the row path $\rho_0$ with the Fock vectors of $(\varrho,\Phi_a)$, and the row lift
$\rowpropsth(\varrho,\Phi_a;\mathrm{form}(u))$ is an automorphic function on $G_{V_{m,0}}(\A)$. The column function $\Psi_a$ is
a column (left equivariant and right $K_+$-invariant), and $M$ is of row type. By
\lref{GlobalTheory.T.besselPeriodCorank_converges_column}, the Bessel pairing \eqref{eq:pf:rowprops:bessel} of the
scalar function of a column with an automorphic function on $G_{V_{m,0}}(\A)$ converges absolutely in the iterated sense.
Both conditions together are the convergence of $\psi_K$ at $(a,u)$.
\end{proof}

\begin{proposition}[Rationality of the class sums]\label{pf:rowprops:sumrat}
$\sfname{SumOKK}$ implies $\sfname{SumRatK}$.
\formallabel{GlobalTheory.RowTest.sumRatK_of_sumOKK}
\formallabel{GlobalTheory.T.rowTestRationalSums_of_finiteSumsMeaningful}
\formallabel{GlobalTheory.RowTest.SumRat.vK_mem}
\formallabel{GlobalTheory.RowTest.SumRat.sAA_mem}
\formallabel{GlobalTheory.RowTest.SumRat.psiC_rat}
\end{proposition}

\begin{proof}
Take the constants of $\sfname{SumOKK}$. Let $\mathrm{ch}$ be admissible at large scale with the field condition, and
fix $t$, $x$, $a$; write $\Phi=\Phi_a$. Since $\rowpropsemb(\rowpropsLL)$ is a subfield of $\C$, it suffices to show that
each factor and summand of
\begin{multline*}
v_K(t,x,a)=s_{AA}\sum_{\beta\in\rowpropsSupp(t,\Phi)}\Bigl(\prod_w\mathrm{coef}_w(\beta_w)\Bigr)\\
\times\sum_{c\in\mathcal C(\beta)}\gamma(c)\,\bigl\langle\psi_C(\Psi_a)(\mathrm{pt}(c)_0),\mathbf{f}(t,j_Xx)(\mathrm{pt}(c)_0)\bigr\rangle_{\tau_\beta,c}
\end{multline*}
lies in it, for a suitable family of expansions $\tau_\beta$ (by $\mathsf{SumOK}$ the value does not depend on this family,
\cref{def:pf:rowprops:classsum}).
\begin{enumerate}[label=(\arabic*)]
\item Each local scalar $s_w$ is an integral power of $\Nm(w)^{1/2}$ (\lref{GlobalTheory.FunObj.aaScalar_eq_pow}, applied
with the conductor property of the additive character at the places of $P_{AA}$), and $\Nm(w)^{1/2}\in\rowpropsemb(\rowpropsLL)$
by the field condition; hence $s_{AA}\in\rowpropsemb(\rowpropsLL)$.
\item For every $\beta$ in the support, which has an expansion by $\mathsf{SumOK}$, there is a \emph{coset expansion at
expansion depth} (\lref{GlobalTheory.RowTest.exists_cosetExpansion_expDepth}); it is an expansion, and each of its
coefficients $\tau_\beta(d)$ is a product over $w$ of values $\mathrm{datum}_w(\beta_w)(X_w,Y_w)$ of the local data of the
terms. These values lie in $\rowpropsemb(\rowpropsLL)$ (\lref{GlobalTheory.T.unfoldedFlipTensor_values_mem_coeffField}).
\item The coefficients $\mathrm{coef}_w(\beta_w)$ are rational, and the weights $\gamma(c)$ are products of inverses of
natural numbers.
\item The vector $\psi_C(\Psi_a)(g)$ has coefficients in $\rowpropsemb(\rowpropsLL)$, being the image under $\rowpropsemb$ of
a vector over $\rowpropsLL$. By \lref{GlobalTheory.T.topRow_normalisedKernelChain_preserves_coeffs}, the column map
$\mathrm{J}^{\mathrm{col}}_d(\mathrm{pt}(c))$ of a datum $d$ in the support of a coset expansion at expansion depth preserves the
property ``coefficients in $\rowpropsemb(\rowpropsLL)$''.
\item The vector $\mathbf{f}(t,j_Xx)(y)$ has coordinates in $\rowpropsemb(\rowpropsLL)$ for every finite point $y$
(\lref{GlobalTheory.T.rationalForm_coordinates_coeff_mem_coefficientField}), and the row map
$\mathrm{J}^{\mathrm{row}}_d(\mathrm{pt}(c))$ preserves this property (\lref{GlobalTheory.T.lastColumn_normalizedKernelChain_coeffField}).
\item The branching pairing of two vectors with coefficients in $\rowpropsemb(\rowpropsLL)$ lies in $\rowpropsemb(\rowpropsLL)$
(\lref{GlobalTheory.T.branchingPairing_mem_coefficientField}).
\end{enumerate}
By (4)--(6) each term $\mathrm{br}(\dots)$ of $\langle\cdot,\cdot\rangle_{\tau_\beta,c}$ lies in the field, by (2) so do
the coefficients $\tau_\beta(d)$, and by (1), (3) the remaining factors do. Hence $v_K(t,x,a)\in\rowpropsemb(\rowpropsLL)$.
\end{proof}

% ---------------------------------------------------------------------------
\subsection{The lift formula from four steps}\label{sec:pf:rowprops:formula}

Fix an admissible choice. For $t\in T$ and $y\in H_X(t)$ we regard the form $F_{t,y}$ on $\Uni(X')$ as a function on
$G_{V_{0,m}}(\A)$ ($V_{0,m}=X'$), and we write $\phi_a$ for the scalar column functions.

\begin{definition}[Flipped finite tensors; definite periods]\label{def:pf:rowprops:flipfin}
\begin{enumerate}[label=(\alph*)]
\item A finite combination $\tau$ of pairs $d=(d_1,d_2)$, with $d_1=(d_{1,i})_{i<m}$, $d_{1,i}\in\rowpropsS_{\chi(i,m)}$, and
$d_2=(d_{2,j})_{j<m}$, $d_{2,j}\in\rowpropsS_{\mu(t,m,j)}$, is a \emph{flipped finite tensor} for $(t,\Phi)$ if there are an
idempotent $e$ of the finite adeles which is a sum of local idempotents at the places of $\mathrm{loc}(t)$, and a map
$d\mapsto\varphi(d)$ such that every component of $d$ is a local vector at $e$ relative to the corresponding component
of $\varphi(d)$ and the vacuum vector of the real match, and for all half-integral arguments $x$, $y$
\[
\sum_d\tau(d)\prod_i\varphi(d)_{1,i}(x_i)\prod_j\varphi(d)_{2,j}(y_j)=\prod_{w\in\mathrm{loc}(t)}T^{\mathrm{flip}}_w(\Phi)\bigl(\arg_w(x,y)\bigr),
\]
where $\arg_w(x,y)$ are the local arguments of the opposite path in the frames of the grid choice.
\item For a finite combination $\tau_i$ of pairs $P=(P_1,P_2)$ of families of polynomials on the last row and the last
column, $\mathrm{opp}(\tau_i,\tau)$ is the finite combination of pure tensors along $\rowpropsopp_t$ with components
$\mathrm{fock}(d_{1,i},P_{1,i})$ and $\mathrm{fock}(d_{2,j},P_{2,j})$ and coefficients $\tau_i(P)\tau(d)$.
\item Let $\iota_{\mathrm t}:G_{V_{m,m}}\to G_{V_{\mathrm t_m}}$ be the stabiliser embedding (\cref{def:pf:rowprops:grid}(c)).
For a pair $P=(P_1,P_2)$ of families of polynomials on the column $\kappa^t$ and on the last row $\rho_m$, a pure finite
datum $d$ along $\kappa^t$ and $\rho_m$ (\cref{def:pf:rowprops:classsum}), and functions $\phi$ on $G_{V_{\mathrm
t_0}}(\A)$ and $F$ on $G_{V_{0,m}}(\A)$, the \emph{definite period} is
\begin{equation}\label{eq:pf:rowprops:defper}
\rowpropsdef(M;P,d;\phi,F)=\int_{\rowpropsfd_{m,m}}\rowpropsL_{\kappa^t}\bigl(\mathrm{fock}(\tilde d_{1,j},P_{1,j})\bigr)\phi\,\bigl(\iota_{\mathrm t}(h)\bigr)\;
\rowpropsL_{\rho_m}\bigl(\mathrm{fock}(\tilde d_{2,i},P_{2,i})\bigr)F\,(h)\,d\nu_{m,m}(h),
\end{equation}
where $\tilde d$ are the vectors obtained from $d$ by the normal datum of the grid choice (``moved vectors''); for finite
combinations $\sigma$ of pairs $P$ and $\tau$ of data $d$, $\rowpropsdef(M;\sigma,\tau;\phi,F)=\sum_{P,d}\sigma(P)\tau(d)
\rowpropsdef(M;P,d;\phi,F)$.
\item For $\beta\in\rowpropsSupp(t,\Phi)$, a volume datum $M$ is a \emph{term volume} for $\beta$ if its measures at the
nodes $\mathrm t_j$ ($j<m$) and $(i,m)$ ($i\le m$) are normalised at the levels determined by the terms of $\beta$ and the
frames (and are Haar measures with fundamental domains). Term volumes exist, since Haar measures and fundamental domains
exist at every open compact level.
\end{enumerate}
\formal{GlobalTheory.RowClaim.IsFlipFin}\formal{GlobalTheory.RowClaim.oppTens}\formal{GlobalTheory.RowClaim.defTot}\formal{GlobalTheory.RowClaim.IsTermVol}
\end{definition}

The spaces $V_{m,m}$, $V_{i,m}$ and $V_{\mathrm t_j}$ are totally definite, so \eqref{eq:pf:rowprops:defper} is an integral
over a compact quotient.

\begin{definition}[The four steps]\label{def:pf:rowprops:steps}
For an admissible choice $\mathrm{ch}$ consider the following properties.
\begin{enumerate}[label=(\alph*)]
\item $\mathsf{ProjAt}$: for all $t$ and all good $x\in\mathrm{FR}(t)$, $e_0\bigl(\mathrm{lift}_t(j_Xx)\bigr)=\mathrm{lift}_t(j_Xx)$.
\item $\mathsf{FlipAt}$: there are a finite combination $\tau_i$ of pairs of polynomial families on the last row and
column and $z:T\to\C^\times$ such that for every $t$: $z_t\in\rowpropsXi_t$, and for every coset datum $\Phi$ there is a
flipped finite tensor $\tau$ for $(t,\Phi)$ such that for every volume datum $M$ of level $t$ and every $y\in H_X(t)$,
$F_{t,y}$ is convergent along $\rowpropsopp_t$ for $\mathrm{opp}(\tau_i,\tau)$ and
\begin{equation}\label{eq:pf:rowprops:flipat}
\rowpropsth\bigl(\varrho,\Phi;\mathrm{lift}_t(y)\bigr)=\bigl(z_t\,\rowpropsctot(t)\bigr)^{-1}\,
\rowpropsL_{\rowpropsopp_t}\bigl(\mathrm{opp}(\tau_i,\tau)\bigr)F_{t,y}.
\end{equation}
\item $\mathsf{UnfoldAt}$: for every $\tau_i$ there are a finite combination $\sigma$ of pairs of polynomial families on the
column $\kappa$ and the last row, and $z':T\to\C^\times$, such that for every $t$: $z'_t\in\rowpropsXi_t$, and for every $a$,
every flipped finite tensor $\tau$ for $(t,\Phi_a)$, every $M$ of level $t$, all families $(M_\beta,\tau_\beta)$ of term
volumes and expansions indexed by $\beta\in\rowpropsSupp(t,\Phi_a)$, and all $y$,
\begin{equation}\label{eq:pf:rowprops:unfoldat}
\rowpropsB\Bigl(\phi_a,\rowpropsL_{\rowpropsopp_t}\bigl(\mathrm{opp}(\tau_i,\tau)\bigr)F_{t,y}\Bigr)=
(z'_t)^{-1}\,s_{AA}\sum_{\beta\in\rowpropsSupp(t,\Phi_a)}\Bigl(\prod_w\mathrm{coef}_w(\beta_w)\Bigr)\,
\rowpropsdef\bigl(M_\beta;\sigma,\tau_\beta;\phi_a,F_{t,y}\bigr).
\end{equation}
\item $\mathsf{JetsAt}$: every $\beta$ has a term volume, and for every $\sigma$ there is $\kappa\in\C$ such that for all
$t$, $a$, $\beta\in\rowpropsSupp(t,\Phi_a)$, term volumes $M$, expansions $\tau$ of $\beta$ and all $y$,
\begin{equation}\label{eq:pf:rowprops:jetsat}
\rowpropsdef\bigl(M;\sigma,\tau;\phi_a,F_{t,y}\bigr)=\kappa\cdot\rowpropsclass_{\Phi_a}\bigl(\beta,\tau;\psi_C(\Psi_a),\mathbf{f}(t,y)\bigr).
\end{equation}
\end{enumerate}
$\sfname{StepProj}$, $\sfname{StepFlip}$, $\sfname{StepUnfold}$ and $\sfname{StepJets}$ are the statements
$\rowpropsBig(\mathsf{ProjAt})$, $\rowpropsBig(\mathsf{FlipAt})$, $\rowpropsBig(\mathsf{UnfoldAt})$, $\rowpropsBig(\mathsf{JetsAt})$.
\formal{GlobalTheory.RowClaim.ProjAt}\formal{GlobalTheory.RowClaim.FlipAt}
\formallabel{GlobalTheory.RowClaim.UnfoldAt}\formallabel{GlobalTheory.RowClaim.JetsAt}
\formallabel{GlobalTheory.RowClaim.StepFlip}
\end{definition}

The essential point is that the constants $\tau_i$, $\sigma$ and $\kappa$ in (b)--(d) do not depend on $t$, $a$, $\beta$ or
$y$; this is what makes $\kappa$ in the lift formula universal.

\begin{lemma}[Values of the twisting character]\label{pf:rowprops:xivals}
Let $t\in T$ and suppose that $\xi_t^n=1$ for some $n\ge1$ and that all local values of $\xi_t$ lie in
$\rowpropsemb(\rowpropsLL)$. Then every $z\in\rowpropsXi_t$ satisfies $z^n=1$ and $z\in\rowpropsemb(\rowpropsLL)$.
\formallabel{GlobalTheory.RowClaim.FormulaTwo.mem_xiVals}
\end{lemma}

\begin{proof}
Induction on the construction of $\rowpropsXi_t$ as the subgroup generated by the local values: a local value $z=\xi_{t,w}(x)$
satisfies $z^n=\xi^n_{t,w}(x)=1$ and lies in the field by hypothesis; $1$ has both properties; and both properties are
preserved by products and inverses, since $\rowpropsemb(\rowpropsLL)$ is a subfield.
\end{proof}

\begin{proposition}[The lift formula]\label{pf:rowprops:formula2K}
$\sfname{StepProj}$, $\sfname{StepFlip}$, $\sfname{StepUnfold}$, $\sfname{StepJets}$, $\sfname{ConstK}$ and $\sfname{SumOKK}$
imply $\sfname{Formula2K}$.
\formallabel{GlobalTheory.T.rowTestLiftFormula_principalLevel}
\formallabel{GlobalTheory.RowClaim.FormulaTwo.formula2K}
\formallabel{GlobalTheory.RowClaim.FormulaTwo.formulaAt_steps}
\end{proposition}

\begin{proof}
By \cref{def:pf:rowprops:big} we may take one set of constants for the six statements. Let $\mathrm{ch}$ be an
admissible choice at large scale with the field condition, at a realisation in Hida setting; then $R$ is a chain and its
level characters satisfy the level conditions at $\mu_0$, in particular every $\xi_t$ has finite order. Let
$\tau_i$, $z$ be given by $\mathsf{FlipAt}$, let $\sigma$, $z'$ be given by $\mathsf{UnfoldAt}$ applied to $\tau_i$, and let
$\kappa$ be given by $\mathsf{JetsAt}$ applied to $\sigma$.

\emph{The roots of unity.} Fix $t$ and $n\ge1$ with $\xi_t^n=1$. The field condition puts the local values of $\xi_t$ in
$\rowpropsemb(\rowpropsLL)$, and $z_tz'_t\in\rowpropsXi_t$. By \cref{pf:rowprops:xivals} there is $\zeta_t\in\rowpropsLL$
with $\rowpropsemb(\zeta_t)=z_tz'_t$, and $\rowpropsemb(\zeta_t^n)=(z_tz'_t)^n=1$, hence $\zeta_t^n=1$ by injectivity of
$\rowpropsemb$. We show $\mathsf{FormulaAt}(\mathrm{ch},\kappa,\zeta)$.

\emph{The formula.} Fix $t$, a good point $x\in\mathrm{FR}(t)$ and $a$; put $y=j_X(t)x$ and $\Phi=\Phi_a$. Let $\tau$ be the
flipped finite tensor for $(t,\Phi)$ given by $\mathsf{FlipAt}$, $M$ a volume datum of level $t$
(\cref{def:pf:rowprops:weil}), $M_\beta$ term volumes and $\tau_\beta$ expansions for $\beta\in\rowpropsSupp(t,\Phi)$ (they
exist by $\mathsf{JetsAt}$ and $\mathsf{SumOK}$). Then:
\begin{enumerate}[label=(\roman*)]
\item By \eqref{eq:pf:rowprops:formTh} and $\mathsf{ProjAt}$, $\mathrm{form}(\Theta_t(y))=e_0(\mathrm{lift}_t(y))=\mathrm{lift}_t(y)$, so
$\psi_K(a,\Theta_t(y))=\rowpropsB\bigl(\phi_a,\rowpropsth(\varrho,\Phi;\mathrm{lift}_t(y))\bigr)$.
\item By \eqref{eq:pf:rowprops:flipat} and the homogeneity of $\rowpropsB$ in its second argument,
$\psi_K(a,\Theta_t(y))=(z_t\rowpropsctot(t))^{-1}\,b$, with $b=\rowpropsB\bigl(\phi_a,\rowpropsL_{\rowpropsopp_t}(\mathrm{opp}(\tau_i,\tau))F_{t,y}\bigr)$.
\item By \eqref{eq:pf:rowprops:unfoldat}, $b=(z'_t)^{-1}s_{AA}\,u$ with $u=\sum_\beta(\prod_w\mathrm{coef}_w(\beta_w))\,
\rowpropsdef(M_\beta;\sigma,\tau_\beta;\phi_a,F_{t,y})$.
\item By \eqref{eq:pf:rowprops:jetsat} applied termwise, $u=\kappa\,v$ with
\[
v=\sum_\beta\Bigl(\prod_w\mathrm{coef}_w(\beta_w)\Bigr)\rowpropsclass_\Phi\bigl(\beta,\tau_\beta;\psi_C(\Psi_a),\mathbf{f}(t,y)\bigr),
\]
and $v=v_{\mathrm{core}}(t,\Phi;\psi_C(\Psi_a),\mathbf{f}(t,y))$ because $v_{\mathrm{core}}$ does not depend on the
expansions (\cref{def:pf:rowprops:classsum}).
\item By definition $v_K(t,x,a)=s_{AA}\,v$.
\end{enumerate}
By $\sfname{ConstK}$ all $c_{ijt}\neq0$, so $\rowpropsctot(t)\ne0$; also $z_t,z'_t\ne0$. Hence
\[
\rowpropsemb(\zeta_t)\,\rowpropsctot(t)\,\psi_K(a,\Theta_t(y))=z_tz'_t\,\rowpropsctot(t)\,(z_t\rowpropsctot(t))^{-1}(z'_t)^{-1}s_{AA}\,\kappa\,v
=\kappa\,v_K(t,x,a).\qedhere
\]
\end{proof}

% ---------------------------------------------------------------------------
\subsection{Projection}\label{sec:pf:rowprops:proj}

Write $\Theta^{\mathrm{dn}}_t(y)$ for the iterated theta lift down the column attached to $y\in H_X(t)$, viewed as a
linear map from smooth vectors of $M_Y(0)$ to functions on $G_{Y_0}(\A)$, so that $\mathrm{lift}_t(y)=\Theta^{\mathrm{dn}}_t(y)(w(d))$.
The finite group $\Gamma$ acts on the rational points by $r(g)$ and on functions on $G_{Y_0}(\A)$ by right translations
$\mathrm{act}(e_G(g))$ by elements $e_G(g)$ of the finite adelic group; by definition the Hecke projector is
$e_0=\sum_g\rowpropsemb(c_g)\,\mathrm{act}(e_G(g))$, where $e_X=\sum_gc_g\,g$.

\begin{proposition}[Projection]\label{pf:rowprops:proj}
Assume the rapid decay of cusp forms (a field of $F$). Then $\sfname{StepProj}$ holds (with $C_{F0}=0$, $x_0=0$): for
every admissible choice at a realisation in Hida setting, every $t$ and every good $x\in\mathrm{FR}(t)$,
$e_0(\mathrm{lift}_t(j_Xx))=\mathrm{lift}_t(j_Xx)$.
\formallabel{GlobalTheory.T.twisted_period_identity}
\formallabel{GlobalTheory.Adopt.twisted_period_identity}
\formallabel{ThetaProjectionIdempotent.CompanionTheory.Member.twisted_period_identity_local}
\formallabel{ThetaProjectionIdempotent.CompanionTheory.Member.ThetaDownProjection.proj_thDn}
\formallabel{ThetaProjectionIdempotent.CompanionTheory.Member.ThetaDownProjection.sum_act}
\end{proposition}

\begin{proof}
We show $e_0(\Theta^{\mathrm{dn}}_t(j_Xx)(v))=\Theta^{\mathrm{dn}}_t(j_Xx)(v)$ for every smooth vector $v$ and apply it to
$v=w(d)$.

(a) \emph{Equivariance.} By \lref{ThetaProjectionIdempotent.CompanionTheory.Member.ThetaDownProjection.thDn_act}, for
$g\in\Gamma$ and $x'\in\mathrm{FR}(t)$,
\[
\Theta^{\mathrm{dn}}_t\bigl(j_X(r(g)x')\bigr)(v)=\mathrm{act}(e_G(g))\,\Theta^{\mathrm{dn}}_t(j_Xx')(v);
\]
this uses the rapid decay of cusp forms (through a finiteness statement for the multiplicity spaces of the column) and
the core condition.

(b) \emph{Good points.} Since $x$ is good, $x=e_Xx'$ for some $x'$ (the level idempotent at the trivial subgroup is $1$,
\lref{Formalism.L2.DIV.level_idempotent_bot}). As $j_X$ is $\rowpropsemb$-semilinear and $\Theta^{\mathrm{dn}}_t$ is
linear, (a) gives
\[
\Theta^{\mathrm{dn}}_t(j_Xx)(v)=\sum_g\rowpropsemb(c_g)\,\mathrm{act}(e_G(g))\,\Theta^{\mathrm{dn}}_t(j_Xx')(v)
=e_0\bigl(\Theta^{\mathrm{dn}}_t(j_Xx')(v)\bigr).
\]

(c) \emph{Idempotence on the image.} For every $y$, $e_0(e_0(\Theta^{\mathrm{dn}}_t(y)(v)))=e_0(\Theta^{\mathrm{dn}}_t(y)(v))$. By
linearity it suffices to take $v$ a pure tensor $f$. By the chain property, $e_0(\Theta^{\mathrm{dn}}_t(y)(f))$ is the value at
$f$ of the form $\Theta_t(y)$ of the multiplicity space $H_E$ (as in \cref{pf:rowprops:formTh}); and the forms of $H_E$ take
values in the $e_0$-fixed space, which is one of the conditions on the multiplicity space in the core condition. Hence
$e_0$ fixes $e_0(\Theta^{\mathrm{dn}}_t(y)(f))$.

By (b) and (c), $e_0(\Theta^{\mathrm{dn}}_t(j_Xx)(v))=e_0(e_0(\Theta^{\mathrm{dn}}_t(j_Xx')(v)))=e_0(\Theta^{\mathrm{dn}}_t(j_Xx')(v))
=\Theta^{\mathrm{dn}}_t(j_Xx)(v)$.
\end{proof}

% ---------------------------------------------------------------------------
\subsection{Commuting squares and the flip}\label{sec:pf:rowprops:flip}

\begin{definition}[Cuspidal listed functions; commuting squares; cusp steps]\label{def:pf:rowprops:commutes}
\begin{enumerate}[label=(\alph*)]
\item A function $F$ on $G_{V_X}(\A)$ is \emph{cuspidal listed} at the node $X$ if it is a smooth automorphic function,
cuspidal, and of the archimedean type prescribed for $X$ at every real place (for the nodes $(i,j)$, its complex
conjugate is of that type).
\item A number $c$ \emph{commutes} with the square $(i,j)$ at $t$ if for every volume datum $M$ of level $t$, every matched
pair $(\Theta,\Theta')$ and every cuspidal listed $F$ at $(i,j+1)$ which is convergent along both routes with $\Theta$
and $\Theta'$,
\begin{equation}\label{eq:pf:rowprops:commute}
\rowpropsL_{\mathrm{big}}(\Theta)F=c\cdot\rowpropsL_{\mathrm{small}}(\Theta')F .
\end{equation}
\item The \emph{cusp step} holds for $(W,M,\mathrm{rm},t)$ if for every arrow $a$ of the grid (the arrows
$\rightarrow_{i,j}$ and $\downarrow^t_{i,j}$ in range), every idempotent $e$ which is a sum of local idempotents at
$\mathrm{loc}(t)$, every local vector $f$ at $e$ relative to the vacuum vector of $\mathrm{rm}$ on the pair of $a$, every
polynomial $Q$ and every cuspidal listed $F$ at the source of $a$, $F$ is convergent for $(a,\mathrm{fock}(f,Q))$ and
$\rowpropsL_a(\mathrm{fock}(f,Q))F$ is cuspidal listed at the target.
\end{enumerate}
\formal{GlobalTheory.FlipRun.IsCuspListed}\formal{GlobalTheory.FunObj.CommutesWith}\formal{GlobalTheory.FlipRun.CuspStep}
\end{definition}

\begin{lemma}[Squares commute with their constants]\label{pf:rowprops:commute}
Assume $F$, $\sfname{SquareSt.WeilConvergence}$ and $\sfname{ThetaIntegralsAgree}$. At large scale (with $C_{F0}=0$,
$x_0=0$), for every admissible choice, every $t$ and all $i,j<m$, the constant $c_{ijt}$ commutes with the square $(i,j)$
at $t$.
\formallabel{GlobalTheory.T.twist_square_commutes}
\formallabel{GlobalTheory.SquareComm.square_commutes_with_all}
\formallabel{GlobalTheory.FunObj.commutesWith_of_isSquareConst}
\formallabel{GlobalTheory.FunObj.commutesWith_of_isSquareConst_of_triple}
\formallabel{GlobalTheory.SquareComm.squareStanding_of_setting}
\end{lemma}

\begin{proof}
The Hida setting gives the square standing condition for every $t$ (as in \cref{pf:rowprops:standing}(b)), hence the
square is realised by an admissible triple (\lref{GlobalTheory.FunObj.squareTriple_of_standing}). Since the choice is
admissible, $c=c_{ijt}$ is a square constant at $(i,j,t)$. Let $M$, $(\Theta,\Theta')$ and $F$ be as in
\cref{def:pf:rowprops:commutes}(b), and fix $g_1\in G_{V_{i+1,j}}(\A)$.
\begin{enumerate}[label=(\arabic*)]
\item By \lref{GlobalTheory.T.thetaLift_eq_integral_routeKernel_fubini} (Fubini's theorem, with the convergence of
the theta integrals of the triple given by $\sfname{SquareSt.WeilConvergence}$, the decay of cusp forms given by $F$, and
$\sfname{ThetaIntegralsAgree}$), for almost every $g\in\rowpropsfd_{i,j+1}$ the kernel $\rowpropsKbig_\Theta(g,g_1)$ is
integrable on $\rowpropsfd_{i+1,j+1}$, and
\[
\rowpropsL_{\mathrm{big}}(\Theta)F(g_1)=\int_{\rowpropsfd_{i,j+1}}F(g)\int_{\rowpropsfd_{i+1,j+1}}\rowpropsKbig_\Theta(g,g_1)(u)\,du\,dg .
\]
\item By \lref{GlobalTheory.T.exists_theta_translates_convergent_lift_degree} there is a finite combination $z$ of
elements of $G_{V_{i,j}}(\A)$ with $\rowpropsdeg z\neq0$ such that for almost every $g$ the function
$z\cdot\rowpropsKsmall_{\Theta'}(g,g_1)$ is integrable on $\rowpropsfd_{i,j}$ and
\[
\rowpropsdeg z\cdot\rowpropsL_{\mathrm{small}}(\Theta')F(g_1)=\int_{\rowpropsfd_{i,j+1}}F(g)\int_{\rowpropsfd_{i,j}}
\bigl(z\cdot\rowpropsKsmall_{\Theta'}(g,g_1)\bigr)(u)\,du\,dg .
\]
\end{enumerate}
For almost every $g$ both integrability conditions hold, and \eqref{eq:pf:rowprops:sqconst} gives
$\rowpropsdeg z\int\rowpropsKbig_\Theta(g,g_1)=c\int z\cdot\rowpropsKsmall_{\Theta'}(g,g_1)$. Multiplying by $F(g)$ and
integrating over $g$,
\[
\rowpropsdeg z\cdot\rowpropsL_{\mathrm{big}}(\Theta)F(g_1)=c\cdot\rowpropsdeg z\cdot\rowpropsL_{\mathrm{small}}(\Theta')F(g_1),
\]
and $\rowpropsdeg z\ne0$ can be cancelled.
\end{proof}

\paragraph{Paths through the grid.}
Monotone paths from $(0,m)$ to $(m,0)$ made of arrows $\rightarrow$ and $\downarrow^t$ are recorded by their sequences of
edges. A path \emph{passes through} the square $(i,j)$ if it contains the small route of $(i,j)$; \emph{flipping} it at
$(i,j)$ replaces the small route by the big route. Let $\mathrm{fo}(m)$ be the list of the $m^2$ squares $(i,j)$,
$0\le i,j<m$, in lexicographic order.

\begin{lemma}[Combinatorics of flips]\label{pf:rowprops:flipcomb}
\begin{enumerate}[label=(\alph*)]
\item The given path $\rowpropsgiven_t$ and the opposite path $\rowpropsopp_t$ are paths on the grid, and a path on the
grid is determined by its source, its target and its sequence of edges.
\item If a path on the grid from $X$ to $Y$ passes through the square $(i,j)$ ($i,j<m$), it is the composite of a path
$A$ from $X$ to $(i,j+1)$, the small route of $(i,j)$ and a path $B$ from $(i+1,j)$ to $Y$, with $A$, $B$ on the grid;
flipping replaces it by $A$, the big route, $B$.
\item The squares of $\mathrm{fo}(m)$ can be flipped in this order starting from the given path (each flip is applied to
a path which passes through the square), and the result is the opposite path.
\item For every function $f$ on pairs, the product of $f$ over $\mathrm{fo}(m)$ is $\prod_{i<m}\prod_{j<m}f(i,j)$.
\end{enumerate}
\formallabel{GlobalTheory.T.given_opposite_gridPaths_and_determined_by_edges}
\formallabel{GlobalTheory.T.gridPath_cut_at_square_swap_route}
\formallabel{GlobalTheory.FlipRun.prod_flipOrder}
\formallabel{GlobalTheory.FlipRun.passesThrough_of_canFlip}
\formallabel{GlobalTheory.FlipRun.lt_of_mem_flipOrder}
\end{lemma}

\begin{proof}
(a) and (b) are statements about lattice paths: a path on the grid is a sequence of unit steps to the right or down,
each arrow being determined by its endpoints, so the edge sequence and the source determine the path; cutting a path at
the first occurrence of the small route of $(i,j)$ gives (b). (c) is \lref{GlobalTheory.FunObj.canFlip_flipOrder}:
in lexicographic order, the square $(i,j)$ is flipped when the current path goes down column $i$ from row $j+1$ to row
$j$ and then right along row $j$, and after all flips the path is the opposite path; the list has length $m^2$. (d) is a
reordering of a finite product.
\end{proof}

\begin{lemma}[One flip]\label{pf:rowprops:oneflip}
Let $c\neq0$ commute with the square $(i,j)$ at $t$, let $A$, $B$ be paths on the grid as in
\cref{pf:rowprops:flipcomb}(b), let $\gamma$ be the path $A$, small route, $B$ and $\gamma'$ the path $A$, big route, $B$.
Let $Q$ be a finite combination of polynomial families, $\tau$ a finite combination of finite tensors along $\gamma$
representing a family of local tensors $(T_w)_{w\in\mathrm{loc}(t)}$, and suppose that the local tensors $T_w$ and their
local flips at $(i,j)$ are tensors of the local grids along $\gamma$, resp.\ $\gamma'$. Then there is a finite combination
$\tau'$ of finite tensors along $\gamma'$ representing the flipped local tensors such that for every volume datum $M$ of
level $t$ for which the cusp step holds and every cuspidal listed $F$ at $(0,m)$ convergent along $\gamma$ with the
tensors of $(Q,\tau)$: $F$ is convergent along $\gamma'$ with those of $(\mathrm{flip}_{ij}Q,\tau')$, and
\[
\rowpropsL_\gamma(Q,\tau)F=c^{-1}\,\rowpropsL_{\gamma'}(\mathrm{flip}_{ij}Q,\tau')F .
\]
\formallabel{GlobalTheory.T.gridSquareFlip_smallRoute_eq_inv_smul_bigRoute}
\formallabel{GlobalTheory.FlipRun.oneFlip}
\end{lemma}

\begin{proof}
Fix an idempotent $e$, a sum of local idempotents at the places of $\mathrm{loc}(t)$, for which $\tau$ represents
$(T_w)_w$.

\emph{Decomposition.} By \lref{GlobalTheory.FlipRun.exists_decomp} there are a finite index set $s$ and, for
$\kappa\in s$, pure tensors $\alpha_\kappa$ along $A$ and $\beta_\kappa$ along $B$ made of local vectors at $e$, and finite
combinations $\sigma^{\mathrm{sm}}_\kappa$, $\sigma^{\mathrm{big}}_\kappa$ of tensors along the small and the big route,
such that: for every $\kappa$, the pair $(\sigma^{\mathrm{big}}_\kappa,\sigma^{\mathrm{sm}}_\kappa)$ has the local flip
property at $(i,j)$ for a natural datum $\Phi_\kappa$ of the local depth at every $w\in\mathrm{loc}(t)$; the combination
$\tau^{\mathrm{sm}}=\sum_\kappa\alpha_\kappa\,\sigma^{\mathrm{sm}}_\kappa\,\beta_\kappa$ along $\gamma$ (concatenation of
tensors) represents $(T_w)_w$; and $\tau'=\sum_\kappa\alpha_\kappa\,\sigma^{\mathrm{big}}_\kappa\,\beta_\kappa$ along $\gamma'$
represents the local flips of the $T_w$ at $(i,j)$. This uses the hypotheses that the $T_w$ and their flips are tensors
along $\gamma$ and $\gamma'$.

\emph{Equal product functions.} Since $\tau$ and $\tau^{\mathrm{sm}}$ represent the same local tensors, their product
functions agree at all half-integral arguments; as the cusp step holds, two such finite tensors give the same lift of every
cuspidal listed function, for every polynomial datum $Q$ (\lref{GlobalTheory.T.totalTensorThetaLift_eq_of_schwartz_eq}).
Hence $\rowpropsL_\gamma(Q,\tau)F=\sum_\kappa\rowpropsL_\gamma(Q,\alpha_\kappa\sigma^{\mathrm{sm}}_\kappa\beta_\kappa)F$, the lift
being linear in the tensor.

\emph{One sandwiched square.} For each $\kappa$,
$\rowpropsL_\gamma(Q,\alpha_\kappa\sigma^{\mathrm{sm}}_\kappa\beta_\kappa)F=c^{-1}\rowpropsL_{\gamma'}(\mathrm{flip}_{ij}Q,
\alpha_\kappa\sigma^{\mathrm{big}}_\kappa\beta_\kappa)F$ (\lref{GlobalTheory.T.sandwichedFlip_routeThroughSmall_eq_inv_const_routeBig}):
the lift of $F$ along $A$ with $\alpha_\kappa$ is again cuspidal listed by the cusp step; the tensors of the two routes built
from $\sigma^{\mathrm{big}}_\kappa$, $\sigma^{\mathrm{sm}}_\kappa$ and the polynomial parts of $\mathrm{flip}_{ij}Q$ and
$Q$ form a matched pair, by the local flip property and the definition of the flip of polynomial data through the
relation of the real match; so \eqref{eq:pf:rowprops:commute} applies to the middle square; finally one lifts along $B$
with $\beta_\kappa$.

Summing over $\kappa$ gives the identity. The convergence along $\gamma'$ with the tensors of
$(\mathrm{flip}_{ij}Q,\tau')$ follows from the cusp step, since all components are local vectors at $e$ and $F$ is cuspidal
listed (\lref{GlobalTheory.FlipRun.totConv_tensE}).
\end{proof}

\begin{lemma}[Representing the given path]\label{pf:rowprops:swapped}
Assume $F$. At large scale, for every admissible choice, every $t$ and every family $f=(f_j)_{j<m}$ of \emph{chain
vectors} at $t$ (the Weil vectors of $R$ on the arrows of the first column: $f_j$ is the vector of $\omega_Y(t,j)$ if
$j+1<m$ and of $\omega_X(t)$ if $j+1=m$), there is $z\in\rowpropsXi_t$ such that for every coset datum $\Phi$ the finite
tensor along the given path with column components $f_j$, row components the coset vectors $\mathrm{cv}_i(\Phi_i)$, and
coefficient $z$, represents the local tensors of the given path attached to $\Phi$.
\formallabel{GlobalTheory.T.twisted_theta_first_occurrence_finite_swapped_pair}
\formallabel{GlobalTheory.FlipStart.theta_first_occurrence_finite_swapped_pair}
\formallabel{GlobalTheory.FlipStart.repFin_of}
\end{lemma}

\begin{proof}
(a) For every column arrow $j<m$ there is $c_j\in\rowpropsXi_t$ such that for every idempotent $e$ which is a sum of local
idempotents at $\mathrm{loc}(t)$ and every $\Phi$, the vector $f_j$ is a local vector at $e$, relative to the vacuum vector of
the real match on $\mu(t,0,j)$, of $c_j$ times the local function of the given path on the arrow $\downarrow^t_{0,j}$: for the
inner arrows ($j+1<m$) this is \lref{GlobalTheory.T.twist_chainColumn_inner_succ_eq}, for the top arrow ($j+1=m$) it is
\lref{GlobalTheory.T.twist_flipStart_highestVector}.

(b) The coset vectors $\mathrm{cv}_i(\Phi_i)$ are local vectors at $e$ of the local functions of the given path on the
row arrows $\rightarrow_{i,0}$ (\lref{GlobalTheory.T.twist_coset_function_isLocalVector}).

(c) Choose such an idempotent $e$ (it exists, \lref{GlobalTheory.FunObj.LocTensor.exists_idemSum_locPl}) and put
$z=(\prod_jc_j)^{-1}$, an element of the group $\rowpropsXi_t$. The product function of the single tensor with coefficient
$z$, computed with the local functions of (a) and (b), is $z\prod_jc_j=1$ times the product of the local functions of
the given path, which is the product over $w\in\mathrm{loc}(t)$ of the local tensors of the given path. This is the
representing property.
\end{proof}

\begin{lemma}[The initial term]\label{pf:rowprops:initial}
Assume $F$ and $\sfname{LiftsCuspidal}$. At large scale, for every admissible choice there is a finite combination $P_g$
of polynomial families such that for every $t$ there is $z_t\in\rowpropsXi_t$ with the following property: for every
coset datum $\Phi$ there is a finite tensor $\tau_g$ along the given path representing the local tensors of the given
path attached to $\Phi$, such that for every $M$ of level $t$ and every $y$, $F_{t,y}$ is convergent along
$\rowpropsgiven_t$ with the tensors of $(P_g,\tau_g)$ and
\[
\rowpropsth\bigl(\varrho,\Phi;\mathrm{lift}_t(y)\bigr)=z_t^{-1}\,\rowpropsL_{\rowpropsgiven_t}(P_g,\tau_g)F_{t,y}.
\]
\formallabel{GlobalTheory.T.flipSequence_initial_term}
\formallabel{GlobalTheory.FlipStart.flipSequence_start_condition}
\end{lemma}

\begin{proof}
Let $q=(q_j)$ be the polynomial data of the chain vector of the choice ($q_j=P_Y(j)$ for $j+1<m$ and $q_j=P_X$ for $j+1=m$),
and let $P_g$ be the polynomial datum which is $q$ on the first column and the real vector $\varrho$ on the first row. For
$t\in T$ let $f$ be the chain vectors at $t$ (\cref{pf:rowprops:swapped}), and $z_t$ the element given by
\cref{pf:rowprops:swapped}. Given $\Phi$, let $\tau_g$ be the single finite tensor of \cref{pf:rowprops:swapped} with
coefficient $z_t$; it represents the local tensors of the given path. With the polynomial datum $P_g$, its Fock tensors
are, on the first column, $\mathrm{fock}(f_j,q_j)$, and on the first row, the row datum $\mathrm{fock}(\mathrm{cv}_i(\Phi_i),\varrho_i)$
of \cref{def:pf:rowprops:bessel}. Let $M$ be of level $t$ and $y\in H_X(t)$.
\begin{enumerate}[label=(\alph*)]
\item By \lref{GlobalTheory.T.flipSequence_column_succ} (which uses $F$ and $\sfname{LiftsCuspidal}$), $F_{t,y}$ is convergent
along the down path $\delta^t_0$ with the tensors $\mathrm{fock}(f_j,q_j)$, and the lift is $\mathrm{lift}_t(y)$.
\item By \lref{GlobalTheory.T.twisted_theta_tower_first_occurrence_exists}, $\mathrm{lift}_t(y)$ is convergent along the
row path $\rho_0$ with the row datum, and the lift computed with $M$ equals the row lift
$\rowpropsth(\varrho,\Phi;\mathrm{lift}_t(y))$ computed with the volume datum of the choice.
\end{enumerate}
Since $\rowpropsgiven_t$ is $\delta^t_0$ followed by $\rho_0$ and lifts compose along composite paths, (a) and (b) give the
convergence along $\rowpropsgiven_t$ and $\rowpropsL_{\rowpropsgiven_t}(P_g,\tau_g)F_{t,y}=z_t\,\rowpropsth(\varrho,\Phi;\mathrm{lift}_t(y))$.
Divide by $z_t\ne0$.
\end{proof}

\begin{lemma}[Local flip stages are tensors]\label{pf:rowprops:localtensor}
At large scale (with $C_{F0}=0$, $x_0=0$), for every admissible choice, every $t$, every coset datum $\Phi$, every
$w\in\mathrm{loc}(t)$ and every $k\le m^2$, the local tensor obtained from the local tensor of the given path at $w$ by the
first $k$ flips of $\mathrm{fo}(m)$ is a tensor of the local grid at $w$ along the path obtained by the same flips.
\formallabel{GlobalTheory.FlipRun.isTensor_flips_take}
\formallabel{GlobalTheory.T.localFlipStage_isTensor_placeDepth}
\formallabel{GlobalTheory.FlipRun.canFlip_of_append}
\end{lemma}

\begin{proof}
The primes $p_{de}$, $p'_{de}$ of the level split in the auxiliary field of the host (a size condition of the Hida
setting), hence are coprime to $Q_S\,q$. With the guard condition this makes the local grid at $w$ good
(\lref{GlobalTheory.FunObj.Good.locGrid_isGood}): in particular every local square $(i,j)$ has a uniformiser, an
additive character of conductor one, a depth at least one, and the character $\nu_{ij}$ satisfies the conductor bound of
the local flip. The local stage of the given path is admissible (\lref{GlobalTheory.FunObj.GivenAdm.givenStage_isAdmissible}),
so its tensor is a tensor along the given path, and the local levels are open
(\lref{GlobalTheory.FunObj.GivenAdm.givenLevel_open}). The first $k$ squares of $\mathrm{fo}(m)$ can be flipped starting
from the given path, since a prefix of a list of flips that can be performed can be performed
(\cref{pf:rowprops:flipcomb}(c)). Under these conditions, performing the local flips of a list of squares which can be
flipped carries a tensor along the initial path to a tensor along the flipped path
(\lref{GlobalTheory.FunObj.isTensor_flips}); this gives the claim.
\end{proof}

\begin{proposition}[Flip]\label{pf:rowprops:flipmain}
Assume the bundle $F$, the statement $\sfname{LiftsCuspidal}$, the square convergence statement, the agreement of theta
integrals and $\sfname{ConstK}$. Then $\sfname{StepFlip}$ holds.
\formallabel{GlobalTheory.T.twisted_lift_flip_invariant}
\formallabel{GlobalTheory.FlipRun.twist_sign_flip_succ}
\formallabel{GlobalTheory.FlipRun.flipAt_of}
\end{proposition}

\begin{proof}
At large scale, for an admissible choice the following hold for all $t$: (A) each $c_{ijt}$ commutes with the square
$(i,j)$ at $t$ (\cref{pf:rowprops:commute}); (B) the cusp step holds for every $M$ of level $t$
(\cref{pf:rowprops:cuspstep}); (C) every $F_{t,y}$ is cuspidal listed at $(0,m)$
(\lref{GlobalTheory.EndCusp.leadingForm_mem_cuspList}); (D) the conclusion of \cref{pf:rowprops:initial}; (E) for every
coset datum $\Phi$, every $w\in\mathrm{loc}(t)$ and every $k\le m^2$, the local tensor obtained from the local tensor of
the given path by the first $k$ flips of $\mathrm{fo}(m)$ is a tensor of the local grid along the current path
(\cref{pf:rowprops:localtensor}); (F) $c_{ijt}\ne0$ ($\sfname{ConstK}$).

Fix $t$ and $\Phi$. Let $P_g$, $z_t$, $\tau_g$ be as in (D), and put $c'(i,j)=c_{ijt}$. For $k\le m^2$ let $S_k$ be the family of
local tensors obtained from those of the given path by the first $k$ flips of $\mathrm{fo}(m)$, and $Q_k$ the polynomial
data obtained from $P_g$ by the corresponding flips of polynomials. We show by induction on $k$: \emph{there are a path $p_k$
on the grid, whose edge sequence is the result of the first $k$ flips applied to the given path, and a finite tensor
$\tau_k$ along $p_k$ representing $S_k$, such that for every $M$ of level $t$ and every $y$, $F_{t,y}$ is convergent along
$p_k$ with the tensors of $(Q_k,\tau_k)$ and}
\begin{equation}\label{eq:pf:rowprops:flipind}
\rowpropsth\bigl(\varrho,\Phi;\mathrm{lift}_t(y)\bigr)=\Bigl(z_t\prod_{(i,j)\in\mathrm{fo}(m)_{<k}}c'(i,j)\Bigr)^{-1}
\rowpropsL_{p_k}(Q_k,\tau_k)F_{t,y},
\end{equation}
where $\mathrm{fo}(m)_{<k}$ denotes the first $k$ squares. For $k=0$ take the given path and $\tau_g$ (\cref{pf:rowprops:flipcomb}(a)
and (D)). For the step $k\to k+1$, let $(i,j)$ be the $(k+1)$-st square. By \cref{pf:rowprops:flipcomb}(c), $p_k$ passes
through $(i,j)$, so by \cref{pf:rowprops:flipcomb}(b) $p_k=A\cdot(\text{small route})\cdot B$; let $p_{k+1}=A\cdot(\text{big
route})\cdot B$. By (E) for $k$ and $k+1$ the local tensors are tensors along $p_k$ and $p_{k+1}$, and $S_{k+1}$ is the local
flip of $S_k$ at $(i,j)$. \Cref{pf:rowprops:oneflip}, applied with $c=c_{ijt}$ (non-zero by (F), commuting by (A)), gives
$\tau_{k+1}$ representing $S_{k+1}$, the convergence along $p_{k+1}$ (using (B), (C) and the convergence along $p_k$), and
$\rowpropsL_{p_k}(Q_k,\tau_k)F_{t,y}=c_{ijt}^{-1}\rowpropsL_{p_{k+1}}(Q_{k+1},\tau_{k+1})F_{t,y}$; substituting in
\eqref{eq:pf:rowprops:flipind} gives the claim for $k+1$.

For $k=m^2$ the edge sequence of $p_{m^2}$ is that of the opposite path (\cref{pf:rowprops:flipcomb}(c)), hence
$p_{m^2}=\rowpropsopp_t$ by \cref{pf:rowprops:flipcomb}(a); the local tensors $S_{m^2}$ are the flipped tensors
$T^{\mathrm{flip}}_w(\Phi)$; and the product of the $c'(i,j)$ over $\mathrm{fo}(m)$ is $\rowpropsctot(t)$
(\cref{pf:rowprops:flipcomb}(d)). Finally, a finite tensor along the opposite path representing $(T^{\mathrm{flip}}_w(\Phi))_w$,
with the polynomial data $Q_{m^2}$, has the form $\mathrm{opp}(\tau_i,\tau)$ for a flipped finite tensor $\tau$ for $(t,\Phi)$,
where $\tau_i$ is the image of $Q_{m^2}$ (\lref{GlobalTheory.T.flippedFiniteTensor_of_representing_oppositePath}).
Since $Q_{m^2}$ and hence $\tau_i$ do not depend on $t$ and $\Phi$, \eqref{eq:pf:rowprops:flipind} for $k=m^2$ is
\eqref{eq:pf:rowprops:flipat}, and $\mathsf{FlipAt}$ holds with $\tau_i$ and $z$.
\end{proof}

% ---------------------------------------------------------------------------
\subsection{Unfolding}\label{sec:pf:rowprops:unfold}

\subsubsection{Iterated quotient integrals and cells}
\begin{definition}[Iterated quotient integrals]\label{def:pf:rowprops:niter}
Let $(H_i)_{i\in\iota}$ be a family of locally compact groups, $R_i\subseteq H_i$ subgroups with measures $\rho_i$, and
$\mu_i$ measures on $H_i$. For a function $q$ on $\prod_iH_i$, a list $l$ of indices, a point $h$ and $z\in\C$, the relation
$\rowpropsNIter(q,l,h,z)$ (``$z$ is the iterated quotient integral of $q$ over the coordinates in $l$, the others being
fixed at $h$'') is defined by recursion on $l$: for $l$ empty it means $z=q(h)$; for $l=i::l'$ it means that there is a
function $f$ on $H_i$ such that $z$ is the integral of $f$ over $R_i\backslash H_i$ (with respect to $\rho_i$ and $\mu_i$,
$f$ being left $R_i$-invariant and integrable on the quotient) and $\rowpropsNIter(q,l',h[i\leftarrow g],f(g))$ for every
$g\in H_i$.
\formal{GlobalTheory.RowVal.NIter}
\end{definition}

\begin{lemma}[Linearity]\label{pf:rowprops:niterlin}
If $\rowpropsNIter(q_n,l,h,z_n)$ for finitely many $n$ and $\alpha_n\in\C$, then
\[
\rowpropsNIter\Bigl(\sum_n\alpha_nq_n,\,l,\,h,\,\sum_n\alpha_nz_n\Bigr).
\]
If $\rowpropsNIter(q,l,h,z)$ and $q'$ agrees with $q$ at all points which agree with $h$ outside $l$, then
$\rowpropsNIter(q',l,h,z)$.
\formallabel{GlobalTheory.T.isIteratedInvariantIntegral_linear_and_congr}
\formallabel{GlobalTheory.RowVal.niter_lin_congr}
\end{lemma}

\begin{proof}
Induction on $l$. For $l$ empty both statements are immediate. For $l=i::l'$, let $f_n$ be the functions attached to
$q_n$; then $f=\sum_n\alpha_nf_n$ is left $R_i$-invariant and integrable on $R_i\backslash H_i$ with integral
$\sum_n\alpha_nz_n$, and the induction hypothesis applied at the points $h[i\leftarrow g]$ gives
$\rowpropsNIter(\sum_n\alpha_nq_n,l',h[i\leftarrow g],f(g))$. For the second statement use the same function $f$: the
points agreeing with $h[i\leftarrow g]$ outside $l'$ agree with $h$ outside $l$.
\end{proof}

\begin{definition}[Split quadruples; cells; cell weights]\label{def:pf:rowprops:cells}
A quadruple $(A,B,R,K)$ of subgroups of a locally compact group $H$ is \emph{split} if $A$, $B$, $R$ are closed, $A$ and
$B$ commute and multiplication $A\times B\to H$ is a homeomorphism, $ab\in R$ with $a\in A$, $b\in B$ implies $a,b\in R$,
$K\subseteq B$ is compact and open in $B$, Haar measures on $R$ and on $R\cap B$ are inversion invariant and Haar measures on
$B$ are right invariant. For families of such quadruples indexed by $\iota$ and a list $l$ of indices, a function $q$ on
$\prod H_i$ is \emph{bi-invariant} if in every coordinate $i\in l$ it is left $R_i$-invariant and right $K_i$-invariant,
and $h$ lies \emph{in the cell} of $b$ if $h_i\in R_iA_ib_iK_i$ for all $i\in l$. For $b\in B$ the \emph{cell weight} is
$\rowpropscw(b)=\mu'(K)/\rho'(\{r\in R\cap B: b^{-1}rb\in K\})$, for Haar measures $\rho'$ on $R\cap B$ and $\mu'$ on $B$.
\formal{GlobalTheory.RowVal.Split}\formal{GlobalTheory.RowVal.BiInv}\formal{GlobalTheory.RowVal.InCells}\formal{GlobalTheory.RowVal.cellWt}
\end{definition}

\begin{lemma}[Cell sums]\label{pf:rowprops:cellsum}
Let $(B_i,R_i,K_i)_{i\in l}$ be locally compact groups with closed subgroups $R_i$ and compact open subgroups $K_i$, with
Haar measures $\rho_i$ on $R_i$ (inversion invariant) and $\mu_i$ on $B_i$ (right invariant). Let $S$ be a finite set,
$b:S\to\prod_iB_i$, distinct elements of $S$ having disjoint cells $R\,b(\beta)K$, and let $q$ be a bi-invariant function on
$\prod_iB_i$ vanishing outside these cells. Then the iterated quotient integral of $q$ over $l$ is
$\sum_{\beta\in S}\bigl(\prod_{i\in l}\rowpropscw(b(\beta)_i)\bigr)\,q(b(\beta))$.
\end{lemma}

This is \lref{GlobalTheory.T.iteratedInvariantIntegral_doubleCoset_cellSum}.

\begin{lemma}[Vanishing from the finite parts]\label{pf:rowprops:innervanish}
Let $(A_i,B_i,R_i,K_i)_{i\in l}$ be split quadruples in $H_i$ ($l$ without repetition) with Haar measures, let $C$ be a finite
set of points of $\prod_iB_i$, and let $q$ be a bi-invariant function on $\prod_iH_i$ which, at points trivial outside $l$,
vanishes outside the cells of the points of $C$, and has an iterated quotient integral $z$ over $l$ (at the point $1$).
If for every $a\in\prod_iA_i$ trivial outside $l$ the iterated quotient integral over $l$ of $b\mapsto q(ab)$ on
$\prod_iB_i$ vanishes, then $z=0$.
\end{lemma}

This is \lref{GlobalTheory.T.isIteratedInvariantIntegral_vanish_of_inner_vanish}.

\begin{proposition}[Relations between cell sums pass to integrals]\label{pf:rowprops:cellvan}
Let $(A_i,B_i,R_i,K_i)_{i\in l}$ be split quadruples with Haar measures $\rho_i,\mu_i$ on $R_i,H_i$ and
$\rho'_i,\mu'_i$ on $R_i\cap B_i,B_i$. For $n\in\{0,1\}$ let $S_n$ be finite sets, $b_n:S_n\to\prod_iB_i$ (trivial outside $l$)
with weights $w_n(\beta)=\prod_{i\in l}\rowpropscw(b_n(\beta)_i)$, distinct elements of $S_n$ having distinct cells; let
$q_n$ be bi-invariant functions which, at points trivial outside $l$, vanish outside the cells of $b_n(S_n)$; let $z_n$ be
iterated quotient integrals of $q_n$ over $l$ at $1$; and let $\alpha_0,\alpha_1\in\C$. If
\[
\sum_n\alpha_n\sum_{\beta\in S_n}w_n(\beta)\,q_n\bigl(a\,b_n(\beta)\bigr)=0
\]
for every $a\in\prod_iA_i$ trivial outside $l$, then $\alpha_0z_0+\alpha_1z_1=0$.
\formallabel{GlobalTheory.T.cellSumRelationsPassToIntegrals_proof}
\formallabel{GlobalTheory.RowVal.cellVan_holds}
\end{proposition}

\begin{proof}
By linearity, $\alpha_0z_0+\alpha_1z_1$ is an iterated quotient integral of $q=\alpha_0q_0+\alpha_1q_1$. The function $q$ is
bi-invariant and, at points trivial outside $l$, vanishes outside the cells of the finite set
$C=b_0(S_0)\cup b_1(S_1)\subseteq\prod_iB_i$. By \cref{pf:rowprops:innervanish} it suffices to show that for every
$a\in\prod_iA_i$ trivial outside $l$, the iterated quotient integral over $l$ of $b\mapsto q_n(ab)$ on $\prod_i B_i$ (with
$\rho'$, $\mu'$) equals $\sum_{\beta\in S_n}w_n(\beta)q_n(a\,b_n(\beta))$; then the inner integral of $q$ is
$\sum_n\alpha_n\sum_\beta w_n(\beta)q_n(ab_n(\beta))=0$ by hypothesis. The function $b\mapsto q_n(ab)$ on $\prod_iB_i$ is
bi-invariant for $(R_i\cap B_i,K_i)$, because $A_i$ and $B_i$ commute and $R_i$ splits; it vanishes outside the cells
$(R_i\cap B_i)b_n(\beta)_iK_i$, because $ab\in R_iA_ib_n(\beta)_iK_i$ with $a\in A_i$ forces $b\in(R_i\cap
B_i)b_n(\beta)_iK_i$ by the splitting; and distinct $\beta$ have disjoint cells in $\prod_iB_i$ (a coincidence there would be
a coincidence of the cells in $\prod_iH_i$). \Cref{pf:rowprops:cellsum} gives the claimed value.
\end{proof}

\begin{lemma}[Proportionality]\label{pf:rowprops:prop}
Let $L,r:I\to\C$ be functions on a set $I$ such that for all $x,x'\in I$ and $\alpha,\alpha'\in\C$ with
$\alpha r(x)+\alpha'r(x')=0$ one has $\alpha L(x)+\alpha'L(x')=0$. Then there is $c\in\C$ with $L=c\,r$.
\formallabel{GlobalTheory.RowRealInt.exists_mul_of_rel}
\end{lemma}

\begin{proof}
If $r\equiv0$, take $x'=x$, $\alpha=1$, $\alpha'=0$: then $L(x)=0$ for all $x$, and $c=0$ works. Otherwise fix $x_0$ with
$r(x_0)\ne0$ and put $c=L(x_0)/r(x_0)$. For $x\in I$ apply the hypothesis to $(x,x_0)$ with $\alpha=r(x_0)$ and
$\alpha'=-r(x)$: then $r(x_0)L(x)-r(x)L(x_0)=0$, i.e.\ $L(x)=c\,r(x)$.
\end{proof}

\begin{corollary}[Iterated integrals are proportional to cell sums]\label{pf:rowprops:propcell}
In the situation of \cref{pf:rowprops:cellvan}, let $(S_x,b_x,w_x,q_x,L_x)_{x\in I}$ be a family of data as there (one
for each $x$), and $r:I\to\C$. Suppose that for every $a\in\prod_iA_i$ trivial outside $l$ there is $e(a)\in\C$ with
$\sum_{\beta\in S_x}w_x(\beta)q_x(ab_x(\beta))=e(a)\,r(x)$ for all $x$. Then there is $c$ with $L_x=c\,r(x)$ for all $x$.
\formallabel{GlobalTheory.RowRealInt.exists_mul_of_cellVan}
\formallabel{GlobalTheory.RowRealInt.exists_mul_of_cellVan'}
\end{corollary}

\begin{proof}
By \cref{pf:rowprops:prop} it suffices to show: if $\alpha r(x)+\alpha'r(x')=0$ then $\alpha L_x+\alpha'L_{x'}=0$. Apply
\cref{pf:rowprops:cellvan} to the two data $x,x'$ with coefficients $\alpha,\alpha'$: for every $a$, the relation of cell
sums is $\alpha e(a)r(x)+\alpha'e(a)r(x')=0$.
\end{proof}

\subsubsection{The last-row integrand}
Fix an admissible choice. Let $I=\{1,\dots,m-2\}$. For $j\in I$ let $H_j=G_{V_{m,j}}(\A)$, $A_j=G_{V_{m,j},\infty}$,
$B_j=G_{V_{m,j},\mathrm f}$, let $R_j\subseteq H_j$ be the Jacobi group (if $m-j$ is odd) or the Bessel group (if $m-j$ is even)
attached to the decomposition of $V_{m,j}$ given by the grid, and $K_j=B_j\cap(\text{the maximal compact level of
}\Lambda_{m,j})$. A point is a family $h=(h_j)_{j\in\N}$, $h_j\in H_j$; it is \emph{real} if $h_j\in A_j$ for all $j$ and
$h_j=1$ for $j\notin I$.

\begin{definition}[The last-row integrand]\label{def:pf:rowprops:tot}
For a volume datum $M$ of level $t$, a pair $P=(P_1,P_2)$ of polynomial families on the last row and the last column, a
finite combination $\tau$ of pairs $d=(d_1,d_2)$ as in \cref{def:pf:rowprops:flipfin}, functions $\phi$ on
$G_{V_{\mathrm t_0}}(\A)$ and $F$ on $G_{V_{0,m}}(\A)$, and a point $h$, the \emph{last-row integrand} is
\[
\rowpropstot(P,\tau;\phi,F)(h)=\sum_d\tau(d)\,\rowpropsper\bigl(\mathrm{con}(\mathrm{fock}(d_2,P_2);h),\,d_1,\,P_1;\,\phi,F\bigr).
\]
Here $\mathrm{con}(f;h)$ is the family of functions $C_j$ of two consecutive top-column arguments obtained from the theta
kernels of the down arrows $\downarrow^t_{m,j}$ of the last column with the vectors $f_j$, evaluated at the points
$h_j,h_{j+1}$ and contracted along the Bessel and Jacobi decompositions of the spaces $V_{m,j}$; and for such a family
$C$ and $f=d_1$, $\rowpropsper(C,f,P_1;\phi,F)$ is the integral over $\rowpropsfd_{m,m}$ of the row lift
$\rowpropsL_{\rho_m}(\mathrm{fock}(f_i,P_{1,i}))F$ against the function obtained from $\phi$ by $m/2$ two-step lifts up the
top column with the kernels $C_0,C_2,\dots$, composed with the embedding $\iota_{\mathrm t}$. The \emph{product function}
of $(C,f)$ is $\mathrm{ev}(C,f)(u,y)=\prod_{k<m/2}C_{2k}(u_{2k},u_{2k+1})\prod_if_i(y_i)$, a function of the top-column
arguments $u$ and the last-row arguments $y$. A pair $(C,f)$ is \emph{good} if the $C_{2k}$ are bi-theta functions and the
$f_i$ theta functions for the Weil representations.
\formal{GlobalTheory.RowVal.tot}\formal{GlobalTheory.RowVal.per}\formal{GlobalTheory.RowVal.conAt}\formal{GlobalTheory.RowVal.evT}\formal{GlobalTheory.RowVal.Good}
\end{definition}

For a flipped finite tensor $\tau$ for $(t,\Phi_a)$ write
$\mathrm{lhs}(\tau_i,\tau;a,y)=\rowpropsB\bigl(\phi_a,\rowpropsL_{\rowpropsopp_t}(\mathrm{opp}(\tau_i,\tau))F_{t,y}\bigr)$ for
the left side of \eqref{eq:pf:rowprops:unfoldat}.

\begin{definition}[The unfolding properties]\label{def:pf:rowprops:unfprops}
For an admissible choice $\mathrm{ch}$ and a sign $s\in\{\pm1\}$ consider the following properties; ``for all data'' means
for all $t$, $a$, flipped finite tensors $\tau$ for $(t,\Phi_a)$, volume data $M$ of level $t$, polynomial pairs $P$ and
$y\in H_X(t)$.
\begin{enumerate}[label=(\alph*)]
\item $\mathsf{GrpAt}$: for $j\in I$ the quadruple $(A_j,B_j,R_j,K_j)$ is split.
\item $\mathsf{IterAt}$: there are Haar measures $\rho_j$ on $R_j$ and $\mu_j$ on $H_j$ ($j\in I$) such that for all data,
$\mathrm{lhs}(\delta_P,\tau;a,y)$ is an iterated quotient integral over $I$ of $\rowpropstot(P,\tau;\phi_a,F_{t,y})$ at $1$
($\delta_P$ the combination with value $1$ at $P$).
\item $\mathsf{LinAt}$: $\mathrm{lhs}(\tau_i,\tau;a,y)=\sum_P\tau_i(P)\,\mathrm{lhs}(\delta_P,\tau;a,y)$.
\item $\mathsf{InvAt}$: for all data, $\rowpropstot(P,\tau;\phi_a,F_{t,y})$ is bi-invariant for $(R_j,K_j)_{j\in I}$.
\item $\mathsf{RepAt}(s)$: for every $t$, every family $\beta$ of local cosets has a \emph{representative point} $b(\beta)$
(trivial outside $I$, with $b(\beta)_j\in B_j$ given by a rational ``tilt'' times an element whose local components at
$\mathrm{loc}(t)$ are the staircase representatives of $\beta$, twisted by $s$), and distinct $\beta$ have disjoint cells.
\item $\mathsf{SuppOn}(s)$: for all data and all choices of representative points, if $h$ is trivial outside $I$ and
$\rowpropstot(P,\tau;\phi_a,F_{t,y})(h)\ne0$, then $h$ lies in the cell of $b(\beta)$ for some
$\beta\in\rowpropsSupp(t,\Phi_a)$.
\item $\mathsf{ValAt}(s)$: there is $z:T\to\C^\times$ with $z_t\in\rowpropsXi_t$ such that for all Haar measures
$\rho'_j,\mu'_j$ on $R_j\cap B_j$, $B_j$, all $P$ and all real points $p$, there is $\sigma$ such that for all data, all
$\beta\in\rowpropsSupp(t,\Phi_a)$ with representative point $b$, all term volumes $M'$ and expansions $\tau'$ of $\beta$:
\begin{equation}\label{eq:pf:rowprops:valat}
w(b)\,\rowpropstot(P,\tau;\phi_a,F_{t,y})(p\,b)=z_t^{-1}\Bigl(\prod_w\mathrm{coef}_w(\beta_w)\Bigr)\rowpropsdef(M';\sigma,\tau';\phi_a,F_{t,y}),
\end{equation}
where $w(b)=\prod_{j\in I}\rowpropscw(b_j)$.
\item $\mathsf{RealLine}$: there is a non-zero linear form $\ell$ on polynomial pairs $P$ (column and last row) such that:
for all $t,a$ with $\mathsf{SumOK}(t,\Phi_a)$, all $\beta\in\rowpropsSupp(t,\Phi_a)$, term volumes $M$, expansions $\tau$ and
$y$, there is $K$ with $\rowpropsdef(M;\sigma,\tau;\phi_a,F_{t,y})=\ell(\sigma)\,K$ for all $\sigma$; and every $z\in\C$ is
$\ell(\sigma)$ for some $\sigma$ (with $\sigma=0$ if $z=0$). Here $\ell(\sigma)=\sum_P\sigma(P)\ell(P)$.
\end{enumerate}
\formal{GlobalTheory.RowVal.GrpAt}\formal{GlobalTheory.RowVal.IterAt}\formal{GlobalTheory.RowVal.LinAt}\formal{GlobalTheory.RowVal.InvAt}\formal{GlobalTheory.RowVal.RepAt}\formal{GlobalTheory.RowVal.SuppOn}\formal{GlobalTheory.RowVal.ValAt}\formal{GlobalTheory.RowReal.RealLine}
\end{definition}

The property $\mathsf{RealLine}$ says that all definite periods, as functions of the archimedean polynomial data, lie on
one line: the archimedean part of the unfolding is the same, up to a scalar, for all cells, twists and test vectors.

\begin{proposition}[The unfolding identity]\label{pf:rowprops:unfoldAt}
Let $\mathrm{ch}$ be an admissible choice with $\mathsf{GrpAt}$, $\mathsf{IterAt}$, $\mathsf{LinAt}$, $\mathsf{InvAt}$,
$\mathsf{RealLine}$, $\mathsf{SumOK}(t,\Phi_a)$ for all $t,a$, and $\mathsf{SuppOn}(s)$, $\mathsf{RepAt}(s)$,
$\mathsf{ValAt}(s)$ for some $s$. Then $\mathsf{UnfoldAt}$ holds.
\formallabel{GlobalTheory.RowRealInt.unfoldAt_of}
\formallabel{GlobalTheory.RowRealInt.exists_lhs_eq}
\formallabel{GlobalTheory.RowRealInt.exists_haar_family}
\formallabel{GlobalTheory.FunObj.sAA_facts}
\end{proposition}

\begin{proof}
Let $z$ be given by $\mathsf{ValAt}(s)$ and $\ell$ by $\mathsf{RealLine}$; fix $\sigma_0$ with $\ell(\sigma_0)=1$. By
$\mathsf{RealLine}$, for $\beta$ in the support, term volumes $M'$, expansions $\tau'$ and all $\sigma$,
\begin{equation}\label{eq:pf:rowprops:hdef}
\rowpropsdef(M';\sigma,\tau';\phi_a,F_{t,y})=\ell(\sigma)\,\rowpropsdef(M';\sigma_0,\tau';\phi_a,F_{t,y}).
\end{equation}
Put $g(t,a,\beta,M',\tau',y)=(\prod_w\mathrm{coef}_w(\beta_w))\rowpropsdef(M';\sigma_0,\tau';\phi_a,F_{t,y})$.

\emph{Step 1: for every polynomial pair $P$ there is $c_P\in\C$ such that for all data and all families $(M_\beta,\tau_\beta)$
of term volumes and expansions,}
\[
\mathrm{lhs}(\delta_P,\tau;a,y)=c_P\,z_t^{-1}\sum_{\beta\in\rowpropsSupp(t,\Phi_a)}g(t,a,\beta,M_\beta,\tau_\beta,y).
\]
Let $\rho_j,\mu_j$ be the measures of $\mathsf{IterAt}$, and choose Haar measures $\rho'_j,\mu'_j$ on the closed subgroups
$R_j\cap B_j$ and $B_j$ (they are locally compact by $\mathsf{GrpAt}$). By $\mathsf{RepAt}(s)$ choose representative points
$b_t(\beta)$; they lie in $\prod B_j$, so their weights are $w(b_t(\beta))=\prod_{j\in I}\rowpropscw(b_t(\beta)_j)$. Index the data by
$x=(t,a,\tau,M,(M_\beta,\tau_\beta),y)$, with $S_x=\rowpropsSupp(t,\Phi_a)$ (finite by $\mathsf{SumOK}$), $b_x=b_t$,
$q_x=\rowpropstot(P,\tau;\phi_a,F_{t,y})$, $L_x=\mathrm{lhs}(\delta_P,\tau;a,y)$ and $r(x)=z_t^{-1}\sum_\beta
g(t,a,\beta,M_\beta,\tau_\beta,y)$. The hypotheses of \cref{pf:rowprops:propcell} hold: distinct cells by $\mathsf{RepAt}(s)$,
bi-invariance by $\mathsf{InvAt}$, the support condition by $\mathsf{SuppOn}(s)$, and $L_x$ is an iterated integral of $q_x$
by $\mathsf{IterAt}$. For the relation of cell sums, let $p$ be real; $\mathsf{ValAt}(s)$ gives $\sigma$ depending only on
$(\rho',\mu',P,p)$ with $w(b_t(\beta))q_x(p\,b_t(\beta))=z_t^{-1}(\prod_w\mathrm{coef}_w(\beta_w))\rowpropsdef(M_\beta;\sigma,\tau_\beta;\phi_a,F_{t,y})
=\ell(\sigma)\,z_t^{-1}g(t,a,\beta,M_\beta,\tau_\beta,y)$ by \eqref{eq:pf:rowprops:hdef}. Summing over $\beta$ gives the relation with
$e(p)=\ell(\sigma)$. \Cref{pf:rowprops:propcell} gives $c_P$.

\emph{Step 2.} The scalar $s_{AA}$ is non-zero: it is a product of powers of $\Nm(w)^{1/2}$
(\lref{GlobalTheory.FunObj.aaScalar_eq_pow}). Given $\tau_i$, choose by $\mathsf{RealLine}$ a $\sigma$ with
$\ell(\sigma)=s_{AA}^{-1}\sum_P\tau_i(P)c_P$, and let $z'=z$. Then by $\mathsf{LinAt}$, Step 1 and \eqref{eq:pf:rowprops:hdef},
\begin{multline*}
\mathrm{lhs}(\tau_i,\tau;a,y)=\sum_P\tau_i(P)\,c_P\,z_t^{-1}\sum_\beta g(t,a,\beta,M_\beta,\tau_\beta,y)\\
=s_{AA}\,\ell(\sigma)\,z_t^{-1}\sum_\beta g(t,a,\beta,M_\beta,\tau_\beta,y)
=z_t^{-1}s_{AA}\sum_\beta\Bigl(\prod_w\mathrm{coef}_w(\beta_w)\Bigr)\rowpropsdef(M_\beta;\sigma,\tau_\beta;\phi_a,F_{t,y}),
\end{multline*}
which is \eqref{eq:pf:rowprops:unfoldat}. The data $\sigma$, $z$ do not depend on $t,a,\tau,M,y$.
\end{proof}

\subsubsection{Values on the cells}
The value property is derived from seven further properties of the choice.

\begin{definition}[Properties entering the cell values]\label{def:pf:rowprops:valprops}
Let $\mathrm{ch}$ be an admissible choice. A family $C$ of functions of two consecutive top-column arguments is a
\emph{real family} if at the even indices $j<m$ it consists of smooth bi-vacuum vectors for the top lattices; a
\emph{combination} of an expansion $\tau$ is a family of finite combinations of half-integral points for every datum in its
support, representing the finite data of $\tau$ through the vacuum vectors (and acting on real families by the
corresponding finite linear combinations of translates, $\mathrm{comb}_C$, $\mathrm{comb}_R$).
\begin{enumerate}[label=(\alph*)]
\item $\mathsf{FlatAt}(s)$ (\emph{flat dictionary}): there is $z^{(1)}$ with $z^{(1)}_t\in\rowpropsXi_t$ such that for every
family $P_2$ of down polynomials and every real point $p$ there is a real family $C$ such that for all $t$, $\Phi$, flipped
finite tensors $\tau$, $M$ of level $t$, $\beta$ with representative point $b$, expansions $\tau'$ of $\beta$ with
combination $cc$, there is $\zeta$, equal to $(z^{(1)}_t)^{-1}$ if $\beta$ is a staircase family, with
\[
\sum_d\tau(d)\,\mathrm{ev}\bigl(\mathrm{con}(\mathrm{fock}(d_2,P_2);p\,b),d_1\bigr)=\zeta\sum_d\tau'(d)\,
\mathrm{ev}\bigl(\mathrm{comb}_C(cc(d))C,\mathrm{comb}_R(cc(d))\bigr)
\]
as functions of all real arguments.
\item $\mathsf{TensAt}$: for Haar volume data and continuous automorphic $\phi$, $F$, a linear relation among the product
functions of finitely many good pairs $(C_k,f_k)$ implies the same relation among the periods
$\rowpropsper(C_k,f_k,P_1;\phi,F)$.
\item $\mathsf{VolAt}$: if the measures of $M$ at the definite nodes other than $\mathrm t_m$ are $s(X)$ times those of $M'$,
then $\rowpropsper_M=\mathrm{sp}(s)\,\rowpropsper_{M'}$ for good pairs, where
$\mathrm{sp}(s)=\prod_{j<m}s(\mathrm t_j)\prod_{i\le m}s(i,m)$.
\item $\mathsf{WtAt}(s)$ (\emph{cell masses}): for all Haar $\rho'$, $\mu'$ there is $c$ such that for all $t,\Phi,\beta$ with
representative point $b$, all $M$ of level $t$ and term volumes $M'$ for $\beta$, there is a scaling $x$ from $M'$ to $M$ with
$w(b)\,\mathrm{sp}(x)=c\prod_w\mathrm{coef}_w(\beta_w)$.
\item $\mathsf{FockAt}$ (\emph{reduction to Fock vectors}): there is $z^{(2)}$ with $z^{(2)}_t\in\rowpropsXi_t$ such that for
every real family $C$ and row polynomials $P_1$ there is $\sigma_0$ such that for all $t,a,\beta$, Haar $M$, with
$\mathsf{SumOK}(t,\Phi_a)$, expansions $\tau$ with combination $cc$, and all $y$:
\[
\sum_d\tau(d)\rowpropsper\bigl(\mathrm{comb}_C(cc(d))C,\mathrm{comb}_R(cc(d)),P_1;\phi_a,F_{t,y}\bigr)
=(z^{(2)}_t)^{-1}\rowpropsdef(M;\sigma_0,\tau;\phi_a,F_{t,y}).
\]
\item $\mathsf{BasicAt}$: the functions $\phi_a$ and $F_{t,y}$ are continuous and automorphic; volume data of level $t$ and term
volumes are Haar measures with fundamental domains at the definite nodes; the pairs occurring in the last-row integrand
and the combinations of real families are good; and every expansion has a combination.
\item $\mathsf{OneAt}$: every $\beta$ in a support is a staircase family (at the places of $W_0$ its cosets are those of the
staircase representatives).
\end{enumerate}
\formal{GlobalTheory.RowVal.FlatAt}\formal{GlobalTheory.RowVal.TensAt}\formal{GlobalTheory.RowVal.VolAt}\formal{GlobalTheory.RowVal.WtAt}\formal{GlobalTheory.RowVal.FockAt}\formal{GlobalTheory.RowVal.BasicAt}\formal{GlobalTheory.RowVal.OneAt}
\end{definition}

\begin{proposition}[Cell values]\label{pf:rowprops:valat}
For an admissible choice, $\mathsf{FlatAt}(s)$, $\mathsf{TensAt}$, $\mathsf{VolAt}$, $\mathsf{WtAt}(s)$, $\mathsf{FockAt}$,
$\mathsf{BasicAt}$, $\mathsf{OneAt}$ and $\mathsf{SumOK}(t,\Phi_a)$ for all $t,a$ imply $\mathsf{ValAt}(s)$, with
$z_t=z^{(1)}_tz^{(2)}_t$.
\formallabel{GlobalTheory.RowVal.valAt_of_steps}
\formallabel{GlobalTheory.RowVal.valAt_big}
\formallabel{GlobalTheory.T.rowTest_cellValue_of_steps_large_depth}
\end{proposition}

\begin{proof}
Clearly $z_t\in\rowpropsXi_t$. Let $\rho',\mu'$, $P=(P_1,P_2)$ and a real point $p$ be given. Let $c_0$ be given by
$\mathsf{WtAt}(s)$, $C$ by $\mathsf{FlatAt}(s)$ for $(P_2,p)$, and $\sigma_0$ by $\mathsf{FockAt}$ for $(C,P_1)$; put
$\sigma=c_0\sigma_0$. Let $t,a,\tau,M,\beta,b,M',\tau',y$ be as in $\mathsf{ValAt}(s)$, and write $\phi=\phi_a$,
$F=F_{t,y}$. By $\mathsf{BasicAt}$, $\tau'$ has a combination $cc$; $\mathsf{FlatAt}(s)$ gives $\zeta$, and $\zeta=(z^{(1)}_t)^{-1}$
by $\mathsf{OneAt}$. $\mathsf{WtAt}(s)$ gives a scaling $x$ from $M'$ to $M$ with $w(b)\,\mathrm{sp}(x)=c_0\prod_w\mathrm{coef}_w(\beta_w)$.
\begin{enumerate}[label=(\arabic*)]
\item Apply $\mathsf{TensAt}$ to the finite family consisting of the pairs $(\mathrm{con}(\mathrm{fock}(d_2,P_2);pb),d_1)$ with
coefficients $\tau(d)$ and the pairs $(\mathrm{comb}_C(cc(d))C,\mathrm{comb}_R(cc(d)))$ with coefficients $-\zeta\tau'(d)$; all
pairs are good by $\mathsf{BasicAt}$, and the linear relation among their product functions is the identity of
$\mathsf{FlatAt}(s)$. Hence
\[
\rowpropstot(P,\tau;\phi,F)(pb)=\zeta\sum_d\tau'(d)\,\rowpropsper_M\bigl(\mathrm{comb}_C(cc(d))C,\mathrm{comb}_R(cc(d)),P_1;\phi,F\bigr).
\]
\item By $\mathsf{VolAt}$, the sum on the right is $\mathrm{sp}(x)$ times the same sum with $M'$.
\item By $\mathsf{FockAt}$ (with $M'$, which is Haar by $\mathsf{BasicAt}$), the sum with $M'$ is equal to
\[
(z^{(2)}_t)^{-1}\rowpropsdef(M';\sigma_0,\tau';\phi,F).
\]
\end{enumerate}
Therefore $w(b)\,\rowpropstot(P,\tau;\phi,F)(pb)=(z^{(1)}_t)^{-1}(z^{(2)}_t)^{-1}w(b)\,\mathrm{sp}(x)\,\rowpropsdef(M';\sigma_0,\tau';\phi,F)
=z_t^{-1}(\prod_w\mathrm{coef}_w(\beta_w))\,\rowpropsdef(M';c_0\sigma_0,\tau';\phi,F)$, using the linearity of $\rowpropsdef$ in $\sigma$.
\end{proof}

\subsubsection{Five of the unfolding properties}
We prove $\mathsf{IterAt}$, $\mathsf{LinAt}$, $\mathsf{FlatAt}$, $\mathsf{WtAt}$, $\mathsf{FockAt}$ and $\mathsf{RealLine}$ from
deeper statements of the development.

\begin{proposition}[Iterated integrals and linearity]\label{pf:rowprops:iterat}
Assume $F$ and $\sfname{LiftsCuspidal}$. At large scale, every admissible choice has the properties $\mathsf{IterAt}$ and
$\mathsf{LinAt}$.
\formallabel{GlobalTheory.T.twisted_iterate_chain}
\formallabel{GlobalTheory.RowChain.twisted_iterate_eq_linear}
\formallabel{GlobalTheory.RowChain.twisted_iterate_eventually_large}
\formallabel{GlobalTheory.RowVal.placeValuation_linAt_large}
\end{proposition}

\begin{proof}
At large scale we have the cusp step for every volume datum of level $t$ (\cref{pf:rowprops:cuspstep}), the forms $F_{t,y}$
are cuspidal listed at $(0,m)$ (\lref{GlobalTheory.EndCusp.leadingForm_mem_cuspList}), and the scalar column
functions $\phi_a$ are continuous and left invariant under rational points (they are built from columns,
\cref{def:pf:rowprops:columns}). Fix $t$, $a$, a flipped finite tensor $\tau$ for $(t,\Phi_a)$ (with idempotent $e$), $M$ of
level $t$, $P$ and $y$; write $\mathrm{opp}(P,d)$ for the pure tensor along $\rowpropsopp_t$ attached to $P$ and a pure datum $d$.

$\mathsf{IterAt}$. Choose Haar measures $\rho_j$ on the row groups $R_j$ and $\mu_j$ on $H_j$ ($j\in I$) compatible with the
measures of volume data of level $t$ at the nodes $(m,j)$ (they exist: \lref{GlobalTheory.RowChain.exists_rowMeas} and
\lref{GlobalTheory.RowChain.exists_vol_row}). For $d$ in the support of $\tau$, the components of $d$ are local vectors at
$e$, and \lref{GlobalTheory.RowChain.pure_core} (which uses the rapid decay of cusp forms, that $M$ is of row type, and
the cusp step) shows: the lift $\rowpropsL_{\rowpropsopp_t}(\mathrm{opp}(P,d))F_{t,y}$ is an automorphic function, so that its
Bessel pairing with $\phi_a$ converges (\lref{GlobalTheory.RowBplus.bplusConv_col}); and this Bessel pairing is the
iterated quotient integral over $I$ of $h\mapsto\rowpropsper(\mathrm{con}(\mathrm{fock}(d_2,P_2);h),d_1,P_1;\phi_a,F_{t,y})$, by
the identity between column periods and iterated integrals of the contracted kernels
(\lref{GlobalTheory.RowVal.nIter_conAt_of_colPer}). Since $\rowpropsL_{\rowpropsopp_t}(\mathrm{opp}(\delta_P,\tau))F_{t,y}=
\sum_d\tau(d)\rowpropsL_{\rowpropsopp_t}(\mathrm{opp}(P,d))F_{t,y}$ and $\rowpropsB$ is linear on convergent arguments,
$\mathrm{lhs}(\delta_P,\tau;a,y)=\sum_d\tau(d)\,\rowpropsB(\phi_a,\rowpropsL_{\rowpropsopp_t}(\mathrm{opp}(P,d))F_{t,y})$, and by
\cref{pf:rowprops:niterlin} this is an iterated quotient integral of $\sum_d\tau(d)\rowpropsper(\dots)=\rowpropstot(P,\tau;\phi_a,F_{t,y})$.

$\mathsf{LinAt}$. Each $\rowpropsL_{\rowpropsopp_t}(\mathrm{opp}(P,d))F_{t,y}$ is continuous and automorphic
(\lref{GlobalTheory.RowVal.opp_contAut}), hence so is every finite sum of them; and for a volume datum of row type the
Bessel pairing with $\phi_a$ is additive on finite sums of continuous automorphic functions
(\lref{GlobalTheory.RowVal.bplus_fsum}). Since $\rowpropsL_{\rowpropsopp_t}(\mathrm{opp}(\tau_i,\tau))F_{t,y}=
\sum_P\tau_i(P)\rowpropsL_{\rowpropsopp_t}(\mathrm{opp}(\delta_P,\tau))F_{t,y}$, $\mathsf{LinAt}$ follows.
\end{proof}

\begin{definition}[Local dictionary; real vacua]\label{def:pf:rowprops:locvac}
Let $\mathrm{ch}$ be an admissible choice. A \emph{base point} for $t$ is a point $b_0$ trivial outside $\{0\}\cup I$ whose
component at $0$ has local components the inverse staircase elements at $\mathrm{loc}(t)$ and whose components at $j\in I$
are finite parts of rational tilts. Let $v^{\mathrm{rm}}$ denote the family of vacuum vectors of the real match on the last
column.
\begin{enumerate}[label=(\alph*)]
\item $\mathsf{LocAt}(s)$: there are $\kappa\in\C$ and $z^{(1)}$ with $z^{(1)}_t\in\rowpropsXi_t$ such that for all $t$, $\Phi$,
flipped finite tensors $\tau$ for $(t,\Phi)$, $M$ of level $t$, $\beta$ with representative point $b$, base points $b_0$,
expansions $\tau'$ of $\beta$ with combination $cc$, there is $\zeta$, equal to $\kappa(z^{(1)}_t)^{-1}$ if $\beta$ is a staircase
family, such that for every family $P_2$ of down polynomials and every real point $p$
\begin{multline*}
\sum_d\tau(d)\,\mathrm{ev}\bigl(\mathrm{con}(\mathrm{fock}(d_2,P_2);pb),d_1\bigr)\\
=\zeta\sum_d\tau'(d)\,\mathrm{ev}\bigl(\mathrm{comb}_C(cc(d))\,\mathrm{con}(\mathrm{fock}(v^{\mathrm{rm}},P_2);pb_0),
\mathrm{comb}_R(cc(d))\bigr).
\end{multline*}
\item $\mathsf{VacAt}$: base points exist for every $t$, and there is $z^{(2)}(t,j)\in\rowpropsXi_t$ such that for every family
$P_2$ of down polynomials and every real point $p$ there is a real family $C$ with
$\mathrm{con}(\mathrm{fock}(v^{\mathrm{rm}},P_2);pb_0)_j=z^{(2)}(t,j)\,C_j$ for all $t$, $M$ of level $t$, base points $b_0$
and even $j<m$.
\end{enumerate}
\formal{GlobalTheory.RowVal.LocAt}\formal{GlobalTheory.RowVal.VacAt}\formal{GlobalTheory.RowVal.IsBase}
\end{definition}

\begin{proposition}[Flat dictionary]\label{pf:rowprops:flatat}
At large scale, every admissible choice satisfies $\mathsf{FlatAt}(s)$ for some sign $s$.
\formallabel{GlobalTheory.RowFlatAsm.flatAt_big}
\formallabel{GlobalTheory.RowFlatAsm.flatAt_of}
\end{proposition}

\begin{proof}
At large scale the choice satisfies $\mathsf{LocAt}(s)$ for some $s$ (\lref{GlobalTheory.RowVal.exists_locAt_big}: at a
cell point, the local tensor of the integrand is a constant multiple of the expansion tensor built from the vacua at a base
point), $\mathsf{VacAt}$ (\lref{GlobalTheory.RowVal.OneVac.big_vacAt}: at a base point, the even components of the
contracted vacuum family are real families up to elements of $\rowpropsXi_t$), and $\mathsf{OneAt}$
(\lref{GlobalTheory.RowVal.oneAt}). Let $\kappa$, $z^{(1)}$ and $z^{(2)}$ be as in these properties and put
$z_t=z^{(1)}_t\bigl(\prod_{k<m/2}z^{(2)}(t,2k)\bigr)^{-1}\in\rowpropsXi_t$. Given $P_2$ and a real point $p$, let $C^0$ be the
real family of $\mathsf{VacAt}$ and $C$ the family with $C_0=\kappa C^0_0$ and $C_j=C^0_j$ for $j\neq0$; it is again a real
family. Let $t,\Phi,\tau,M,\beta,b,\tau',cc$ be as in $\mathsf{FlatAt}(s)$ and choose a base point $b_0$.

If $\beta$ is a staircase family, $\mathsf{LocAt}(s)$ gives the identity of \cref{def:pf:rowprops:locvac}(a) with
$\zeta=\kappa(z^{(1)}_t)^{-1}$. Substitute $\mathrm{con}(\mathrm{fock}(v^{\mathrm{rm}},P_2);pb_0)_j=z^{(2)}(t,j)C^0_j$ for the even $j$
(the product function involves only the even components): the product function of a combination of a family is multilinear
in its even components (\lref{GlobalTheory.RowFlatAsm.evT_combC_smul}), so the right side becomes
\begin{multline*}
\kappa(z^{(1)}_t)^{-1}\prod_kz^{(2)}(t,2k)\sum_d\tau'(d)\,\mathrm{ev}\bigl(\mathrm{comb}_C(cc(d))C^0,\mathrm{comb}_R(cc(d))\bigr)\\
=z_t^{-1}\sum_d\tau'(d)\,\mathrm{ev}\bigl(\mathrm{comb}_C(cc(d))C,\mathrm{comb}_R(cc(d))\bigr),
\end{multline*}
which is the identity of $\mathsf{FlatAt}(s)$ with $\zeta=z_t^{-1}$.

If $\beta$ is not a staircase family, then $\beta$ is not in any support by $\mathsf{OneAt}$, and the zero combination is an
expansion of $\beta$ (\lref{GlobalTheory.RowFlatAsm.isExpansion_zero}). Applying $\mathsf{LocAt}(s)$ to the zero expansion
shows that the left side of the identity of $\mathsf{FlatAt}(s)$ vanishes, so $\zeta=0$ works (there is no condition on $\zeta$
in this case).
\end{proof}

\begin{proposition}[Cell masses]\label{pf:rowprops:wtat}
At large scale, every admissible choice satisfies $\mathsf{WtAt}(s)$ for every sign $s$.
\formallabel{GlobalTheory.T.cell_masses_against_level_volumes_large_depth}
\formallabel{GlobalTheory.RowWt.wtAt_big}
\end{proposition}

\begin{proof}
Two statements hold at large scale.
(A) For every $t$, $\Phi$, $\beta$, every $M$ of level $t$ and every term volume $M'$ for $\beta$, there is a scaling $x$ from
$M'$ to $M$ (that is, $\nu_X(M)=x(X)\,\nu_X(M')$ at the definite nodes $X\ne\mathrm t_m$) with $x(i,m)=1$ for all $i\le m$
and $x(\mathrm t_j)=1$ for $j\notin I$ (\lref{GlobalTheory.T.rowTest_tupleTerm_measures_differ_by_scalars}).
(C) For all Haar measures $\rho',\mu'$ and every $j\in I$ there is a constant $c_j$ such that for all $t,\Phi,\beta$ with
representative point $b$, all $M$, $M'$ and every scaling $x$ as in (A),
$\rowpropscw(b_j)\,x(\mathrm t_j)=c_j\,\mathrm{coef}^{\mathrm{col}}(t,\beta,j)$, where the numbers $\mathrm{coef}^{\mathrm{col}}(t,\beta,j)$ are the
contributions of the column $j$ to the local coefficients:
$\prod_{j\in I}\mathrm{coef}^{\mathrm{col}}(t,\beta,j)=\prod_w\mathrm{coef}_w(\beta_w)$
(\lref{GlobalTheory.T.unfolding_cell_mass_eq_const_mul_local_coefficients}, \lref{GlobalTheory.RowWt.coef_prod}).

Given $\rho',\mu'$, put $c=\prod_{j\in I}c_j$. For $t,\Phi,\beta,b,M,M'$ take $x$ from (A). The components $b_j$, $j\in I$, of a
representative point are finite points, so $w(b)=\prod_{j\in I}\rowpropscw(b_j)$; and
$\mathrm{sp}(x)=\prod_{j\in I}x(\mathrm t_j)$ because the other factors are $1$. Hence
$w(b)\,\mathrm{sp}(x)=\prod_{j\in I}\rowpropscw(b_j)x(\mathrm t_j)=\prod_{j\in I}c_j\,\mathrm{coef}^{\mathrm{col}}(t,\beta,j)=c\prod_w\mathrm{coef}_w(\beta_w)$.
\end{proof}

\begin{definition}[Moved vectors and column Fock combinations]\label{def:pf:rowprops:moved}
Let $\mathrm{ch}$ be an admissible choice.
\begin{enumerate}[label=(\alph*)]
\item $\mathsf{MovedAt}$: there are column vacuum vectors $v=(v_j)$, non-zero constants $\kappa_j$, $\kappa'_i$ and
$z(t,j)\in\rowpropsXi_t$ such that for all $t$, $a$, $\beta$, all expansions $\tau$ of $\beta$ with combination $cc$ and all $d$
in the support of $\tau$, the moved vectors of $d$ (\cref{def:pf:rowprops:flipfin}(c)) are $\kappa_jz(t,j)\,\mathrm{comb}(cc(d)_{1,j})v_j$
on the column and $\kappa'_i\,\mathrm{comb}_R(cc(d))_i$ on the last row.
\item $\mathsf{ColFock}$: for every family $v$ of column vacuum vectors and every real family $C$ there is a finite combination
$\sigma_c$ of column polynomial families such that for every $t$, Haar volume datum $M$, $a$, finite combination $c$ of points
and every row datum $f$ of theta functions with polynomials $P_1$ whose row lift of $F_{t,y}$ is continuous and automorphic,
$\rowpropsper(\mathrm{comb}_C(c)C,f,P_1;\phi_a,F_{t,y})$ equals $\sum_Q\sigma_c(Q)$ times the integral over
$\rowpropsfd_{m,m}$ of the column lift of $\phi_a$ with the Fock vectors $\mathrm{fock}(\mathrm{comb}(c_j)v_j,Q_j)$, composed with
$\iota_{\mathrm t}$, against the row lift of $F_{t,y}$ with $\mathrm{fock}(f_i,P_{1,i})$.
\end{enumerate}
\formal{GlobalTheory.RowVal.MovedAt}\formal{GlobalTheory.RowVal.ColFock}
\end{definition}

\begin{proposition}[Reduction to Fock vectors]\label{pf:rowprops:fockat}
At large scale, every admissible choice satisfies $\mathsf{FockAt}$.
\formallabel{GlobalTheory.T.periodReductionToFockVectors_large_depth}
\formallabel{GlobalTheory.FockRed.fockAt_big}
\formallabel{GlobalTheory.FockRed.fockAt_of_steps}
\end{proposition}

\begin{proof}
At large scale the choice satisfies $\mathsf{BasicAt}$ (\lref{GlobalTheory.T.lastRow_unfolding_basicHypotheses}),
$\mathsf{MovedAt}$ (\lref{GlobalTheory.T.movedVectorsAreCombinationsOfVacua_of_large_depth}) and $\mathsf{ColFock}$
(\lref{GlobalTheory.T.columnRealDatum_fockCombination_largeDepth}); and for every Haar volume datum and every row datum of
theta functions the row lift of $F_{t,y}$ along $\rho_m$ is continuous and automorphic
(\lref{GlobalTheory.RowTens.row_invCont}, using that the spaces of the last row are positive definite and that $F_{t,y}$
is continuous and automorphic by $\mathsf{BasicAt}$).

Let $v,\kappa,\kappa',z$ be as in $\mathsf{MovedAt}$ and put $z^{(2)}_t=\prod_jz(t,j)\in\rowpropsXi_t$ and
$K=(\prod_j\kappa_j)^{-1}(\prod_i\kappa'_i)^{-1}$. Given a real family $C$ and row polynomials $P_1$, let $\sigma_c$ be given by
$\mathsf{ColFock}$ for $(v,C)$, and put $\sigma_0=K\cdot\sum_Q\sigma_c(Q)\,\delta_{(Q,P_1)}$. Let $t,a,\beta,M$, an expansion $\tau$
with combination $cc$ and $y$ be given. For $d$ in the support of $\tau$, $\mathsf{ColFock}$ (with the combination $cc(d)$, whose
row part is a datum of theta functions by $\mathsf{BasicAt}$) gives
\[
\rowpropsper\bigl(\mathrm{comb}_C(cc(d))C,\mathrm{comb}_R(cc(d)),P_1;\phi_a,F_{t,y}\bigr)=\sum_Q\sigma_c(Q)\,(\dots),
\]
and by $\mathsf{MovedAt}$ and the multilinearity of the definite period in the Fock vectors of its column and row lifts
(\lref{GlobalTheory.FockRed.defPer_of_smul}) each integral $(\dots)$ equals
$(z^{(2)}_t)^{-1}K\,\rowpropsdef(M;(Q,P_1),d;\phi_a,F_{t,y})$: replacing the vectors $\mathrm{comb}(cc(d)_{1,j})v_j$ and
$\mathrm{comb}_R(cc(d))_i$ by the moved vectors of $d$ multiplies the integral by $\prod_j\kappa_jz(t,j)\prod_i\kappa'_i$. Multiplying by
$\tau(d)$ and summing over $d$ gives $\sum_d\tau(d)\rowpropsper(\dots)=(z^{(2)}_t)^{-1}\rowpropsdef(M;\sigma_0,\tau;\phi_a,F_{t,y})$,
which is $\mathsf{FockAt}$.
\end{proof}

\begin{definition}[Invariant lines]\label{def:pf:rowprops:invline}
For a finite set $\iota$ (of real places), sizes $S$ (ranks of a row chain and a column chain of compact unitary groups) and
good entries $P=(P_\tau)_{\tau\in\iota}$ (archimedean weights and exponents), let $\mathrm{OM}(P)$ be the tensor product of the
Fock spaces of the row chain and of the column chain with two irreducible algebraic representations at the ends, as a
representation of the product $\mathrm{Gc}$ of the compact unitary groups of the chains over $\iota$, and $\mathrm{inv}(P)$ the
space of $\mathrm{Gc}$-invariant linear forms on $\mathrm{OM}(P)$.
\begin{enumerate}[label=(\alph*)]
\item $\mathsf{InvLine}$ for a choice: there are $\iota,S,P$ as above, with all entries good, and a map $e$ from polynomial
pairs to $\mathrm{OM}(P)$, such that for all $t,a$ with $\mathsf{SumOK}(t,\Phi_a)$, $\beta$ in the support, term volumes $M$,
expansions $\tau$, $d$ in the support of $\tau$ and $y$, there is $l\in\mathrm{inv}(P)$ with
$\rowpropsdef(M;Q,d;\phi_a,F_{t,y})=l(e(Q))$ for all polynomial pairs $Q$.
\item $\mathsf{OneLine}$: for all $\iota,S$ and good $P$, $\dim\mathrm{inv}(P)=1$.
\item $\mathsf{PeriodsOnLine}$ for a choice: there is a function $\ell$ on polynomial pairs such that for all
$t,a,\beta,M,\tau,d,y$ as in (a) there is $\kappa\in\C$ with $\rowpropsdef(M;Q,d;\phi_a,F_{t,y})=\kappa\,\ell(Q)$ for all $Q$.
\end{enumerate}
\formal{GlobalTheory.RowReal.InvLine}\formal{GlobalTheory.RowReal.PeriodsOnLine}
\end{definition}

\begin{proposition}[The definite periods lie on one line]\label{pf:rowprops:realline}
At large scale, every admissible choice satisfies $\mathsf{RealLine}$.
\formallabel{GlobalTheory.RowReal.realLine}
\formallabel{GlobalTheory.RowReal.realLine_of_periodsOnLine}
\formallabel{GlobalTheory.RowReal.periodsOnLine_big}
\formallabel{GlobalTheory.RowReal.periodsOnLine_of_invLine}
\formallabel{GlobalTheory.T.rowTest_realPlaces_definitePeriods_onOneLine}
\end{proposition}

\begin{proof}
\emph{Step 1.} At large scale every admissible choice satisfies $\mathsf{InvLine}$
(\lref{GlobalTheory.RowReal.invLine_big}): the archimedean part of a definite period is an invariant linear form on a
fixed tensor product of Fock spaces and algebraic representations, evaluated at the image of the polynomial pair.

\emph{Step 2.} $\mathsf{OneLine}$ holds (\lref{GlobalTheory.OneFunctional.oneLine}). This is a multiplicity-one statement
for branching along the chains, proved by choosing, along the row chain and the column chain, branching weights of
Gan--Gross--Prasad type attached to the entries
(\lref{GlobalTheory.T.exists_ganGrossPrasad_row_weights}, \lref{GlobalTheory.T.exists_ganGrossPrasad_column_weights}), contracting
the tensor products of Fock spaces along the chains step by step to the representations at the ends
(\lref{GlobalTheory.T.contracts_tensorProduct_chain}), and identifying the end representations with Weyl modules of
positive weight (\lref{GlobalTheory.T.outer_irr_iso_weylRep_of_pos}); from these identifications the dimension of
$\mathrm{inv}(P)$ is computed to be one.

\emph{Step 3: $\mathsf{InvLine}$ and $\mathsf{OneLine}$ imply $\mathsf{PeriodsOnLine}$.} With $(\iota,S,P,e)$ from $\mathsf{InvLine}$,
let $L_0$ span $\mathrm{inv}(P)$ and put $\ell(Q)=L_0(e(Q))$. For given $t,a,\beta,M,\tau,d,y$, the form $l$ of $\mathsf{InvLine}$
is $\kappa L_0$ for some $\kappa$, so $\rowpropsdef(M;Q,d;\phi_a,F_{t,y})=\kappa\,\ell(Q)$.

\emph{Step 4: $\mathsf{PeriodsOnLine}$ implies $\mathsf{RealLine}$.} Let $\ell_0$ be given by $\mathsf{PeriodsOnLine}$. If $\ell_0=0$,
all definite periods vanish; put $\ell\equiv1$ (non-zero), and take $\kappa=0$ below. Otherwise put $\ell=\ell_0$. Then
$\ell\ne0$, and every definite period of a pure datum is $\kappa_d\ell(Q)$. For an expansion $\tau$ and a finite combination
$\sigma$, $\rowpropsdef(M;\sigma,\tau;\phi_a,F_{t,y})=\sum_d\tau(d)\sum_Q\sigma(Q)\kappa_d\ell(Q)=\ell(\sigma)\,K$ with
$K=\sum_d\tau(d)\kappa_d$, which does not depend on $\sigma$. Finally, if $\ell(Q_0)\ne0$, then $\sigma=(z/\ell(Q_0))\delta_{Q_0}$
satisfies $\ell(\sigma)=z$, and $\sigma=0$ if $z=0$.
\end{proof}

\subsubsection{Assembly of the unfolding}
\begin{proposition}[Unfolding]\label{pf:rowprops:unfoldmain}
Assume $F$, $\sfname{LiftsCuspidal}$ and $\sfname{SumOKK}$. Then $\sfname{StepUnfold}$ holds.
\formallabel{GlobalTheory.T.twisted_theta_integral_unfolds}
\formallabel{GlobalTheory.RowUnfLast.twisted_theta_integral_unfolding}
\formallabel{GlobalTheory.T.lastRow_real_iterated_integral_eq_const_mul_finite}
\formallabel{GlobalTheory.RowRealInt.unfold_of_steps}
\formallabel{GlobalTheory.RowReal.big_mono}
\end{proposition}

\begin{proof}
\Cref{pf:rowprops:cellvan} holds unconditionally. At large scale, every admissible choice has the following properties:
\begin{itemize}
\item $\mathsf{GrpAt}$ (\lref{GlobalTheory.RowGrpSplit.grpAt_big}): the archimedean and finite parts of $G_{V_{m,j}}(\A)$ and
the row groups and lattice groups form split quadruples;
\item $\mathsf{IterAt}$ and $\mathsf{LinAt}$ (\cref{pf:rowprops:iterat}, which uses $F$ and
$\sfname{LiftsCuspidal}$): the Bessel pairing of a column with the lift along the opposite path is the iterated integral of
the last-row integrand over the row groups, and it is linear in $\tau_i$;
\item $\mathsf{InvAt}$ (\lref{GlobalTheory.T.unfolded_integrand_biInvariant_large_depth});
\item $\mathsf{FlatAt}(s)$ for some sign $s$ (\cref{pf:rowprops:flatat});
\item $\mathsf{OffAt}$ (\lref{GlobalTheory.T.singleCosetOffTwistPlaces_largeDepth}): if the last-row integrand of a pure
datum is non-zero at a point $h$ with $h_0=1$, then for $j\in I$ the component $h_j$ lies in $R_jA_j\,\delta x\,K_j$ for a
rational tilt $\delta$ and an element $x$ with prescribed local components at $\mathrm{loc}(t)$;
\item $\mathsf{TensAt}$ (\lref{GlobalTheory.T.definitePeriod_depends_on_productFunction}), $\mathsf{VolAt}$
(\lref{GlobalTheory.T.rowTest_definitePeriod_change_of_measures}), $\mathsf{BasicAt}$
(\lref{GlobalTheory.T.lastRow_unfolding_basicHypotheses}), $\mathsf{OneAt}$ (\lref{GlobalTheory.RowVal.oneAt}),
$\mathsf{WtAt}(s)$ for all $s$ (\cref{pf:rowprops:wtat}), $\mathsf{FockAt}$ (\cref{pf:rowprops:fockat}) and $\mathsf{RepAt}(s)$ for all $s$
(\lref{GlobalTheory.T.representation_at_every_place});
\item $\mathsf{SuppOn}(s)$, given $\mathsf{FlatAt}(s)$, $\mathsf{OffAt}$, $\mathsf{InvAt}$, $\mathsf{TensAt}$ and $\mathsf{BasicAt}$
(\lref{GlobalTheory.T.integrand_supported_on_cells_largeDepth});
\item $\mathsf{RealLine}$ (\cref{pf:rowprops:realline});
\item $\mathsf{SumOK}(t,\Phi_a)$ for all $t,a$ ($\sfname{SumOKK}$).
\end{itemize}
Taking $s$ from $\mathsf{FlatAt}$, \cref{pf:rowprops:valat} gives $\mathsf{ValAt}(s)$, and \cref{pf:rowprops:unfoldAt} gives
$\mathsf{UnfoldAt}$. The combination of these statements at large scale is the monotonicity of $\rowpropsBig$ under
implication together with \lref{GlobalTheory.RowClaim.Big.and}.
\end{proof}

% ---------------------------------------------------------------------------
\subsection{Jets: definite periods as class sums}\label{sec:pf:rowprops:jets}

On a totally definite unitary group the archimedean group is compact, and the functions occurring in the definite
periods are matrix coefficients of algebraic representations at the archimedean places. The theta kernels act on these
algebraic data by explicit linear maps; we call the vectors so obtained \emph{jets}.

\begin{definition}[Finite-real functions; good arrows]\label{def:pf:rowprops:jetdata}
Let $V_0,V_1,\dots$ be hermitian spaces, $\varrho_n$ representations of $G_{V_n}(\A)$ on finite-dimensional spaces
$\mathcal M_n$, $\Lambda_n$ linear forms on $\mathcal M_n$, $\ker_n$ functions on $G_{V_n}(\A)\times G_{V_{n+1}}(\A)$, and
$\mathrm{comp}_n(x,y):\mathcal M_n\to\mathcal M_{n+1}$ linear maps depending on finite points $x$, $y$.
\begin{enumerate}[label=(\alph*)]
\item A function $F$ on $G_{V_0}(\A)$ is \emph{finite-real} for $(\varrho_0,\Lambda_0,v)$, where $v$ is an
$\mathcal M_0$-valued function on the finite points, if $F(h)=\Lambda_0\bigl(\varrho_0(h)^{-1}v(h_{\mathrm f})\bigr)$ for all $h$.
\item The arrow $n$ is \emph{good} if $V_n$ is totally definite (positive definite at every complex embedding), $\varrho_n$
is trivial on the finite points, the matrix coefficients $k\mapsto\Lambda_n(\varrho_n(k)^{-1}u)$ are continuous on
$G_{V_n,\infty}$, the kernel $\ker_n$ is invariant under rational points on both sides and continuous in the archimedean
variable, and for every Haar probability measure $\kappa$ on $G_{V_n,\infty}$, every finite point $x$, every $h$ and every
$u\in\mathcal M_n$,
\begin{equation}\label{eq:pf:rowprops:arrowok}
\int_{G_{V_n,\infty}}\ker_n(xk,h)\,\Lambda_n\bigl(\varrho_n(k)^{-1}u\bigr)\,d\kappa(k)=\Lambda_{n+1}\bigl(\varrho_{n+1}(h)^{-1}\,
\mathrm{comp}_n(x,h_{\mathrm f})\,u\bigr).
\end{equation}
\end{enumerate}
\formal{GlobalTheory.JetsOne.ArrowOK}
\end{definition}

By \eqref{eq:pf:rowprops:arrowok}, if $F$ is finite-real then its lift along a good arrow is finite-real, with the vector
obtained by averaging $\mathrm{comp}_n$ over the finite class representatives; iterating along a chain of good arrows,
the lift of a finite-real function is a weighted sum over chains of class representatives of the composites of the maps
$\mathrm{comp}_n$, the weights being the inverse orders $\gamma$ of the finite stabilisers. This is the mechanism of the
following gluing statement.

\begin{lemma}[Gluing]\label{pf:rowprops:glue}
Let $(V_n)_{n\le N}$ and $(V'_n)_{n\le N}$ be two chains of totally definite hermitian spaces with data as in
\cref{def:pf:rowprops:jetdata}, all arrows $n<N$ good, for a finite combination $\tau$ of pure data $d$ (each giving kernels
$\ker_n^d,\ker_n'^d$ and maps $\mathrm{comp}^d_n,\mathrm{comp}'^d_n$). Let $\phi$ on $G_{V_0}(\A)$ and $F$ on $G_{V'_0}(\A)$ be
finite-real for $(\varrho_0,\Lambda_0,v)$ and $(\varrho'_0,\Lambda'_0,v')$, left invariant under rational points and right
invariant under a neighbourhood of $1$ in the finite points, and suppose that the kernels are locally constant to the
same extent. Suppose that at the last node $\Lambda_N$ and $\Lambda'_N$ are compatible, through an embedding
$s:G_{V'_N}\to G_{V_N}$ agreeing with the stabiliser embedding, with a bilinear pairing $B$ of $\mathcal M_N$ and
$\mathcal M'_N$. Let the measures and fundamental domains be those of definite spaces at given open compact levels
$K_n$, $K'_n$, and suppose that the \emph{pairing of points} (the pairing by $B$ of the propagated vectors at chains of
finite points) is invariant under right multiplication of each point by its level subgroup. Then
\begin{multline*}
\sum_d\tau(d)\int_{\rowpropsfd_{V'_N}}\mathrm{lift}_{\ker^d}(\phi)\bigl(s(h)\bigr)\,\mathrm{lift}_{\ker'^d}(F)(h)\,dh\\
=\sum_{c}\gamma(c)\,\bigl(\text{pairing of points at the representatives of }c\bigr),
\end{multline*}
the sum running over the families $c$ of double cosets of the levels.
\formallabel{GlobalTheory.JetsOne.glue}
\end{lemma}

\begin{proof}
\emph{Finer levels.} Finitely many conditions are involved (one for each datum $d$ in the support of $\tau$ and each node),
each of which holds on a neighbourhood of $1$ in the finite points; so there are open compact subgroups $J_n\subseteq K_n$
($n<N$) and $J'_n\subseteq K'_n$ ($n\le N$) such that all kernels $\ker^d_n$, $\ker'^d_n$ and the functions $\phi$, $F$ are
invariant under right translation by $J_n$, resp.\ $J'_n$, and at the last node also under $s(J'_N)$.

\emph{One datum.} For each $d$, the definite integral $\int_{\rowpropsfd_{V'_N}}\mathrm{lift}_{\ker^d}(\phi)(s(h))
\mathrm{lift}_{\ker'^d}(F)(h)\,dh$ equals the sum, over the families of double cosets of the finer levels $J$, $J'$, of the
products of the inverse orders of the finite stabilisers times the pairing by $B$ of the vectors propagated along the two
chains at the representatives (\lref{GlobalTheory.JetsOne.one_term}). This rests on the decomposition of the quotient
$G_V(F)\backslash G_V(\A)$ of a definite group into the finitely many images of the sets $G_{V,\infty}\,x\,J$, $x$ running over
class representatives, and on \eqref{eq:pf:rowprops:arrowok}, by which the integrals over the compact archimedean groups
produce the propagated vectors.

\emph{Regrouping.} The pairing of points $\Pi=\sum_d\tau(d)\,(\text{pairing of } d)$ is invariant under left multiplication
of the points by rational elements (\lref{GlobalTheory.JetsOne.pairP_rat}, from the invariance of the kernels and of $\phi$,
$F$ under rational points), and by hypothesis under right multiplication by the coarse levels $K_n$, $K'_n$. For such a
function, the weighted sum over the classes of the finer levels equals the weighted sum over the classes of the coarser
levels: each coarse double coset is a disjoint union of fine ones, and the weights add up
(\lref{GlobalTheory.ClassRegroup.finCl_regroup_prod}). Summing the previous identity over $d$ with the coefficients
$\tau(d)$ and regrouping gives the claim.
\end{proof}

\begin{theorem}[Level invariance of the pairing of points]\label{pf:rowprops:levelinv}
Let $S$ be a grid choice for a grid satisfying the level conditions at $\mu_0$, such that $p_{de}$ and $p'_{de}$ are coprime to
$Q_S\,q$ and the guard condition holds at $t\in U$. Let $\beta\in\rowpropsSupp(t,\Phi)$ and let $\tau$ be an expansion of $\beta$
whose data have vectors. For vector-valued functions $\Psi$ on the finite points of $G_{V_{\mathrm t_0}}$ and $F$ on those of
$G_{V_{0,m}}$, put $\Pi(x,y)=\langle\Psi(x_0),F(y_0)\rangle_{\tau}(x,y)$, the pairing by $\mathrm{br}$ of the vectors propagated
by the jet maps along the column points $x=(x_j)$ and the last-row points $y=(y_i)$ (as in \cref{def:pf:rowprops:classsum}).
Then:
\begin{enumerate}[label=(\alph*)]
\item $\Pi(x[j\leftarrow x_jk],y)=\Pi(x,y)$ for $0<j<m$ and $k$ in the level of $\mathrm t_j$ attached to $\beta$;
\item if $\Psi$ is right $K_+$-invariant, $\Pi(x[0\leftarrow x_0k],y)=\Pi(x,y)$ for $k\in K_+$;
\item $\Pi(x,y[i\leftarrow y_ik])=\Pi(x,y)$ for $0<i<m$ and $k$ in the level of $(i,m)$;
\item if $x_m=\iota_{\mathrm t}(y_m)$, then $\Pi(x[m\leftarrow\iota_{\mathrm t}(y_mk)],y[m\leftarrow y_mk])=\Pi(x,y)$ for $k$ in the
level of $(m,m)$.
\end{enumerate}
\formallabel{GlobalTheory.T.normalisedKernelExpansionSum_levelInvariant}
\formallabel{GlobalTheory.ExpansionInvariance.levels_invariant}
\formallabel{GlobalTheory.T.expansionPairing_invariant_inner_vertex_level}
\formallabel{GlobalTheory.ExpansionInvariance.pairPts_update_top}
\formallabel{GlobalTheory.T.expansionPairing_invariant_parabolic_basePoint}
\formallabel{GlobalTheory.TopZeroInv.pairPts_update_top_zero}
\formallabel{GlobalTheory.T.expansionPairing_invariant_innerVertex_level}
\formallabel{GlobalTheory.RowVertexInv.pairPts_update_row}
\formallabel{GlobalTheory.T.expansion_pairing_invariant_last_vertex_level}
\formallabel{GlobalTheory.LastVertexInv.pairPts_last_invariant}
\end{theorem}

\begin{proof}
The four statements are proved by the same three steps, which we describe for (a). The pairing $\Pi$ is a sum over the
data $d$ of $\tau$ of pairings of vectors propagated by the jet maps
$k_n=\mathrm{scal}_n\cdot\mathrm{fc}(\mathrm{vec}(d_n);x_nu_n,x_{n+1}u_{n+1})$ of the normal datum, where $u_n$ are the base
points of the standard datum, $\mathrm{scal}_n$ are normalising scalars, and $\mathrm{fc}(f;x,x')$ is the matrix, in the
algebraic representations, of the Fock component of the Schwartz function $f$ at the finite points $x$, $x'$.
\begin{enumerate}[label=(\arabic*)]
\item \emph{Translation.} Replacing $x_j$ by $x_jk$ changes only the jet maps of the two arrows adjacent to $\mathrm t_j$. The
Fock component of a theta function at translated points is the Fock component of the translated Schwartz function at the
original points (\lref{GlobalTheory.FockTranslate.fockComp_translate}). Hence $\Pi(x[j\leftarrow x_jk],y)$ is the pairing
at the original points $(x,y)$ computed with the \emph{moved} vectors of the data $d$: the Weil translates of $\mathrm{vec}(d_{j-1})$
and $\mathrm{vec}(d_j)$ by $g=u^{-1}ku$, $u$ the base point of the standard datum at $\mathrm t_j$.
\item \emph{Dependence on product functions.} The pairing at fixed points, as a function of the families of vectors of the
data, depends only on the finite combination of product functions
$\sum_d\tau(d)\prod_n\mathrm{vec}(d_n)$ (\lref{GlobalTheory.ExpansionInvariance.pairV_eq_of_prodFun_eq}).
\item \emph{Invariance of the product function.} The product function of the moved vectors equals that of the original
vectors, because the product function of the expansion is the product of the local data of the terms of $\beta$, which are
invariant under the level at $\mathrm t_j$ attached to $\beta$ (\lref{GlobalTheory.T.expansion_product_invariant_inner_vertex_level}).
\end{enumerate}
For (b) the same steps use the moved vectors at $\mathrm t_0$ (\lref{GlobalTheory.TopZeroInv.pairVec_move},
\lref{GlobalTheory.TopZeroInv.pairVec_sum_eq_of_prod_eq}) and the invariance of the product function under $K_+$
(\lref{GlobalTheory.T.unfoldingExpansion_productFunction_fixed_by_Kplus}); the right $K_+$-invariance of $\Psi$ takes care of the
argument $\Psi(x_0k)=\Psi(x_0)$. For (c) they use the moved vectors on the last row
(\lref{GlobalTheory.RowVertexInv.pairFn_update_moveV}), the multilinear dependence on product functions
(\lref{GlobalTheory.finsupp_sum_multilinear_eq_of_prod_eq}) and the invariance of the expansion tensor under the levels of
the last row (\lref{GlobalTheory.T.expansionTensorSum_invariant_gridRowLevel}). For (d) the points of $\mathrm t_m$ and $(m,m)$
move together through the stabiliser embedding; the last jet maps of the column and of the row change by the moved vectors
(\lref{GlobalTheory.LastColumnMove.colMapV_update_last}, \lref{GlobalTheory.LastRowMove.rowMapV_update_last}), the pairing depends
multilinearly on the product functions (\lref{GlobalTheory.funProd_multilinear_sum_eq}), and the product function of the
total tensor is invariant under the level of $(m,m)$ acting on the last arrows of both chains
(\lref{GlobalTheory.T.productFunction_totalTensor_invariant_lastArrows}).
\end{proof}

\begin{proposition}[First-row invariance]\label{pf:rowprops:firstrow}
Assume $F$ and $\sfname{SumOKK}$. At large scale, for every admissible choice, all $t$, $a$, $\beta\in\rowpropsSupp(t,\Phi_a)$, all
expansions $\tau$ of $\beta$, $u\in H_X(t)$, $v$, points $x,y$ and every $k$ in the level of $(0,m)$ attached to $t$,
\[
\Pi_v\bigl(x,y[0\leftarrow y_0k]\bigr)=\Pi_v(x,y),\qquad
\Pi_v(x,y)=\bigl\langle v,\mathbf{f}(t,u)(y_0)\bigr\rangle_\tau(x,y).
\]
\formallabel{GlobalTheory.FirstRowLaw.first_row_invariant_under_twist}
\end{proposition}

\begin{proof}
At large scale: the level of $(0,m)$ attached to $t$ is contained in the level group $J$ of the forms of $R$ on $X'$ at $t$
(\lref{GlobalTheory.T.gridBottomLevel_le_formLevel_largeDepth}); the vector-valued forms transform under $J$ by its
character $\chi$: $\mathbf{f}(t,u)(yk)=\chi(k)\,\mathbf{f}(t,u)(y)$ for $k\in J$
(\lref{GlobalTheory.FirstRowLaw.formC_mul}); and the pairing of points has the compensating transformation law
$\langle v,\chi(k)w\rangle_\tau(x,y[0\leftarrow y_0k])=\langle v,w\rangle_\tau(x,y)$ for $k$ in the level of $(0,m)$
(\lref{GlobalTheory.T.vertex_pair_points_transformation_law}, which uses $F$ and $\sfname{SumOKK}$). With $w=\mathbf{f}(t,u)(y_0)$
these three facts give the claim.
\end{proof}

\begin{proposition}[Jets]\label{pf:rowprops:jetsmain}
Assume $F$ and $\sfname{SumOKK}$. Then $\sfname{StepJets}$ holds. More precisely, at large scale, for every admissible
choice and every polynomial pair $P$ there is $\kappa_P\in\C$ such that for all $t$, $a$, $\beta\in\rowpropsSupp(t,\Phi_a)$,
term volumes $M$, expansions $\tau$ of $\beta$ and $y$,
\begin{equation}\label{eq:pf:rowprops:jetsone}
\rowpropsdef(M;\delta_P,\tau;\phi_a,F_{t,y})=\kappa_P\,\rowpropsclass_{\Phi_a}\bigl(\beta,\tau;\psi_C(\Psi_a),\mathbf{f}(t,y)\bigr).
\end{equation}
\formallabel{GlobalTheory.T.twisted_jets_assemble}
\formallabel{GlobalTheory.RowClaim.JetsAssembly.twistedJet_succ_eq_assembly}
\formallabel{GlobalTheory.RowClaim.JetsAssembly.twisted_jet_one_of_all}
\formallabel{GlobalTheory.T.twisted_pairing_identity}
\formallabel{GlobalTheory.RowClaim.JetsAssembly.jetsAt_of_one}
\formallabel{GlobalTheory.RowClaim.JetsAssembly.defTot_eq_sum}
\end{proposition}

\begin{proof}
\emph{Step 1: the identity \eqref{eq:pf:rowprops:jetsone}.} Fix an admissible choice at large scale; $\mathsf{SumOK}(t,\Phi_a)$
holds for all $t,a$, and the first-row invariance (\cref{pf:rowprops:firstrow}, which uses $F$ and $\sfname{SumOKK}$) holds: the
pairing of points of an expansion is invariant under right multiplication of the first point of the last row by the level
subgroup of $(0,m)$. The ingredients are:
\begin{enumerate}[label=(\alph*)]
\item the scalar column functions $\phi_a$ are finite-real at $\mathrm t_0$ for a fixed linear form $\Lambda_0$ and the
vectors $\psi_C(\Psi_a)$ transported by the normal frame (\lref{GlobalTheory.JetsOne.col_real}), and the forms
$F_{t,y}$ are finite-real at $(0,m)$ for a fixed $\Lambda'_0$ and the vectors $\mathbf{f}(t,y)$, read through the
algebraic representation of $(0,m)$ (\lref{GlobalTheory.JetsOne.form_real}); both are left invariant under rational
points and locally constant on finite points (the column is right $K_+$-invariant, and the forms of $R$ are automorphic);
\item for the pair $P$, there are linear forms $\Lambda_j$ on the representations of the top nodes $\mathrm t_j$ with
$\Lambda_0$ as given, such that for every datum $d$ of an expansion (with vectors) every arrow $j<m$ of the column chain,
with the kernels of the column lift with polynomial $P_1$ and the jet maps $\mathrm{J}^{\mathrm{col}}$ of $d$, is good
(\lref{GlobalTheory.JetsOne.col_chain}); and similarly linear forms $\Lambda'_i$ on the representations of the
nodes $(i,m)$ of the last row, starting from $\Lambda'_0$, making the arrows of the row chain good with the jet maps
$\mathrm{J}^{\mathrm{row}}$ (\lref{GlobalTheory.JetsOne.row_chain});
\item there are $c_P\in\C$ and a homomorphism $s:G_{V_{m,m}}\to G_{V_{\mathrm t_m}}$ agreeing with the stabiliser embedding
such that $\Lambda_m$ and $\Lambda'_m$ are compatible through $s$ with the pairing $c_P\cdot\mathrm{br}$
(\lref{last_functional_compatible});
\item the pairing of points of an expansion is invariant under the level subgroups at the nodes $\mathrm t_j$ ($0<j<m$)
and $(i,m)$ ($0<i<m$), under $K_+$ at $\mathrm t_0$, and under the level of $(m,m)$ acting jointly on the last points
(\cref{pf:rowprops:levelinv}, applicable since the Hida setting gives the level conditions, the coprimality of $p_{de}$,
$p'_{de}$ with $Q_S\,q$ and the guard condition), and at $(0,m)$ by the first-row invariance;
\item the term volumes are measures of definite spaces at open compact levels, the levels $K_+$ at $\mathrm t_0$ being
part of them; the top and last-row spaces are positive definite at every complex embedding.
\end{enumerate}
By \cref{pf:rowprops:glue}, applied with $\tau$, the chains $(V_{\mathrm t_j})_{j\le m}$ and $(V_{i,m})_{i\le m}$, the
functions $\phi_a$ and $F_{t,y}$, the pairing $c_P\cdot\mathrm{br}$ and the pairing of points multiplied by $c_P$, the sum
$\sum_d\tau(d)\rowpropsdef(M;P,d;\phi_a,F_{t,y})$ (the definite period is the integral of the product of the column lift
composed with the stabiliser embedding and the row lift) equals $c_P\sum_c\gamma(c)\langle\psi_C(\Psi_a)(\mathrm{pt}(c)_0),
\mathbf{f}(t,y)(\mathrm{pt}(c)_0)\rangle_{\tau,c}$, i.e.\ $c_P$ times the class sum. Since $c_P$ depends only on $P$ (and
the choice), \eqref{eq:pf:rowprops:jetsone} holds with $\kappa_P=c_P$.

\emph{Step 2: $\mathsf{JetsAt}$.} Term volumes exist. Given $\sigma$, put $\kappa=\sum_P\sigma(P)\kappa_P$. Since
$\rowpropsdef(M;\sigma,\tau;\phi,F)=\sum_P\sigma(P)\rowpropsdef(M;\delta_P,\tau;\phi,F)$, \eqref{eq:pf:rowprops:jetsone} gives
\eqref{eq:pf:rowprops:jetsat}.
\end{proof}

% ---------------------------------------------------------------------------
\subsection{Smooth cuspidal lifts}\label{sec:pf:rowprops:cusp}

The convergence statements used in the flip (\cref{pf:rowprops:flipmain}) rest on the following theorem on lifts of
smooth cusp forms of prescribed archimedean type, which complements the statement $\sfname{LiftsCuspidal}$ on
vector-valued cusp forms of the multiplicity spaces.

\begin{definition}[Cusp forms of type]\label{def:pf:rowprops:cusptype}
For a node $N$ of a grid pre-datum (a hermitian space with archimedean type data $S$), a function $\varphi$ on
$G_{V_N}(\A)$ is a \emph{cusp form of type} if it is a smooth automorphic function, cuspidal, and of type $S(\tau)$ at every
real place $\tau$. For nodes $A$, $B$, measure data on $A$ and a kernel $K$ on $G_{V_B}(\A)\times G_{V_A}(\A)$, the kernel
\emph{lifts smoothly} if for every cusp form $\varphi$ of type at $A$ and every $h$, $g\mapsto K(h,g)\varphi(g)$ is integrable
on the fundamental domain and $h\mapsto\int K(h,g)\varphi(g)\,dg$ is a cusp form of type at $B$. A datum $D$ satisfies the
\emph{size condition} if $r>1800$ and consecutive Harish-Chandra parameters of the column differ by more than $2N^3$ for
the relevant bound $N$.
\formal{GlobalTheory.IndefiniteRowSideFullHecke.CuspAlongSmoothLifts.LiftsSmooth}\formal{GlobalTheory.IndefiniteRowSideFullHecke.CuspAlongSmoothLifts.Sizes}
\end{definition}

\begin{theorem}[Smooth cuspidal lifts]\label{pf:rowprops:cuspsmooth}
Assume $F$ and $\sfname{LiftsCuspidal}$. Let $D$ satisfy the size condition, $R$ a realisation and $\rowpropsG$ a grid pre-datum
over $R$ satisfying the base conditions (level conditions at $\mu_0$, archimedean data, measure models of the column). Then
for every theta function $\Phi$ of the Weil representation of the relevant pair, the theta kernels of $\Phi$ lift smoothly
along: (X) the rightward arrows $(i,j)\to(i+1,j)$; (th) the downward arrows $(i,j+1)\to(i,j)$; (up), (dn) the two arrows
between consecutive core nodes; and (col) the arrows of the top column.
\formallabel{GlobalTheory.T.cuspidal_lift_smooth_at_cusps}
\formallabel{GlobalTheory.CuspSmoothLift.cuspidal_lift_smooth_at_of_cuspidal_lift}
\end{theorem}

\begin{proof}
Each of the five clauses is the conjunction of four statements: integrability of the kernel against cusp forms of type
on the fundamental domain, smoothness and automorphy of the lift, its cuspidality, and its archimedean type. They are
\begin{itemize}
\item for (X): \lref{GlobalTheory.T.twisted_theta_kernel_mul_cusp_form_integrableOn},
\lref{GlobalTheory.T.twisted_theta_lift_isSmoothAutomorphic}, \lref{GlobalTheory.T.twisted_theta_lift_isCuspidal},
\lref{GlobalTheory.T.twisted_theta_lift_hasKType};
\item for (th): \lref{GlobalTheory.T.twisted_theta_integral_integrableOn_automorphic_quotient},
\lref{GlobalTheory.T.twisted_theta_lift_of_cusp_form_isSmoothAutomorphic},
\lref{GlobalTheory.T.twisted_theta_lift_of_cusp_form_isCuspidal},
\lref{GlobalTheory.T.twisted_theta_lift_of_cusp_form_hasKType};
\item for (up): \lref{GlobalTheory.T.twisted_upward_theta_integral_converges},
\lref{GlobalTheory.T.twisted_upward_theta_lift_isSmoothAutomorphic},
\lref{GlobalTheory.T.twisted_upward_theta_lift_isCuspidal}, \lref{GlobalTheory.T.twisted_upward_theta_lift_hasKType};
\item for (dn): \lref{GlobalTheory.T.twisted_downward_theta_integral_integrableOn},
\lref{GlobalTheory.T.twisted_downward_theta_lift_isSmoothAutomorphic},
\lref{GlobalTheory.T.twisted_downward_theta_lift_isCuspidal},
\lref{GlobalTheory.T.twisted_downward_theta_lift_hasKType};
\item for (col): \lref{GlobalTheory.T.columnConvolution_integrableOn},
\lref{GlobalTheory.T.columnLift_isAutomorphicFunction}, \lref{GlobalTheory.T.columnLift_isCuspForm},
\cref{pf:rowprops:coltype}.
\end{itemize}
In each case the integrability statement is proved first and the other three are derived from it and from
$\sfname{LiftsCuspidal}$ and $F$.
\end{proof}

The analytic weight of \cref{pf:rowprops:cuspsmooth} lies in a few common ingredients, which reduce statements about
smooth cusp forms of type to the statement $\sfname{LiftsCuspidal}$ about vector-valued cusp forms of the multiplicity
spaces; the Casselman--Wallach globalisation provides the framework of smooth vectors
(\lref{GlobalTheory.TypeSum.casselmanWallach_banach}, \cite{Casselman1989,Wallach1992}). We prove the type statement of
clause (col) and three of these ingredients: the description of functions in the closure of the discrete series coefficients
(\cref{pf:rowprops:closure}), which is used to identify archimedean types; the approximation of lifts along rightward arrows by
lifts with Fock vectors (\cref{pf:rowprops:approxfock}), which enters the cuspidality and the type of these lifts; and the
$K$-finite decomposition of averaged cusp forms (\cref{pf:rowprops:kfinite}), which enters the approximation of cusp forms of
type by values of vector-valued cusp forms.

\begin{proposition}[Types of column lifts]\label{pf:rowprops:coltype}
Assume $\sfname{LiftsCuspidal}$ and $F$. In the situation of \cref{pf:rowprops:cuspsmooth}, let $j<m$ and $t\in T$, let $\Phi$ be a
theta function for the Weil representation of the arrow $\uparrow^t_j:\mathrm t_j\to\mathrm t_{j+1}$, and let $\varphi$ be a cusp
form of type at $\mathrm t_j$ such that $g\mapsto K(\Phi;g,h)\varphi(g)$ is integrable on the fundamental domain for every $h$.
Then the lift $h\mapsto\int K(\Phi;g,h)\varphi(g)\,dg$ (over the fundamental domain) is of the type of $\mathrm t_{j+1}$ at every
real place.
\formallabel{GlobalTheory.T.columnLift_hasTypeAtPlaces}
\formallabel{GlobalTheory.ColTypeLift.columnLift_isTypeAt_place}
\end{proposition}

\begin{proof}
The spaces $V_{\mathrm t_j}$ and $V_{\mathrm t_{j+1}}$ satisfy the signature conditions and are totally definite
(\lref{GlobalTheory.ColTypeLift.xp_definite}). On a definite group a cusp form of type is a finite sum
$\varphi=\sum_iF_i(w_i)$ of vector-valued cusp forms $F_i$ of the multiplicity space of the node evaluated at smooth vectors
$w_i$ of the discrete series model (\lref{GlobalTheory.T.definite_unitary_cuspidal_decomposition}, which uses $F$). The
lift of such a finite sum, when convergent, is again a finite sum $\sum_iF'_i(z_i)$ of the same kind at $\mathrm t_{j+1}$
(\lref{GlobalTheory.T.smooth_of_theta_lift_finite_sum}, which uses $\sfname{LiftsCuspidal}$ and $F$). Finally, at a definite
node the local discrete series models are finite dimensional, and a finite sum of values of vector-valued cusp forms of the
multiplicity space at smooth vectors is of the type of the node at every real place
(\lref{GlobalTheory.ColTypeLift.isTypeAt_sum}).
\end{proof}

\begin{theorem}[Closure of discrete series coefficients]\label{pf:rowprops:closure}
Assume the theory of Harish-Chandra modules and the Casselman--Wallach globalisation ($\sfname{HarishChandraModules}$,
$\sfname{CasselmanWallach}$, fields of $F$). Let $\varepsilon\in\{\pm\}^n$, let $\lambda$ be a Harish-Chandra parameter, and let
$\sigma$ be a discrete series representation of $\Uni(J_\varepsilon)$ on a Hilbert space $H$ with parameter $(\varepsilon,\lambda)$.
Let $f$ be a continuous function on $\Uni(J_\varepsilon)$ of uniform moderate growth which is, on every compact set, a uniform
limit of finite sums $\sum_j\ell_j(\sigma(y)v_j)$ with Sobolev-continuous linear forms $\ell_j$ and smooth vectors $v_j$. Then
$f$ is the sum of a convergent series of matrix coefficients of $\sigma$: there are a finite set $S$ of words in the Lie
algebra whose Sobolev norm $\|\cdot\|_S$ satisfies $\|\sigma(g)w\|_S\le B(g)\|w\|_S$ for smooth $w$ and a continuous function
$B$, smooth vectors $v_i$, linear forms $\ell_i$ and constants $C_i\ge0$ with $|\ell_i(w)|\le C_i\|w\|_S$ for smooth $w$
($i\in\N$), such that $\sum_iC_i\|v_i\|_{S'}<\infty$ for every finite set $S'$ of words and $f(g)=\sum_i\ell_i(\sigma(g)v_i)$
for all $g\in\Uni(J_\varepsilon)$.
\formallabel{GlobalTheory.ClosureSeries.closure_model}
\formallabel{GlobalTheory.T.hasAllKTypes_of_closure_discreteSeriesCoefficients}
\formallabel{GlobalTheory.ClosureSeries.hasTypeAll_of_closure}
\end{theorem}

\begin{proof}
The representation $\sigma$ satisfies the standing conditions of the development for discrete series (in particular its
$K$-finite vectors are smooth; \lref{GlobalTheory.TypeSum.std_of_isDS}).
By \lref{GlobalTheory.T.discreteSeriesCoeffSeries_of_kTypeProjections}, a continuous function $f$ of uniform moderate
growth is the sum of such a series as soon as, for every $K$-finite function $e$ on the maximal compact subgroup $K$ (that is,
$e$ lies in a finite-dimensional space of continuous functions stable under left and right translations), the average
$e*f(y)=\int_Ke(k)f(yk)\,dk$ is a finite sum $\sum_j\ell_j(\sigma(y)x_j)$ on $\Uni(J_\varepsilon)$ with Sobolev-continuous
$\ell_j$ and $K$-finite $x_j$. Fix $e$.
\begin{enumerate}[label=(\arabic*)]
\item $e*f$ is continuous and of uniform moderate growth
(\lref{GlobalTheory.T.compactAverage_continuousOn_uniformModerateGrowth}).
\item There are finitely many $K$-finite vectors $w_1,\dots,w_d$ such that on every compact set $e*f$ is a uniform limit of
finite sums $\sum_aL_a(\sigma(y)w_a)$ with Sobolev-continuous $L_a$
(\lref{GlobalTheory.T.kProjection_limit_discreteSeriesCoeffs}, applied to the approximations of $f$ in the hypothesis).
\item A continuous function of uniform moderate growth which is such a locally uniform limit of matrix coefficients of the
fixed $K$-finite vectors $w_a$ is itself equal to $\sum_aL_a(\sigma(y)w_a)$ for suitable Sobolev-continuous $L_a$
(\lref{GlobalTheory.T.eq_matrixCoeff_of_locallyUniform_limit_kFinite}, which uses the theory of Harish-Chandra modules and
the Casselman--Wallach globalisation).
\end{enumerate}
By (1)--(3), $e*f$ has the required form, and the criterion gives the expansion of $f$.
\end{proof}

\begin{proposition}[Approximation of lifts by Fock lifts]\label{pf:rowprops:approxfock}
Assume $\sfname{LiftsCuspidal}$ and $F$. In the situation of \cref{pf:rowprops:cuspsmooth}, let $(i,j)$ with $i<m$, $j\le m$, let
$\Phi$ be a theta function for the rightward arrow $(i,j)\to(i+1,j)$ and $\varphi$ a cusp form of type at $(i,j)$ such that
$g\mapsto K(\Phi;g,h)\varphi(g)$ is integrable on the fundamental domain for every $h$ (integrals below are over the
fundamental domain). For every compact set $C\subseteq G_{V_{i+1,j}}(\A)$ and $\delta>0$ there are finitely many finite Fock
vectors $f_l$ with polynomials $p_l$, vector-valued cusp forms $F_l$ of the multiplicity space of $(i,j)$ and smooth vectors
$v_l$ of its discrete series model such that, for all $h\in C$,
\[
\Bigl|\int K(\Phi;g,h)\varphi(g)\,dg-\sum_l\int K\bigl(\mathrm{fock}(f_l,p_l);g,h\bigr)\,F_l(v_l)(g)\,dg\Bigr|<\delta .
\]
\formallabel{GlobalTheory.T.cusp_type_approximation_by_fock_lift}
\formallabel{GlobalTheory.CuspTypeApproxX.cuspType_lift_approx_fockSpace}
\end{proposition}

\begin{proof}
By \lref{GlobalTheory.T.cusp_forms_approximated_by_theta_lifts} there are finitely many vector-valued cusp forms $F_k$ of
the multiplicity space and smooth vectors $v_k$ ($k<m'$), with $K(\Phi;\cdot,h)F_k(v_k)$ integrable, such that
\[
\Bigl|\int K(\Phi;g,h)\varphi(g)\,dg-\sum_{k<m'}\int K(\Phi;g,h)F_k(v_k)(g)\,dg\Bigr|<\delta/2\qquad(h\in C).
\]
By \lref{GlobalTheory.T.theta_lifts_dense_in_cusp_forms}, for each $k$ there are finitely many finite Fock vectors
$f_{k,q}$ and polynomials $p_{k,q}$ such that, for $h\in C$,
\[
\Bigl|\int K(\Phi;g,h)F_k(v_k)(g)\,dg-\sum_q\int K\bigl(\mathrm{fock}(f_{k,q},p_{k,q});g,h\bigr)F_k(v_k)(g)\,dg\Bigr|
<\frac{\delta}{2(m'+1)}.
\]
Take the family indexed by the pairs $(k,q)$, with
$f_{(k,q)}=f_{k,q}$, $p_{(k,q)}=p_{k,q}$, $F_{(k,q)}=F_k$ and $v_{(k,q)}=v_k$. By the triangle inequality the error on $C$
is less than $\delta/2+m'\delta/(2(m'+1))<\delta$.
\end{proof}

\begin{proposition}[$K$-finite decomposition of averaged cusp forms]\label{pf:rowprops:kfinite}
Assume $F$ and let $\theta$ be admissible archimedean data. Let $N$ be a node with discrete series model $M$, let $U$ be an
open subgroup of $G_{V_N}(\A)$ and let $\varphi$ be a cusp form of type at $N$ with $\varphi(gk)=\varphi(g)$ for all $g$ and all
$k\in U\cap G_{V_N}(\A_f)$. For every real place $\tau$ let $e_\tau$ be a $K$-finite function on the maximal compact subgroup
$K_\tau$ (as in the proof of \cref{pf:rowprops:closure}), and write
$\mathrm{av}_\tau\psi(g)=\int_{K_\tau}e_\tau(k)\,\psi(gk)\,dk$, with $k$ placed at $\tau$ through the frame of $\theta$. Let
$l=(\tau_1,\dots,\tau_r)$ be a list containing every real place and $\varphi_l=\mathrm{av}_{\tau_1}\cdots\mathrm{av}_{\tau_r}\varphi$.
Then there are finitely many vector-valued cusp forms $F_j$ of the multiplicity space of $N$ and smooth vectors $v_j$ of $M$
such that each $F_j(v_j)$ is right invariant under $U\cap G_{V_N}(\A_f)$ and $\varphi_l=\sum_jF_j(v_j)$.
\formallabel{GlobalTheory.T.twist_cusp_average_k_finite_decomposition}
\formallabel{GlobalTheory.CuspKFinDecomp.cusp_average_k_finite_decomposition}
\end{proposition}

\begin{proof}
We show by induction on lists $l'$ that $\varphi_{l'}$ is a cusp form of type, right invariant under $U\cap G_{V_N}(\A_f)$,
and $K_\tau$-finite at every $\tau\in l'$. For the empty list this is the hypothesis. For $l'=(\tau,l'')$ we have
$\varphi_{l'}=\mathrm{av}_\tau\varphi_{l''}$. It is a cusp form of type (\lref{GlobalTheory.T.averaging_preserves_cusp_type}).
It is right invariant under $U\cap G_{V_N}(\A_f)$, since for such $k$ the element placed at $\tau$ commutes with $k$, so that
$\mathrm{av}_\tau\varphi_{l''}(gk)=\int e_\tau(k')\varphi_{l''}(gk'k)\,dk'=\mathrm{av}_\tau\varphi_{l''}(g)$. It is
$K_\tau$-finite by \lref{GlobalTheory.T.averaged_automorphic_function_K_finite}, and $K_{\tau'}$-finite for
$\tau'\in l''$ because averaging at $\tau$ preserves $K_{\tau'}$-finiteness
(\lref{GlobalTheory.T.average_preserves_K_finiteness}).

Applied to $l$, which contains every real place, $\varphi_l$ is a cusp form of type, right invariant under
$U\cap G_{V_N}(\A_f)$ and $K$-finite at all real places. Choose a Haar measure and a fundamental domain for the rational
points (\lref{GlobalTheory.Mem.AdelicGroupFund.rightInvariant_and_exists_fundamentalDomain}). By
\lref{GlobalTheory.T.cuspForm_decomposes_into_kFinite_components}, $\varphi_l=\sum_jT_j(u_j)$ with finitely many tuples
$u_j=(u_{j,\tau})_\tau$ of $K$-finite vectors of the local models $M_\tau$ and maps $T_j$ on such tuples which are
$(\mathfrak g,K)$-module maps into the cusp forms right invariant under $U\cap G_{V_N}(\A_f)$. By the Casselman--Wallach
globalisation each $T_j$ extends to a map $\tilde T_j$ on tuples of smooth vectors with values in the multiplicity space,
agreeing with $T_j$ on $K$-finite tuples and preserving the invariance under $U\cap G_{V_N}(\A_f)$
(\lref{GlobalTheory.GKExt.extExists}). The $K$-finite vectors of a discrete series are smooth
(\lref{GlobalTheory.TypeSum.std_of_isDS}, \lref{GlobalTheory.T.mem_smoothVectors_discreteSeriesModel_iff}). Let $F_j$
be the linear map on the tensor product $\bigotimes_\tau M_\tau^\infty$ induced by the multilinear map $\tilde T_j$ and
$v_j=\bigotimes_\tau u_{j,\tau}$. Then $F_j$ is a vector-valued cusp form of the multiplicity space, $F_j(v_j)=T_j(u_j)$ is right
invariant under $U\cap G_{V_N}(\A_f)$, and $\varphi_l=\sum_jF_j(v_j)$.
\end{proof}

\begin{lemma}[Cusp steps]\label{pf:rowprops:cuspstep}
Assume $\sfname{LiftsCuspidal}$ and $F$. At large scale, for every admissible choice, every $t$ and every volume datum $M$
of level $t$, the cusp step holds for $(W,M,\mathrm{rm},t)$.
\formallabel{GlobalTheory.FlipRun.cuspStep_of_large_flipRun}
\formallabel{GlobalTheory.T.flipRun_cusp_transition}
\end{lemma}

\begin{proof}
By \cref{pf:rowprops:cuspsmooth}, smooth cuspidal lifting holds for every datum. The arrows of the grid are rightward or
downward. For rightward arrows the cusp step follows from smooth cuspidal lifting and $F$ by
\lref{GlobalTheory.T.cusp_transition_rightward}, and for downward arrows by
\lref{GlobalTheory.T.flipRun_descending_cusp_transition}.
\end{proof}

% ---------------------------------------------------------------------------
\subsection{Remarks on the formal statements}\label{sec:pf:rowprops:remarks}

\begin{remark}[Unused hypotheses]\label{pf:rowprops:rem:unused}
In the proof of \cref{pf:rowprops:rowthm} the hypotheses $\sfname{SquareSt.LocalMatching}$ and
$\sfname{SquareSt.ThetaIntegralsAgree}$ are not used: they are passed to the proof of $\sfname{ConstK}$
(\cref{pf:rowprops:constK}), but the three lemmas used there require only $\sfname{SquareSt.WeilConvergence}$
(\cref{pf:rowprops:bigint}), $\sfname{ThetaIntegralsAgree}$ (\cref{pf:rowprops:smallint}) and $F$,
$\sfname{ThetaIntegralsAgree}$, $\sfname{LocalMatching}$ (\cref{pf:rowprops:bignz}). The statement $\sfname{LocalMatching}$
is used only in \cref{pf:rowprops:bignz}.
\end{remark}

\begin{remark}[Real matches]\label{pf:rowprops:rem:rel}
The relation $\mathrm{rel}_{ij}$ of a real match is only required to project onto both factors and to be a proper
subspace. Such relations exist: if $\lambda$ and $\lambda'$ are the linear forms ``constant coefficient'' on the two
spaces of polynomial tensors, the subspace $\{(a,b):\lambda(a)=\lambda'(b)\}$ projects onto both factors and is proper.
The existence of real matches is proved with this relation. The square constants of a choice are defined relative to its
real match, and they are part of the data of the choice: admissibility requires each $c_{ijt}$ to be a square constant.
Thus $\sfname{ConstK}$ proves that these given constants are non-zero and that the squares are non-zero; it does not
construct the constants.
\end{remark}

\begin{remark}[The constant $\kappa$]\label{pf:rowprops:rem:kappa}
The lift formula does not assert $\kappa\ne0$. If $\rowpropsctot(t)\ne0$ and $\psi_K(a,\Theta_t(j_Xx))\ne0$ for some good $x$,
then $\kappa\ne0$: otherwise the formula gives $\rowpropsemb(\zeta_t)\rowpropsctot(t)\psi_K(a,\Theta_t(j_Xx))=0$, while
$\rowpropsemb(\zeta_t)\ne0$ because $\zeta_t$ is a root of unity. The roots of unity $\zeta_t$ are elements of
$\rowpropsLL$; they come from the subgroup $\rowpropsXi_t$ of local values of the twisting character, which consists of
roots of unity under the field condition (\cref{pf:rowprops:xivals}). The rationality statement $\sfname{SumRatK}$
concerns all rational points, the lift formula only the good ones.
\end{remark}

\begin{remark}[Uniformity]\label{pf:rowprops:rem:uniform}
$\sfname{LinK}$, $\sfname{StepProj}$ and the commutation of squares hold with $C_{F0}=0$ and $x_0=0$, i.e.\ for all admissible
choices at all scales; the integrability of the big route uses the threshold $x_0=12^{5/\eps}$. The size conditions of the
Hida setting enter through the square standing condition (\cref{pf:rowprops:standing}) and through the large-scale
statements cited in \cref{sec:pf:rowprops:unfold,sec:pf:rowprops:jets}.
\end{remark}
